%%%%%%%%%%%%%%%%%%%%%%%%%%%%%%%%%%%%%%%%%%%%%%%%%%%%%%%%%%%%%%%%%%%%%%
% Lớp 12 | Ứng dụng đạo hàm để khảo sát hàm số | Khảo sát sự biến thiên và vẽ đồ thị hàm số
% Dạng 2D1?5-1: Nhận dạng đồ thị
% Số câu: 111  (sinh tự động bởi tach_cau.py; từ nay thầy sửa trực tiếp file này)
%%%%%%%%%%%%%%%%%%%%%%%%%%%%%%%%%%%%%%%%%%%%%%%%%%%%%%%%%%%%%%%%%%%%%%

%%%%%%%%%%%%%%%% TRẮC NGHIỆM 4 PHƯƠNG ÁN %%%%%%%%%%%%%%%%

% nguồn: Phan1_Lop_12_Dai_so_Chuong_1
\begin{ex}%[2D1N5-1]%[ma:d4fa5a34]
	 \immini{
		Đường cong trong hình vẽ bên là  đồ thị của hàm số nào dưới đây?
	 \choice
			{\True  $y=\dfrac{x+1}{x-1}$}
			{$y=x^3-3x-1$} 
			{$y=\dfrac{2x-1}{x-1}$} 
			 {$y=\dfrac{x^2-x+1}{x-1}$}
			 }
			 {
			\begin{tikzpicture}[scale=.6,font=\footnotesize,line join=round,line cap=round,>=stealth]
				\draw[->]  (-4,0)--(0,0) node[below right]{$O$}--(6,0) node[below]{$x$};
				\draw[->]  (0,-3.6) --(0,4) node[right]{$y$};
				 \draw [domain=1.7:5.5,samples=100]
					plot (\x,{(\x+1)/((\x)-1)}); 
				\draw [domain=-4:0.55,samples=100] 
					plot (\x,{(\x+1)/((\x)-1)}); 
				 \draw[fill] (0,0) circle (1.5pt);
				\draw  (-4,1)--(6,1);
				\draw  (1,-3.6)--(1,4);
				\draw (-1,0) node[below] {$-1$}  circle (1.5pt);
				\draw (0,-1)  node[left] {$-1$} circle (1.5pt);
				\draw (1,0) node[below right] {$1$}  circle (1.5pt);
				\draw (0,1)  node[above left] {$1$} circle (1.5pt);
			\end{tikzpicture}
			 }
	\loigiai{ 
			Từ đồ thị hàm số đã cho, ta thấy $x=1$ là tiệm cận đứng của đồ thị hàm số, $y=1$ là tiệm cận ngang của  đồ thị hàm số.\\
			Do đó, đây là đồ thị  của hàm số $y=\dfrac{x+1}{x-1}$.
		 }
	\end{ex}

% nguồn: Phan1_Lop_12_Dai_so_Chuong_1
\begin{ex}%[2D1N5-1]%[ma:b004278f]
		Đường cong ở hình  vẽ dưới đây là đồ thị của hàm số $y=\dfrac{x+a}{bx+c}$, $(a,b,c\in \mathbb{Z})$.
		\begin{center}
			\begin{tikzpicture}[scale=0.7,font=\footnotesize,line join=round,line cap=round,>=stealth]
				\def\xmin{-3}; \def\ymin{-3}; \def\xmax{5};  \def\ymax{4};
				 \coordinate (O) at (0,0);
				\draw[->] (\xmin,0)--(\xmax,0) node[below]{$x$}; 
				 \draw[->] (0,\ymin)--(0,\ymax) node[left]{$y$};
				\fill (O)  node[below left]{$O$} circle (1pt);
				\clip ({\xmin-0.1},{\ymin-0.1}) rectangle ({\xmax+0.1},{\ymax+0.1}); 
				 \foreach \x in {2}{
					\fill (\x,0) node[below]{$\x$}  circle (1pt);
				} 
				 \foreach \y in {2}{
					\fill (0,\y)  node[left]{$\y$} circle (1pt);
				} 
				\draw[samples=100] plot[domain=\xmin:0.99](\x,{(\x-2)/(\x-1)}); 
				\draw[samples=100]  plot[domain=1.01:\xmax](\x,{(\x-2)/(\x-1)});
				\draw[dashed]  (1.0,\ymin)--(1.0,\ymax);
				\draw[dashed] (\xmin,1.0)--(\xmax,1.0); 
				\draw (1,3) node[right]{$x=1$} 
					 (-2,0.2) node[above]{$y=1$};
			\end{tikzpicture}
		\end{center}
		Giá trị biểu thức $T=a-3b-2c$  bằng
	 \choice
			{$3$} 
			 {$2$}
			 {$0$}
			{\True $-3$} 
	\loigiai{ 
			Đồ thị hàm số có tiệm cận  ngang $y=1\Rightarrow \dfrac{1}{b}=1\Leftrightarrow b=1$.\\
			Đồ thị hàm số có tiệm cận đứng $x=1\Rightarrow-\dfrac{c}{b}=1\Leftrightarrow  b=-c\Rightarrow c=-1$.\\
			Đồ thị hàm số đi qua điểm $(0;2)\Rightarrow  2=\dfrac{0+a}{1\cdot 0+(-1)} \Leftrightarrow a=-2$.\\
			Vậy $T=a-3b-2c=(-2)-3\cdot  1-2\cdot (-1)=-3$.
		 }
	\end{ex}

% nguồn: Phan1_Lop_12_Dai_so_Chuong_1
\begin{ex}%[2D1N5-1]%[ma:76133325]
		Cho hàm số $f(x)=\dfrac{ax+1}{bx+c}$ ($a,b,c\in \mathbb{R}$)  có bảng biến thiên như sau
		\begin{center}
			\begin{tikzpicture}[scale=1,font=\footnotesize]
				\tkzTabInit[nocadre=false,lgt=1.2,espcl=2.5,deltacl=0.7] 
					{$x$/0.8,$f'(x)$/0.6,$f  (x)$/2}
					 {$-\infty$,$3$,$+\infty$};
				 \tkzTabLine{,-,d,-,};
				\tkzTabVar{+/$\dfrac{1}{2}$,-D+/$-\infty$/$+\infty$,-/$\dfrac{1}{2}$}; 
			\end{tikzpicture}
		\end{center}
		Khẳng định  nào dưới đây đúng?
	 \choice
			{$\hoac{&b>\dfrac{2}{3}\\&b <  0}$}
			 {\True $\hoac{&b > 0\\&b <-\dfrac{2}{3}}$}
			{$-\dfrac{2}{3} <  b < 0$}
			{$0\le b  < \dfrac{2}{3}$}
	 \loigiai{
			Tập  xác định $\mathscr{D}=\mathbb{R}\setminus \left\{\dfrac{-c}{b} \right\}$.\\
			Từ bảng biến thiên ta có 
			\begin{itemize}
				\item $\lim\limits_{x\to 3^{+}} f(x)=+\infty$ nên đường thẳng $x=3$ là tiệm cận đứng của đồ  thị hàm số.
				\item $\lim\limits_{x\to \pm  \infty} f(x)=\dfrac{1}{2}$ nên đường thẳng $y=\dfrac{1}{2}$ là tiệm cận ngang của đồ thị hàm số.
			\end{itemize}
			Nghĩa là ta có, $\heva{&\dfrac{-c}{b}=3\\&\dfrac{a}{b}=\dfrac{1}{2}} \Rightarrow \heva{&c=-3b \\&a=\dfrac{1}{2}b}\quad  (1).$\\
			Mặt  khác $y'=\dfrac{ac-b}{\left(bx+c\right)^2} < 0\Leftrightarrow ac-b < 0\quad (2)$.\\
			 Từ $(1)$ và $(2)$ ta có $\dfrac{1}{2} b\cdot (-3b)-b < 0\Leftrightarrow-\dfrac{3}{2} b^2-b < 0\Leftrightarrow \hoac{&b > 0\\
			&b  <-\dfrac{2}{3}.}$
		} 
	\end{ex}

% nguồn: Phan1_Lop_12_Dai_so_Chuong_1
\begin{ex}%[2D1N5-1]%[ma:0d0c7198]
	\immini{Đường cong trong hình bên là đồ thị của hàm số  nào?
	 \choice
			 {$y=\dfrac{x^2+x+2}{x+1}$}
			 {$y=\dfrac{x^2-x+2}{-x+1}$}
			 {$y=\dfrac{x^2-x-2}{x-1}$}
			 {\True $y=\dfrac{x^2-x+2}{x-1}$}}
			{\begin{tikzpicture}[line join = round,line cap = round,>=stealth,font=\footnotesize,scale=0.5]
		\def\xt{-3.5} \def\xp{5.5} \def\yd{-5.5}  \def\yt{5.5}
		 \def\a{1}
		\def\b{0} 
		 \def\c{2}
		 \def\d{1}
		\def\e{-1} 
		 \draw[->] (\xt,0)--(\xp,0) node[below]{$x$};
		\draw[->] (0,\yd)--(0,\yt)  node[left]{$y$};
		\foreach \x/\y/\label/\pos in {0/0/$O$/above left,1/0/$1$/above right,0/1/$1$/above  left,0/-2/$-2$/below left}
			{ 
		\fill  (\x,\y) circle (1pt) node [\pos]{\label};
		 }
		\clip (\xt+0.1,\yd+0.1) rectangle (\xp-0.1,\yt-0.1); 
		 \draw[samples=150,smooth,domain=\xt:0.99] plot (\x,{\a*\x+\b+\c/(\d*\x+\e)});
		\draw[samples=150,smooth,domain=1.01:\xp] plot (\x,{\a*\x+\b+\c/(\d*\x+\e)}); 
		 \draw[dashed,samples=150,smooth,domain=\xt:\xp] plot (\x,{\a*\x+\b});
		\draw[dashed]  (1,\yd) -- (1,\yt);
		\draw[dashed] (0,1) -- (1,1); 
	\end{tikzpicture}
			} 
	\loigiai{Vì đồ thị hàm số có đường tiệm cận đứng là $x=1$ và đường tiệm cận xiên là $y=x$ nên hàm số  có thể viết được dưới dạng: $y=x+\dfrac{c}{x-1}$.\\
			Vì đồ thị hàm số cắt trục tung tại  điểm $(0;-2)$ nên ta có: $\dfrac{c}{-1}=-2\Leftrightarrow c=2$.\\
			Suy ra $y=x+\dfrac{2}{x-1}$ hay  $y=\dfrac{x^2-x+2}{x-2}$.
		} 
	\end{ex}

% nguồn: Phan1_Lop_12_Dai_so_Chuong_1
\begin{ex}%[2D1N5-1]%[ma:c797edf5]
		Cho hàm số $y=\dfrac{ax+b}{cx+d}$ ($a < 0$) có đồ thị  như sau
		\begin{center}
			\begin{tikzpicture}[scale=0.7,font=\footnotesize,line join=round,line cap=round,>=stealth]
				\def\xmin{-5}; \def\ymin{-5}; \def\xmax{4}; \def\ymax{4}; 
				\coordinate  (O) at (0,0);
				\draw[->] (\xmin,0)--(\xmax,0)  node[below]{$x$};
				\draw[->]  (0,\ymin)--(0,\ymax) node[left]{$y$};
				\fill (O) node[below  left]{$O$} circle (1pt);
				\clip  ({\xmin-0.1},{\ymin-0.1}) rectangle ({\xmax+0.1},{\ymax+0.1});
				\draw[samples=100] plot[domain=\xmin:-1.01](\x,{(-\x+1)/(\x+1)}); 
				\draw[samples=100]  plot[domain=-0.99:\xmax](\x,{(-\x+1)/(\x+1)});
				\draw[dashed]  (-1.0,\ymin)--(-1.0,\ymax);
				\draw[dashed]  (\xmin,-1.0)--(\xmax,-1.0);
			\end{tikzpicture}
		\end{center}
		Mệnh đề nào sau  đây là đúng?
	\choice 
			{$ab<0$, $bc<0$, $cd > 0$} 
			{$ab>0$, $bc<0$, $cd  > 0$}
			{$ab>0$,  $bc>0$, $cd > 0$}
			{\True  $ab<0$, $bc>0$, $cd > 0$}
	 \loigiai{
			Dựa vào đồ thị  ta thấy
			\begin{itemize}
				\item Đồ thị hàm số có tiệm cận  đứng là $x=-\dfrac{d}{c}<0\Leftrightarrow \dfrac{d}{c}>0$. Suy ra, $d$ và $c$ cùng dấu.
				\item Đồ  thị hàm số có tiệm cận ngang là $y=\dfrac{a}{c}<0$. Suy ra, $a$ và $c$ trái dấu.
			\end{itemize}
			Mà theo giả thiết $a<0$ nên $c>0$ và  suy ra $d>0$.\\
			Đồ thị hàm  số cắt trục tung tại điểm $\left(0;\dfrac{b}{d} \right)$ và $\dfrac{b}{d}>0$ nên $b$ và $d$ cùng dấu dẫn đến $b>0$.\\
			Do đó  $ab<0$, $bc>0$, $cd>0$.
		 }
	\end{ex}

% nguồn: Phan1_Lop_12_Dai_so_Chuong_1
\begin{ex}%[2D1N5-1]%[ma:4015077f]
		Đồ thị hàm số nào dưới đây có dạng như đường cong  trong hình vẽ?
		\begin{center}
			\begin{tikzpicture}[scale=0.7,font=\footnotesize,line join=round,line cap=round,>=stealth]
				\def\xmin{-3}; \def\ymin{-3}; \def\xmax{5};  \def\ymax{4};
				 \coordinate (O) at (0,0);
				\draw[->] (\xmin,0)--(\xmax,0) node[below]{$x$}; 
				\draw[->] (0,\ymin)--(0,\ymax)  node[left]{$y$};
				\fill (O) node[below left]{$O$}  circle (1pt);
				\clip ({\xmin-0.1},{\ymin-0.1}) rectangle  ({\xmax+0.1},{\ymax+0.1});
				\draw[samples=100]  plot[domain=\xmin:0.99](\x,{(\x-2)/(\x-1)});
				\draw[samples=100]  plot[domain=1.01:\xmax](\x,{(\x-2)/(\x-1)});
				\draw[dashed]  (1.0,\ymin)--(1.0,\ymax);
				 \draw[dashed] (\xmin,1.0)--(\xmax,1.0);
				\draw  (1,3.5) node[right]{$x=1$}
					 (-2,0.2) node[above]{$y=1$};
			\end{tikzpicture}
		\end{center}
	\choice 
			 {\True $y=\dfrac{x-2}{x-1}$}
			 {$y=\dfrac{x-2}{x+1}$}
			 {$y=\dfrac{2x+1}{x-1}$}
			 {$y=-x^3+3x+2$}
	 \loigiai{
			Gọi hàm  số cần tìm có dạng $y=\dfrac{ax+b}{cx+d}$, với $a, b, c, d\in\mathbb{R}$ có độ thị $(C)$.\\
			Đồ thị có các đường tiệm cận đứng $x=1$ và tiệm cận ngang $y=1$  nên chỉ có hàm số $y=\dfrac{x-2}{x-1}$ thỏa mãn yêu cầu bài toán.
		 }
	\end{ex}

% nguồn: Phan1_Lop_12_Dai_so_Chuong_1
\begin{ex}%[2D1N5-1]%[ma:f28fcf67]
		Đồ thị của hàm số nào dưới đây có dạng như đường cong trong  hình dưới?
		\begin{center}
			\begin{tikzpicture}[scale=0.7,font=\footnotesize,line join=round,line cap=round,>=stealth]
				\draw[->]  (-2,0)--(4,0) node[below left] {$x$};
				\draw[->]  (0,-1.5)--(0,4) node[below left] {$y$};
				\draw (0,0) node [shift={(135:3mm)}]  {$O$};
				 \foreach \x/\g in {-1/-90,1/-90,2/-90}
				 \fill (\x,0) circle (1pt) node[shift={(\g:3mm)}]{\scriptsize $\x$};
				\foreach  \y/\g in {-1/-30,1/0,2/0,3/30}
				\fill  (0,\y) circle (1pt) node[shift={(\g:3mm)}]{\scriptsize $\y$};
				\draw[dashed,thin] 
					 (-1,0)|-(0,3)-|(2,0);
				\begin{scope}
					\clip (-2,-1.5) rectangle  (4,4);
					 \draw[smooth,samples=200,domain=-1.8:3.5,smooth,variable=\x] plot (\x,{-1*((\x)^3)+3*((\x)^2)-1});
				\end{scope}
			\end{tikzpicture}
		\end{center}
	\choice 
			{$y=-x^3+3x^2+1$} 
			 {$y=x^3+3x^2-1$}
			{\True $y=-x^3+3x^2-1$} 
			{$y=-x^3-3x-1$} 
	\loigiai{ 
			Dựa vào hàm số ta  thấy,\\
			+) Hình dáng đồ thị cho ta hệ số  $a<0$.\\
			+) Đồ thị cắt trục tung  tại $-1$ nên hệ số tự do $-1$.\\
			+) Đồ thị có hai cực trị nên $\Delta_{y'}  = b^2-3ac > 0$.\\
			Như vậy, hàm số  cần tìm là $y=-x^3+3x^2-1$.
		} 
	\end{ex}

% nguồn: Phan1_Lop_12_Dai_so_Chuong_1
\begin{ex}%[2D1N5-1]%[ma:2e290b25]
		Cho hàm số $f(x)=\dfrac{ax+1}{bx+c}$ ($a,b,c\in \mathbb{R}$) có bảng biến  thiên như sau
		\begin{center}
			\begin{tikzpicture}[scale=1,font=\footnotesize]
				 \tkzTabInit[nocadre=false,lgt=1.2,espcl=2.5,deltacl=0.7]
					 {$x$/0.8,$f'(x)$/0.6,$f (x)$/2}
					 {$-\infty$,$-1$,$+\infty$};
				\tkzTabLine{,-,d,-,}; 
				\tkzTabVar{+/$2$,-D+/$-\infty$/$+\infty$,-/$2$}; 
			\end{tikzpicture}
		\end{center}
		Trong  các số $a,b,c$ có bao nhiêu số dương?
	\choice 
			 {$2$}
			{$1$} 
			 {$0$}
			 {\True $3$}
	 \loigiai{
			Từ bảng biến  thiên ta thấy đồ thị hàm số có đường tiệm cận đứng là $x=-1$ và tiệm cận ngang $y=2$.\\
			Suy ra $\heva{&\dfrac{-c}{b}=-1 \\&\dfrac{a}{b}=2}  \Rightarrow \heva{&c=b \\&a=2b.}$\\
			Hàm số đã cho nghịch biến trên từng khoảng xác định  nên $y'=\dfrac{ac-b}{(bx+c)^2}<0$.\\
			Suy ra $ac-b < 0\Rightarrow 2b^2-b <  0\Rightarrow 0< b < \dfrac{1}{2}$.\\
			Suy ra cả ba số $a,b,c$ đều dương. 
		} 
	\end{ex}

% nguồn: Phan1_Lop_12_Dai_so_Chuong_1
\begin{ex}%[2D1N5-1]%[ma:90f98762]
		Cho  hàm số $y=\dfrac{ax-b}{x-1}$, với $a, b\in\mathbb{R}$ có đồ thị như hình vẽ.
		\begin{center}
			\begin{tikzpicture}[scale=0.7,font=\footnotesize,line join=round,line cap=round,>=stealth]
				\def\xmin{-4}; \def\ymin{-5};  \def\xmax{4}; \def\ymax{3};
				 \coordinate (O) at (0,0);
				 \draw[->] (\xmin,0)--(\xmax,0) node[below]{$x$};
				\draw[->] (0,\ymin)--(0,\ymax)  node[left]{$y$};
				\fill (O) node[below left]{$O$}  circle (1pt);
				\clip ({\xmin-0.1},{\ymin-0.1}) rectangle ({\xmax+0.1},{\ymax+0.1}); 
				\foreach \x  in {2}{
					\fill  (\x,0) node[below]{$\x$} circle (1pt);
				} 
				\foreach \y in  {-2}{
					\fill (0,\y)  node[left]{$\y$} circle (1pt);
				} 
				\draw[samples=100] plot[domain=\xmin:0.99](\x,{(-\x+2)/(\x-1)}); 
				\draw[samples=100]  plot[domain=1.01:\xmax](\x,{(-\x+2)/(\x-1)});
				\draw[dashed] (1.0,\ymin)--(1.0,\ymax); 
				\draw[dashed]  (\xmin,-1.0)--(\xmax,-1.0);
				\draw  (1,-4) node[right]{$x=1$}
					(-3,-0.9) node[above]{$y=-1$}; 
			\end{tikzpicture}
		\end{center}
		 Khẳng định nào dưới đây là đúng?
	 \choice
			{\True $b < a <  0$}
			{$0< b < a$} 
			{$b <  0< a$}
			{$0< a  < b$}
	\loigiai{ 
			Đồ thị hàm số có tiệm cận ngang là đường thẳng  $y=-1\Rightarrow\dfrac{a}{c}=-1\Rightarrow a=-1$.\\
			Đồ thị hàm số đi qua  điểm $(2;0)\Rightarrow 2a-b=0\Rightarrow b=-2$.\\ Vậy $b<a<0$.
		} 
	\end{ex}

% nguồn: Phan1_Lop_12_Dai_so_Chuong_1
\begin{ex}%[2D1N5-1]%[ma:cd95b799]
		Cho hàm số $y=\dfrac{ax-b}{x-1}$, với $a, b \in\mathbb{R}$ có đồ thị như hình  vẽ dưới đây
		\begin{center}
			\begin{tikzpicture}[scale=0.7,font=\footnotesize,line join=round,line cap=round,>=stealth]
				 \def\xmin{-4}; \def\ymin{-4}; \def\xmax{4}; \def\ymax{4};
				\coordinate  (O) at (0,0);
				\draw[->] (\xmin,0)--(\xmax,0) node[below]{$x$}; 
				\draw[->] (0,\ymin)--(0,\ymax)  node[left]{$y$};
				\fill (O) node[below left]{$O$}  circle (1pt);
				\clip ({\xmin-0.1},{\ymin-0.1})  rectangle ({\xmax+0.1},{\ymax+0.1});
				\foreach  \x in {2}{
					\fill (\x,0) node[below]{$\x$} circle  (1pt);
				} 
				\foreach \y  in {-2}{
					\fill  (0,\y) node[left]{$\y$} circle (1pt);
				 }
				 \draw[samples=100] plot[domain=\xmin:0.99](\x,{(-\x+2)/(\x-1)});
				\draw[samples=100]  plot[domain=1.01:\xmax](\x,{(-\x+2)/(\x-1)});
				 \draw[dashed] (1.0,\ymin)--(1.0,\ymax);
				\draw[dashed] (\xmin,-1.0)--(\xmax,-1.0); 
				 \draw (-3,-1) node[above]{$y=-1$}
					(1.0,\ymin) node[right]{$x=1$}; 
			\end{tikzpicture}
		\end{center}
		Tích $a.b$ bằng 
	\choice 
			{\True $2$} 
			{$-3$} 
			 {$-2$}
			 {$3$}
	\loigiai{ 
			Hàm số $y=\dfrac{ax-b}{x-1}$  có đồ thị là $(C)$.\\
			Từ đồ thị ta có  tiệm cận ngang $y=-1$ nên $a=-1$.\\
			Đồ thị hàm số cắt trục hoành tại điểm có hoành độ bằng $2$  nên $M\left(2;0\right)\in (C)$.\\
			 Suy ra, $0=\dfrac{-2-b}{2-1} \Leftrightarrow b=-2$.\\
			Vậy $a\cdot b=2$. 
		} 
	\end{ex}

% nguồn: Phan1_Lop_12_Dai_so_Chuong_1
\begin{ex}%[2D1N5-1]%[ma:6c9516a0]
		Cho hàm số $y=f(x)$ liên tục trên $[-3;2]$ và có bảng biến thiên như sau. Gọi $M$, $m$ lần lượt là giá trị lớn nhất và giá trị nhỏ nhất của hàm  số $y=f(x)$ trên đoạn $[-1;1]$. Tính $M+m$.
		\begin{center}
			\begin{tikzpicture}
				 \tkzTabInit[nocadre=true,lgt=1.2,espcl=2,deltacl=0.5]
					 {$x$/0.7,$f (x)$/2}
					 {$-3$,$-1$,$0$,$1$,$2$}
				 \tkzTabVar{-/$-2$,+/$3$,-/$0$,+/$2$,-/$1$}
			\end{tikzpicture}
		\end{center}
	\choice 
			{$0$} 
			{\True  $3$}
			 {$2$}
			{$1$} 
	\loigiai{Trên đoạn $[-1;1]$ ta có giá trị lớn nhất $M=3$ khi $x=-1$ và giá  trị nhỏ nhất $m=0$ khi $x=0$.\\
		 Khi đó $M=+m=3+0=3$.}
	\end{ex}

% nguồn: Phan1_Lop_12_Dai_so_Chuong_1
\begin{ex}%[2D1N5-1]%[ma:4778a72b]
	\immini{Đường cong hình bên là đồ thị của  một trong bốn hàm số dưới đây. Hàm số đó là hàm số nào?
	\choice 
			 {$y=\dfrac{x-2}{x-1}$}
			{\True  $y=x^3-3x+2$}
			{$y=-x^3+3x+2$} 
			 {$y=x^2+1$}}
			{\begin{tikzpicture}[line cap=butt,line join=miter,>=stealth,xscale=0.83,yscale=0.5]
		\tikzset{declare function={xmin=-3;  xmax=3;
			ymin=-3; ymax=5; 
			a=1; b=0;  c=-3; d=2;
			 f (\x)=a*(\x)^3+b*(\x)^2+c*(\x)+d;
			}, 
			 smooth,samples=50
		} 
		\draw[->] (xmin-0.25,0)--(xmax+0.5,0) 
			node[shift={(-100:7pt)},font=\normalsize]{$ x $}; 
		 \draw[->] (0,ymin-0.25)--(0,ymax+0.5)
			node[shift={(-5:7pt)},font=\normalsize]{$ y  $};
		\fill  (0,0) node[shift={(-45:9pt)},font=\normalsize]{$ O $};
		\begin{scope}
			\clip  (xmin,ymin) rectangle (xmax,ymax);
			\draw plot[domain=xmin:xmax] (\x,{f  (\x)});
		\end{scope}
	\end{tikzpicture}
			 }
	\loigiai{Ta thấy đường cong trên là đồ thị của hàm số bậc ba có hệ  số $a>0$.}
	\end{ex}

% nguồn: Phan1_Lop_12_Dai_so_Chuong_1
\begin{ex}%[2D1N5-1]%[ma:8b901ba0]
		Cho hàm số $y=\dfrac{ax+b}{cx+d}$, với  $a, b, c, d\in\mathbb{R}$ có đồ thị như hình bên dưới.
		\begin{center}
			\begin{tikzpicture}[scale=0.7,font=\footnotesize,line join=round,line cap=round,>=stealth]
				\def\xmin{-5}; \def\ymin{-4}; \def\xmax{4}; \def\ymax{5}; 
				\coordinate (O)  at (0,0);
				 \draw[->] (\xmin,0)--(\xmax,0) node[below]{$x$};
				\draw[->] (0,\ymin)--(0,\ymax) node[left]{$y$}; 
				\fill (O) node[below left]{$O$} circle  (1pt);
				\clip ({\xmin-0.1},{\ymin-0.1}) rectangle ({\xmax+0.1},{\ymax+0.1}); 
				\draw[samples=100]  plot[domain=\xmin:-1.01](\x,{(\x-1)/(\x+1)});
				\draw[samples=100] plot[domain=-0.99:\xmax](\x,{(\x-1)/(\x+1)}); 
				\draw[dashed]  (-1.0,\ymin)--(-1.0,\ymax);
				 \draw[dashed] (\xmin,1.0)--(\xmax,1.0);
			\end{tikzpicture}
		\end{center}
		Biết rằng $a$ là một số thực  dương, hỏi trong các số $b, c, d$ có tất cả bao nhiêu số dương?
	\choice 
			{$3$} 
			{$0$} 
			{\True $2$} 
			{$1$} 
	\loigiai{ 
			Đồ thị hàm số có đường tiệm cận ngang là  $y=\dfrac{a}{c} > 0$ nên $c > 0$.\\
			Đồ thì  hàm số có đường tiệm cận đứng là $x=-\dfrac{d}{c} < 0$ nên $d > 0$.\\
			Đồ thị hàm số cắt trục hoành tại điểm có hoành độ $x=-\dfrac{b}{a} > 0$ nên $b  < 0$.\\
			Vậy trong các số $b, c, d$ có $2$ số  dương.
		 }
	\end{ex}

% nguồn: Phan1_Lop_12_Dai_so_Chuong_1
\begin{ex}%[2D1N5-1]%[ma:753e2c47]
		Cho hàm số $y=\dfrac{ax+4-b}{cx+b}$, với $a, b\in\mathbb{R}$  có đồ thị là đường cong trong hình vẽ dưới đây.
		\begin{center}
			\begin{tikzpicture}[scale=0.7,font=\footnotesize,line join=round,line cap=round,>=stealth]
				 \def\xmin{-3}; \def\ymin{-5}; \def\xmax{5}; \def\ymax{3};
				\coordinate (O) at  (0,0);
				\draw[->] (\xmin,0)--(\xmax,0)  node[below]{$x$};
				\draw[->] (0,\ymin)--(0,\ymax)  node[left]{$y$};
				\fill (O) node[below left]{$O$} circle (1pt); 
				\clip ({\xmin-0.1},{\ymin-0.1})  rectangle ({\xmax+0.1},{\ymax+0.1});
				 \draw[samples=100] plot[domain=\xmin:0.99](\x,{(-2*\x+6)/(2*\x-2)});
				 \draw[samples=100] plot[domain=1.01:\xmax](\x,{(-2*\x+6)/(2*\x-2)});
				 \draw[dashed] (1.0,\ymin)--(1.0,\ymax);
				\draw[dashed] (\xmin,-1.0)--(\xmax,-1.0); 
			\end{tikzpicture}
		\end{center}
		Khẳng định nào sau đây là đúng? 
	\choice 
			{$a>0$, $b >  4$, $c < 0$}
			{\True $a>0$, $0< b  < 4$, $c < 0$}
			 {$a>0$, $b < 0$, $c < 0$}
			{$a<0$, $0< b < 4$,  $c < 0$}
	\loigiai{ 
			 Hàm số $y=\dfrac{ax+4-b}{cx+b}$ có đồ thị là đường cong $(C)$.\\
			Đồ thị có tiệm  cận ngang $y=\dfrac{a}{c}<0$.\\
			Đồ thị có tiệm cận đứng $x=\dfrac{-b}{c}>0\Leftrightarrow  \dfrac{b}{c}<0$.\\
			Đồ thị hàm số  cắt trục hoành tại điểm có hoành độ $x=\dfrac{b-4}{a}>0$.\\
			Đồ thị cắt trục tung tại điểm  có tung độ $y=\dfrac{4-b}{b}<0$.\\
			Suy ra $\heva{&{\dfrac{a}{c} < 0}  \\ &{-\dfrac{b}{c} > 0} \\& {\dfrac{4-b}{b} > 0} \\& {\dfrac{b-4}{a} < 0}}\Leftrightarrow \heva{& {0< b < 4} \\ &{a > 0} \\ &{c < 0}.}$
		 }
	\end{ex}

% nguồn: Phan1_Lop_12_Dai_so_Chuong_1
\begin{ex}%[2D1N5-1]%[ma:30865136]
		Cho hàm số $y=\dfrac{ax+b}{x+c}$ với $a,b,c\in  \mathbb{R}$ có đồ thị như hình vẽ dưới đây
		\begin{center}
			\begin{tikzpicture}[scale=0.7,font=\footnotesize,line join=round,line cap=round,>=stealth]
				 \def\xmin{-4}; \def\ymin{-5}; \def\xmax{4}; \def\ymax{3};
				\coordinate (O)  at (0,0);
				\draw[->] (\xmin,0)--(\xmax,0) node[below]{$x$}; 
				\draw[->]  (0,\ymin)--(0,\ymax) node[left]{$y$};
				 \fill (O) node[below left]{$O$} circle (1pt);
				\clip ({\xmin-0.1},{\ymin-0.1}) rectangle  ({\xmax+0.1},{\ymax+0.1});
				\foreach \x in  {2}{
					\fill  (\x,0) node[below]{$\x$} circle (1pt);
				 }
				\foreach \y  in {-2}{
					\fill (0,\y) node[left]{$\y$} circle (1pt); 
				 }
				\draw[samples=100]  plot[domain=\xmin:0.99](\x,{(-\x+2)/(\x-1)});
				\draw[samples=100]  plot[domain=1.01:\xmax](\x,{(-\x+2)/(\x-1)});
				\draw[dashed] (1.0,\ymin)--(1.0,\ymax); 
				\draw[dashed]  (\xmin,-1.0)--(\xmax,-1.0);
				\draw (1.0,-4)  node[right]{$x=1$}
					 (-3,-1.1) node[above]{$y=-1$};
			\end{tikzpicture}
		\end{center}
		 Giá trị của $a+2b+3c$ bằng
	 \choice
			{$2$} 
			 {\True $0$}
			{$6$} 
			{$-8$} 
	\loigiai{ 
			Tiệm cận đứng của đồ thị hàm số $x=-c$ suy ra $-c=1$ $\Leftrightarrow c=-1$.\\ 
			Tiệm cận ngang của đồ thị hàm số $y=a$ suy ra $a=-1$.\\ 
			Đồ thị cắt trục $Oy$ tại điểm  $\left(0;\dfrac{b}{c} \right)$ nên $\dfrac{b}{c}=-2$ mà $c=-1$ suy ra $b=2$.\\
			Vậy  $a+2b+3c=0$.
		 }
	\end{ex}

% nguồn: Phan1_Lop_12_Dai_so_Chuong_1
\begin{ex}%[2D1N5-1]%[ma:578aa23e]
		Cho hàm số $f(x)=\dfrac{ax-4}{bx+c}$ ($a,b,c\in \mathbb{R}$) có  bảng biến thiên như sau
		\begin{center}
			\begin{tikzpicture}[scale=1,font=\footnotesize]
				 \tkzTabInit[nocadre=false,lgt=1.2,espcl=2.5,deltacl=0.7]
					{$x$/0.8,$f'(x)$/0.6,$f  (x)$/2}
					 {$-\infty$,$1$,$+\infty$};
				 \tkzTabLine{,+,d,+,};
				 \tkzTabVar{-/$1$,+D-/$+\infty$/$-\infty$,+/$1$};
			\end{tikzpicture}
		\end{center}
		 Trong các số $a,b,c$ có bao nhiêu số dương?
	\choice 
			 {$3$}
			 {$4$}
			 {\True $2$}
			 {$1$}
	\loigiai{ 
			 Tiệm cận ngang $y=\dfrac{a}{b}=1\Rightarrow a=b \quad (1)$.\\
			Tiệm cận đứng $x=-\dfrac{c}{b}=1\Rightarrow c=-b\quad  (2)$.\\
			Ta  có $y'=\dfrac{ac+4b}{\left(bx+c\right)^2}>0\Rightarrow ac+4b>0\quad (3)$.\\
			 Từ (1), (2), (3) suy ra $-b^2+4b>0\Leftrightarrow 0<b<4$.\\
			Suy ra $a=b>0$, $c<0$. 
		 }
	\end{ex}

% nguồn: Phan1_Lop_12_Dai_so_Chuong_1
\begin{ex}%[2D1N5-1]%[ma:58959e0e]
		Bảng biến thiên sau đây của hàm  số nào?
		\begin{center}
			\begin{tikzpicture}[>=stealth]
				 \tkzTabInit[nocadre=false,lgt=1,espcl=2.5,deltacl=0.5]
					 {$x$/.7,$y'$/.7,$y$/2}
					{$-\infty$,$-1$,$+\infty$} 
				\tkzTabLine{,+,d,+,} 
				 \tkzTabVar{-/$2$,+D-/$+\infty$/$-\infty$,+/$2$}
			\end{tikzpicture}
		\end{center}
	\choice 
			{$y=\dfrac{2x+3}{x+1}$} 
			{$y=\dfrac{2x-1}{x-1}$} 
			{\True  $y=\dfrac{2x-1}{x+1}$}
			{$y=\dfrac{x+1}{2x-1}$} 
	\loigiai{ 
			Gọi hàm số cần tìm có  dạng $y=\dfrac{ax+b}{cx+d}$ có đồ thị $(C)$.\\
			Dựa vào bảng biến thiên, ta  thấy đồ thị $(C)$ có
			\begin{itemize}
				\item Tiệm cận ngang $y=\dfrac{a}{c}=2$. 
				\item Tiệm cận đứng  $x=\dfrac{-d}{c}=-1$.
				\item Hàm số đồng biến trên từng khoảng xác định nên $ad-bc>0$. 
			\end{itemize}
			Vậy hàm số  thỏa mãn là $y=\dfrac{2x-1}{x+1}$.
		} 
	\end{ex}

% nguồn: Phan1_Lop_12_Dai_so_Chuong_1
\begin{ex}%[2D1N5-1]%[ma:7b3122da]
		Cho hàm số  $y=\dfrac{ax+b}{x+1}$, với $a, b\in\mathbb{R}$ có đồ thị như hình vẽ dưới đây.
		\begin{center}
			\begin{tikzpicture}[scale=0.7,font=\footnotesize,line join=round,line cap=round,>=stealth]
				\def\xmin{-4};  \def\ymin{-3}; \def\xmax{3}; \def\ymax{4};
				\coordinate (O) at  (0,0);
				\draw[->] (\xmin,0)--(\xmax,0)  node[below]{$x$};
				\draw[->] (0,\ymin)--(0,\ymax) node[left]{$y$}; 
				\fill (O) node[below  left]{$O$} circle (1pt);
				\clip ({\xmin-0.1},{\ymin-0.1})  rectangle ({\xmax+0.1},{\ymax+0.1});
				\draw[samples=100]  plot[domain=\xmin:-1.01](\x,{(\x+2)/(\x+1)});
				 \draw[samples=100] plot[domain=-0.99:\xmax](\x,{(\x+2)/(\x+1)});
				 \draw[dashed] (-1.0,\ymin)--(-1.0,\ymax);
				\draw[dashed]  (\xmin,1.0)--(\xmax,1.0);
				\draw  (-1,3.5) node[left]{$x=-1$};
			\end{tikzpicture}
		\end{center}
		Khẳng định nào sau đây là  đúng?
	\choice 
			 {$0<a<b$}
			 {$b<0<a$}
			{$0<b<a$} 
			 {\True $a<b<0$}
	 \loigiai{
			 Hàm số $y=\dfrac{ax+b}{x+1}$ xác định trên $\mathscr{D}=\left(-\infty;-1\right)\cup \left(-1;+\infty \right)$.\\
			Ta có $y'=\dfrac{a-b}{\left(x+1\right)^2}$, $\forall x\ne-1$. Dựa  vào đồ thị ta thấy.
			\begin{itemize}
				\item Hàm số  có tiệm cận ngang là đường thẳng $y=a$. Dựa vào đồ thị ta có $a > 0$.\\
				\item Đồ thị hàm số nghịch biến trên khoảng $\left(-\infty;-1\right)$ và $\left(-1;+\infty \right)$ nên $y'  < 0$, $\forall x\ne-1\Leftrightarrow a-b<0\Leftrightarrow a<b\Rightarrow 0<a<b$.
			\end{itemize}
		 }
	\end{ex}

% nguồn: Phan1_Lop_12_Dai_so_Chuong_1
\begin{ex}%[2D1N5-1]%[ma:7064da16]
		Hàm số  nào có đồ thị như đường cong trong hình vẽ dưới đây?
		\begin{center}
			\begin{tikzpicture}[scale=0.7,font=\footnotesize,line join=round,line cap=round,>=stealth]
				\def\xmin{-3}; \def\ymin{-3}; \def\xmax{5};  \def\ymax{4};
				\coordinate (O) at (0,0); 
				\draw[->] (\xmin,0)--(\xmax,0) node[below]{$x$}; 
				 \draw[->] (0,\ymin)--(0,\ymax) node[left]{$y$};
				\fill (O)  node[below left]{$O$} circle (1pt);
				\clip  ({\xmin-0.1},{\ymin-0.1}) rectangle ({\xmax+0.1},{\ymax+0.1});
				\foreach \x in {2}{ 
					\fill (\x,0) node[below]{$\x$}  circle (1pt);
				 }
				\foreach  \y in {2}{
					\fill (0,\y) node[left]{$\y$} circle (1pt); 
				} 
				\draw[samples=100] plot[domain=\xmin:0.99](\x,{(\x-2)/(\x-1)}); 
				\draw[samples=100] plot[domain=1.01:\xmax](\x,{(\x-2)/(\x-1)}); 
				\draw[dashed]  (1.0,\ymin)--(1.0,\ymax);
				 \draw[dashed] (\xmin,1.0)--(\xmax,1.0);
				 \draw (1,3.5) node[right]{$x=1$}
					(-2,0.2)  node[above]{$y=1$};
			\end{tikzpicture}
		\end{center}
	 \choice
			 {$y=x^3-3x+2$}
			 {$y=\dfrac{x+2}{x-1}$}
			 {\True $y=\dfrac{x-2}{x-1}$}
			{$y=-x^4+5x^2-1$} 
	\loigiai{ 
			Gọi hàm số cần tìm là $y=\dfrac{ax+b}{cx+d}$, với $a, b, c, d\in\mathbb{R}$ có đồ thị là  $(C)$.\\
			Dựa vào  đồ thị ta có
			\begin{itemize}
				 \item Tiệm cận ngang của $(C)$ là $y=\dfrac{a}{c}=1$.
				\item Tiệm cận đứng của $(C)$ là $x=\dfrac{-d}{c}=1$. 
			\end{itemize}
			Mặt khác quan sát đồ thị thì hàm số  $y=f(x)$ đồng biến trên từng khoảng xác định nên $y=\dfrac{x-2}{x-1}$.
		 }
	\end{ex}

% nguồn: Phan1_Lop_12_Dai_so_Chuong_1
\begin{ex}%[2D1N5-1]%[ma:5e5a21f0]
		Hàm số nào dưới đây có đồ thị là đường cong trong hình vẽ? 
		\begin{center}
			\begin{tikzpicture}[scale=0.7,font=\footnotesize,line join=round,line cap=round,>=stealth]
				\def\xmin{-3}; \def\ymin{-2}; \def\xmax{4};  \def\ymax{6};
				 \coordinate (O) at (0,0);
				\draw[->]  (\xmin,0)--(\xmax,0) node[below]{$x$};
				\draw[->]  (0,\ymin)--(0,\ymax) node[left]{$y$};
				\fill (O) node[below left]{$O$} circle (1pt); 
				\clip ({\xmin-0.1},{\ymin-0.1}) rectangle  ({\xmax+0.1},{\ymax+0.1});
				\draw[samples=100] plot[domain=\xmin:0.99](\x,{(2*\x-3)/(\x-1)}); 
				\draw[samples=100] plot[domain=1.01:\xmax](\x,{(2*\x-3)/(\x-1)}); 
				\draw[dashed]  (1.0,\ymin)--(1.0,\ymax);
				\draw[dashed] (\xmin,2.0)--(\xmax,2.0); 
				\draw (1,5)  node[right]{$x=1$}
					(-2,1.2)  node[above]{$y=2$};
			\end{tikzpicture}
		\end{center}
	\choice 
			 {\True $y=\dfrac{2x-3}{x-1}$}
			 {$y=\dfrac{2x-1}{x-1}$}
			 {$y=\dfrac{x-3}{x-2}$}
			 {$y=\dfrac{2x+3}{x-1}$}
	\loigiai{ 
			Đồ thị hàm số có tiệm cận đứng $x=1$ và tiệm  cận ngang $y=2$.\\
			Từ hình vẽ ta được hàm số đồng biến trên  các khoảng xác định.
			\begin{itemize}
				\item Hàm số $y=\dfrac{2x-3}{x-1}$ có đạo hàm $y'=\dfrac{1}{\left(x-1\right)^2} > 0$,  $\forall x\ne 1$.
				\item Hàm số $y=\dfrac{2x-1}{x-1}$ có đạo hàm $y'=\dfrac{-1}{\left(x-1\right)^2} < 0$,  $\forall x\ne 1$.
				\item Hàm số $y=\dfrac{2x+3}{x-1}$ có đạo hàm $y'=\dfrac{-5}{\left(x-1\right)^2}  < 0$, $\forall x\ne 1$.
			\end{itemize}
			Do đó hàm số $y=\dfrac{2x-3}{x-1}$ thỏa  mãn bài toán.
		 }
	\end{ex}

% nguồn: Phan1_Lop_12_Dai_so_Chuong_1
\begin{ex}%[2D1N5-1]%[ma:ef33e84c]
		Cho hàm số $y=\dfrac{bx-c}{x-a}$ ($a\ne 0$ và $a,b,c\in \mathbb{R})$ có đồ thị như hình vẽ dưới đây. 
		\begin{center}
			\begin{tikzpicture}[scale=0.7,font=\footnotesize,line join=round,line cap=round,>=stealth]
				\def\xmin{-3};  \def\ymin{-3}; \def\xmax{5}; \def\ymax{4};
				\coordinate  (O) at (0,0);
				\draw[->] (\xmin,0)--(\xmax,0) node[below]{$x$}; 
				\draw[->] (0,\ymin)--(0,\ymax)  node[left]{$y$};
				\fill (O) node[below  right]{$O$} circle (1pt);
				\clip ({\xmin-0.1},{\ymin-0.1}) rectangle ({\xmax+0.1},{\ymax+0.1}); 
				\draw[samples=100]  plot[domain=\xmin:0.99](\x,{(\x+0.5)/(\x-1)});
				\draw[samples=100]  plot[domain=1.01:\xmax](\x,{(\x+0.5)/(\x-1)});
				\draw[dashed] (1.0,\ymin)--(1.0,\ymax); 
				 \draw[dashed] (\xmin,1.0)--(\xmax,1.0);
			\end{tikzpicture}
		\end{center}
		Khẳng định nào sau đây  đúng?
	\choice 
			{$a>0$, $b < 0$, $c-ab  < 0$}
			{\True $a > 0$, $b > 0$, $c-ab < 0$} 
			{$a < 0$, $b > 0$, $c-ab  < 0$}
			{$a < 0$,  $b < 0$, $c-ab > 0$}
	 \loigiai{
			Hàm số  $y=\dfrac{bx-c}{x-a}$ có đồ thị $(C)$.\\
			Dựa vào  đồ thị ta có
			\begin{itemize}
				\item Đồ thị có tiệm  cận ngang $y=b>0$.
				\item Đồ thị có tiệm cận  đứng $x=a>0$.
			\end{itemize}
			Vì hàm số $y=\dfrac{bx-c}{x-a}$ nghịch biến trên tập xác định nên $-ab+c < 0\Leftrightarrow c-ab <  0$.
		} 
	\end{ex}

% nguồn: Phan1_Lop_12_Dai_so_Chuong_1
\begin{ex}%[2D1N5-1]%[ma:49306f24]
	 \immini{Hàm số nào sau đây có đồ thị là đường cong như hình bên?
	\choice 
			{$y=x-\dfrac{1}{x-1}$} 
			{$y=-x+\dfrac{1}{x-1}$} 
			 {$y=-x-\dfrac{1}{x-1}$}
			{\True $y=x+\dfrac{1}{x-1}$}} 
			{\resizebox{!}{4cm}{ 
			\begin{tikzpicture}[line cap=butt,line join=miter,>=stealth,scale=0.7]
				 \draw[->] (-3,0)--(5,0)
					node[shift={(-100:7pt)},font=\normalsize]{$x$}; 
				\draw[->]  (0,-4)--(0,5)
					node[shift={(180:7pt)},font=\normalsize]{$y$}; 
				\draw  (5pt,0) |- (0,5pt)
					(0,0) node[shift={(145:11pt)},font=\normalsize]{$O$}; 
				\tikzset{declare function={f  (\x)=(\x)+1/((\x)-1);}}
				\foreach \x in  {2}{
					\draw[dashed] (\x,0) |- (0,{f  (\x)});
					\fill (\x,{f (\x)})  circle (1pt);
				 }
				\begin{scope}
					\clip (-3,-4)  rectangle (5,5);
					\draw[samples=100] plot[domain=-3:5]  (\x,{f (\x)});
					\draw[samples=50] plot[domain=-3:5] (\x,{\x}); 
					 \draw (1,-4)--(1,5);
					\draw[dashed,fill] (0,1)--(1,1)  circle (1pt);
				\end{scope}
				\foreach \x/\a in {1/below right,2/below} \fill (\x,0) circle (1pt) node [\a] {$\x$}; 
				\foreach \y/\a  in {-1/below left,1/above left,3/left} \fill (0,\y) circle (1pt) node [\a] {$\y$};
			\end{tikzpicture}
			 }}
	\loigiai{ 
			Ta có điểm $(2;3)$ thuộc đồ thị hàm số.\\ 
			Lần lượt thay $x=2$, $y=3$ vào các  phương án, ta có điểm $(2;3)$ chỉ thuộc đồ thị hàm số $y=x+\dfrac{1}{x-1}$.
		} 
	\end{ex}

% nguồn: Phan1_Lop_12_Dai_so_Chuong_1
\begin{ex}%[2D1N5-1]%[ma:3b743896]
		Hình vẽ dưới là đồ thị của hàm số $y=\dfrac{ax+b}{cx+d}$, với $a, b, c, d\in\mathbb{R}$. 
		\begin{center}
			\begin{tikzpicture}[scale=0.7,font=\footnotesize,line join=round,line cap=round,>=stealth]
				 \def\xmin{-4}; \def\ymin{-4}; \def\xmax{4}; \def\ymax{4};
				\coordinate (O)  at (0,0);
				 \draw[->] (\xmin,0)--(\xmax,0) node[below]{$x$};
				\draw[->]  (0,\ymin)--(0,\ymax) node[left]{$y$};
				\fill (O) node[below left]{$O$} circle (1pt); 
				\clip ({\xmin-0.1},{\ymin-0.1}) rectangle  ({\xmax+0.1},{\ymax+0.1});
				 \draw[samples=100] plot[domain=\xmin:-1.01](\x,{(\x-0.5)/(\x+1)});
				 \draw[samples=100] plot[domain=-0.99:\xmax](\x,{(\x-0.5)/(\x+1)});
				\draw[dashed] (-1.0,\ymin)--(-1.0,\ymax); 
				 \draw[dashed] (\xmin,1.0)--(\xmax,1.0);
			\end{tikzpicture}
		\end{center}
		Khẳng  định nào sau đây là đúng?
	 \choice
			{$bd>0$,  $ad>0$}
			{$bd<0$, $ab>0$} 
			{\True $ad>0$,  $ab<0$}
			{$ab<0$, $ad<0$} 
	 \loigiai{
			Đồ thị của hàm số $y=\dfrac{ax+b}{cx+d}$ là  $(C)$.\\
			Đồ thị có tiệm  cận ngang $y=\dfrac{a}{c}>0$. Suy ra, $a$ và $c$ cùng dấu.\\
			Đồ thị  có tiệm cận đứng $x=\dfrac{-d}{c}<0\Leftrightarrow \dfrac{d}{c}>0$. Suy ra, $d$ và $c$ cùng dấu.\\
			Dẫn đến $a$ và $d$ cùng  dấu.\\
			Đồ thị cắt trục hoành tại điểm $A\left(-\dfrac{b}{a};0\right)$ và  $-\dfrac{b}{a}>0\Leftrightarrow \dfrac{b}{a}<0$. Suy ra, $b$ và $a$ trái dấu. \\
			Vậy ta có được  $ad>0$ và $ab<0$.
		} 
	\end{ex}

% nguồn: Phan1_Lop_12_Dai_so_Chuong_1
\begin{ex}%[2D1N5-1]%[ma:0f62a56a]
		Đường cong  trong hình vẽ bên là đồ thị của hàm số nào?
		\begin{center}
			\begin{tikzpicture}[scale=0.7,font=\footnotesize,line join=round,line cap=round,>=stealth]
				 \def\xmin{-3}; \def\ymin{-3}; \def\xmax{5}; \def\ymax{4};
				\coordinate (O)  at (0,0);
				 \draw[->] (\xmin,0)--(\xmax,0) node[below]{$x$};
				\draw[->]  (0,\ymin)--(0,\ymax) node[left]{$y$};
				\fill (O) node[below right]{$O$} circle (1pt); 
				\clip ({\xmin-0.1},{\ymin-0.1}) rectangle ({\xmax+0.1},{\ymax+0.1}); 
				\draw[samples=100]  plot[domain=\xmin:0.99](\x,{(\x+1)/(\x-1)});
				\draw[samples=100] plot[domain=1.01:\xmax](\x,{(\x+1)/(\x-1)}); 
				\draw[dashed]  (1.0,\ymin)--(1.0,\ymax);
				\draw[dashed]  (\xmin,1.0)--(\xmax,1.0);
			\end{tikzpicture}
		\end{center}
	\choice 
			 {$y=x^4+x^2+1$}
			{\True $y=\dfrac{x+1}{x-1}$} 
			{$y=x^3-3x-1$} 
			{$y=\dfrac{2x+1}{x-1}$} 
	 \loigiai{
			Ta thấy đồ thị hàm số có một đường tiệm cận đứng $x=1$ và một đường tiệm cận ngang  $y=1$.\\
			Vậy  đường cong trong hình vẽ là đồ thị của hàm số $y=\dfrac{x+1}{x-1}$.
		} 
	\end{ex}

% nguồn: Phan1_Lop_12_Dai_so_Chuong_1
\begin{ex}%[2D1N5-1]%[ma:4c64ab12]
	\immini{ 
		Đồ thị của hàm  số nào dưới đây có dạng như đường cong trong hình sau.
		 }
		{ 
		\begin{tikzpicture}[scale=0.6,font=\footnotesize,line join=round,line cap=round,>=stealth]
			\tikzset{ethicknode/.style={scale=0.9}} 
			\draw[->] (-5.1,0)--(4.1,0) node[below  left] {$x$};
			 \draw[->] (0,-3.1)--(0,4.1) node[below left] {$y$};
			\draw (0,0) node [above left] {$O$}  circle (1pt);
			 \foreach \x/\nx in {-1/-1}
			\draw[thin]  (\x,1pt)--(\x,-1pt) node [below left] {$\nx$};
			 \foreach \y/\ny in {2/2}
			\draw[thin]  (1pt,\y)--(-1pt,\y) node [above left] {$\ny$};
			 \draw (-0.99,-3)--(-0.99,4);
			\begin{scope}
				\clip (-5,-3) rectangle (4,4); 
				\draw[samples=200,domain=-5:-1.01,smooth,variable=\x] plot  (\x,{(2*(\x)+-1)/(1*(\x)+1)});
				 \draw[samples=200,domain=-0.99:4,smooth,variable=\x] plot (\x,{(2*(\x)+-1)/(1*(\x)+1)});
				\draw (-5,2/1)--(4,2/1); 
			\end{scope}
		\end{tikzpicture}
		} 
	\choice 
			 {$y=\dfrac{2x+1}{x+1}$}
			 {$y=x^3-3x-1$}
			 {\True $y=\dfrac{2x-1}{x+1}$}
			 {$y=\dfrac{2x^2+3x+2}{x+1}$}
	 \loigiai{
			Quan sát đồ thị hàm số ta  thấy
			\begin{itemize}
				\item Đồ thị  dạng bậc nhất trên bậc nhất.
				\item Với $x=0\Rightarrow y(0)<0$. Suy ra hàm số cần tìm là  $y=\dfrac{2x-1}{x+1}$.
			\end{itemize}
		} 
	\end{ex}

% nguồn: Phan1_Lop_12_Dai_so_Chuong_1
\begin{ex}%[2D1N5-1]%[ma:bfad001b]
	\immini{ 
		Một  trong bốn hàm số sau đây có đồ thị như ở bên. Hàm số đó là hàm số nào?
	\choice 
			{\True  $y=x-1+\dfrac{1}{x-2}$}
			{$y=x-1-\dfrac{1}{x-2}$} 
			{$y=x-1+\dfrac{1}{2(x-2)}$} 
			{$y=x+\dfrac{1}{x-2}$} 
			}{\resizebox{!}{5cm}{\begin{tikzpicture}[scale=0.7,>=stealth]
		\draw[cyan!50,thin]  (-4.5,-4) grid (5,5);
		\draw[->] (-4,0)--(0,0) 
			 node[below right]{$O$}--(5,0) node[below]{$x$};
		\draw[->] (0,-4) --(0,6)  node[right]{$y$};
		 \foreach \x in {-4,-3,-2,-1,1,2,3,4}{
			\draw (\x,-.1)--(\x,.1) node[above]{$\x$}; 
			 }]
		\foreach \y in {-4,-3,-2,-1,1,2,3,4,5}{ 
			\draw  (-.1,\y)--(.1,\y) node[right]{$\y$};
		 }
		\draw[samples=100,domain=-3.1:1.8] plot  (\x,{(\x)-1+1/(\x-2)});
		\draw[samples=100,domain=2.27:4.9]  plot (\x,{(\x)-1+1/(\x-2)});
		\draw [domain=-3:4.9] plot  (\x,{(\x-1)});
		\draw (2,-4.1)--(2,5.1); 
		\draw[fill]  (2,1) circle (1pt);
	\end{tikzpicture}
			 }
			} 
	\loigiai{ 
			Quan sát đồ thị ấ thấy đây là  đồ thị của hàm số $y=\dfrac{ax^2+bx+c}{mx+n}$.\\
			Có $3$ hàm số có tiệm cận xiên và tiệm cận đứng thoả mãn với hình nên ta không sử dụng được để đưa ra đáp án chính xác ngay.\\ 
			 Thứ hai nữa là hàm số đi qua $2$ điểm $(3;3)$ và $(1;-1)$ nên chỉ có $y=x-1+\dfrac{1}{x-2}$ thoả mãn.
		 }
	\end{ex}

% nguồn: Phan1_Lop_12_Dai_so_Chuong_1
\begin{ex}%[2D1N5-1]%[ma:eb406350]
		Đường cong trong hình vẽ dưới đây là đồ thị của hàm  số nào trong các hàm số sau?
		\begin{center}
			\begin{tikzpicture}[scale=0.7,font=\footnotesize,line join=round,line cap=round,>=stealth]
				\def\xmin{-4}; \def\ymin{-4}; \def\xmax{4};  \def\ymax{4};
				 \coordinate (O) at (0,0);
				\draw[->] (\xmin,0)--(\xmax,0)  node[below]{$x$};
				\draw[->] (0,\ymin)--(0,\ymax) node[left]{$y$}; 
				\fill (O) node[below  right]{$O$} circle (1pt);
				\clip ({\xmin-0.1},{\ymin-0.1}) rectangle ({\xmax+0.1},{\ymax+0.1}); 
				 \draw[samples=100] plot[domain=\xmin:0.99](\x,{(\x+1)/(2*\x-2)});
				 \draw[samples=100] plot[domain=1.01:\xmax](\x,{(\x+1)/(2*\x-2)});
				 \draw[dashed] (1.0,\ymin)--(1.0,\ymax);
				 \draw[dashed] (\xmin,0.5)--(\xmax,0.5);
				 \draw (1.0,-3.5) node[right]{$x=1$}
					(-3,0.7)  node[above]{$y=\dfrac{1}{2}$};
			\end{tikzpicture}
		\end{center}
	 \choice
			{$y=\dfrac{x+1}{x-1}$} 
			{$y=\dfrac{2x-4}{x-1}$} 
			 {\True $y=\dfrac{x+1}{2x-2}$}
			 {$y=\dfrac{2x}{3x-3}$}
	\loigiai{ 
			Gọi hàm số cần tìm có dạng $y=\dfrac{ax+b}{cx+d}$, với $a, b, c, d\in\mathbb{R}$ có đồ thị $(C)$.\\ 
			Dựa vào đồ thị ta thấy đồ thị hàm số có tiệm  cận ngang là $y=\dfrac{1}{2}$.\\
			Do đó, $\lim\limits_{x\to -\infty}y=\dfrac{a}{c}=\dfrac{1}{2}$  và $\lim\limits_{x\to +\infty}y=\dfrac{a}{c}=\dfrac{1}{2}$.\\
			Vậy hàm số  cần tìm là $y=\dfrac{x+1}{2x-2}$.
		 }
	\end{ex}

% nguồn: Phan1_Lop_12_Dai_so_Chuong_1
\begin{ex}%[2D1N5-1]%[ma:d6797310]
		Đồ thị của hàm  số nào dưới đây có dạng như đường cong trong hình bên?
		\begin{center}
			\begin{tikzpicture}[line join=round,line cap=round,>=stealth]
				 \tikzset{ethicknode/.style={scale=0.9}}
				\draw[->] (-2.5,0)--(2.5,0) node[below left]  {$x$};
				\draw[->] (0,-2.5)--(0,2.5) node[below left] {$y$}; 
				\draw  (0,0) node [below left] {$O$};
				\begin{scope}
					\clip (-2.5,-2.5) rectangle (2.5,2.5); 
					 \draw[samples=200,domain=-2:2,smooth,variable=\x] plot (\x,{1*((\x)^3)+0*((\x)^2)+-3*(\x)+0});
				\end{scope}
			\end{tikzpicture}
		\end{center}
	\choice 
			{\True  $y=x^3-3x$}
			 {$y=-x^3+3x$}
			{$y=x^2-2x$} 
			 {$y=-2x^2+3$}
	\loigiai{ 
			Đường cong có dạng của đồ thị hàm số  bậc ba với hệ số $a>0$ nên chỉ có hàm số $y=x^3-3x$ thỏa yêu cầu bài toán.
		} 
	\end{ex}

% nguồn: Phan1_Lop_12_Dai_so_Chuong_1
\begin{ex}%[2D1N5-1]%[ma:8a0f98c9]
	\immini{Đồ thị của hàm số  nào dưới đây có dạng như đường cong trong hình dưới?
	\choice 
			 {$y=x^4-3x^2+1$}
			 {$y=x^4-2x^2+1$}
			 {\True$y=-x^4+2x^2+1$}
			{$y=2x^4-x^2+1$}}{ 
			\begin{tikzpicture}[>=stealth,x=1cm,y=1cm,scale=1]
				 \def\a{-1}
				 \def\b{2}
				 \def\c{1}
				\draw[->]  (-2,0) -- (2,0) node[below] {\scriptsize $x$};
				\draw[->] (0,-.5)  -- (0,2.25) node[left] {\scriptsize $y$};
				\draw[dashed] (0,0) node[below left]{\scriptsize  $O$}
					(1,0) node[below]{\scriptsize $1$}|-(0,2)  node[below left]{\scriptsize $2$}-|(-1,0) node[below]{\scriptsize $-1$}
					(0,1) node[below left]{\scriptsize  $1$};
				\clip (-2,-.5)  rectangle (2,2.25);
				\draw[samples=150,smooth,domain=-4:4]  plot (\x,{\a*(\x)^4+(\b)*(\x)^2+(\c)});
				\fill  (0,0) circle (1.25pt);
				 \fill (-1,2) circle (1.25pt);
				\fill (1,2) circle (1.25pt); 
				\fill (0,1) circle  (1.25pt);
			\end{tikzpicture}
			} 
	 \loigiai{
			Dựa vào đồ thị  hàm số ta thấy\\
			+ Hình dáng đồ  thị cho ta hệ số $a<0$.\\
			+ Phương trình $y'=0$ có 3  nghiệm phân biệt $x_1=-1$, $x_2=0$, $x_3=1$.\\
			Như vậy,  hàm số cần tìm là $y=x^4-2x^2+1$.
		 }
	\end{ex}

% nguồn: Phan1_Lop_12_Dai_so_Chuong_1
\begin{ex}%[2D1N5-1]%[ma:0ef1f960]
		Cho hàm số $y=\dfrac{ax+b}{cx+d}$, với $a, b,  c, d\in\mathbb{R}$ có đồ thị như hình vẽ bên dưới.
		\begin{center}
			\begin{tikzpicture}[scale=0.7,font=\footnotesize,line join=round,line cap=round,>=stealth]
				\def\xmin{-4}; \def\ymin{-3};  \def\xmax{3}; \def\ymax{4};
				\coordinate (O)  at (0,0);
				\draw[->] (\xmin,0)--(\xmax,0) node[below]{$x$}; 
				\draw[->] (0,\ymin)--(0,\ymax) node[left]{$y$}; 
				\fill (O) node[below left]{$O$} circle (1pt); 
				\clip ({\xmin-0.1},{\ymin-0.1}) rectangle ({\xmax+0.1},{\ymax+0.1}); 
				\draw[samples=100] plot[domain=\xmin:-1.01](\x,{(-0.5*\x+0.2)/(-\x-1)}); 
				\draw[samples=100] plot[domain=-0.99:\xmax](\x,{(-0.5*\x+0.2)/(-\x-1)}); 
				 \draw[dashed] (-1.0,\ymin)--(-1.0,\ymax);
				\draw[dashed] (\xmin,0.5)--(\xmax,0.5); 
			\end{tikzpicture}
		\end{center}
		Khẳng  định nào sau đây là khẳng định đúng?
	 \def\dotEX{}
	\choice 
			{$\heva{&ad<0\\&bc<0.}$} 
			 {$\heva{&ad<0\\&bc>0.}$}
			 {\True $\heva{&ad>0\\&bc<0.}$}
			{$\heva{&ad>0\\&bc>0.}$} 
	\loigiai{ 
			 Đồ thị có tiệm cận ngang $y=\dfrac{a}{c}>0\Rightarrow ac>0\quad (1)$.\\
			Đồ thị có  tiệm cận đứng $x=-\dfrac{d}{c}<0\Leftrightarrow \dfrac{d}{c} > 0\Rightarrow cd > 0\quad (2)$.\\
			Đồ thị cắt trục $Oy$ tại điểm có tung độ âm $y=\dfrac{b}{d}<0\Rightarrow bd < 0\quad  (3)$.\\
			Đồ thị cắt trục $Ox$ tại  điểm có hoành độ dương $x=-\dfrac{b}{a} > 0\Leftrightarrow \dfrac{b}{a} < 0\Rightarrow ab < 0\quad (4)$.\\
			Lấy $(3)\cdot  (4)\Rightarrow ad\cdot b^2> 0\Rightarrow ad>0\quad (5)$.\\
			Lấy $(1)\cdot  (4)\Rightarrow bc\cdot a^2< 0\Rightarrow bc<0\quad (6)$.
		} 
	\end{ex}

% nguồn: Phan1_Lop_12_Dai_so_Chuong_1
\begin{ex}%[2D1H5-1]%[ma:fdbaff57]
		Hàm số nào dưới đây có  đồ thị như hình vẽ?
		\begin{center}
			\begin{tikzpicture}[line cap=round,line join=round,font=\footnotesize,>=stealth,scale=1,x=0.5cm,y=0.5cm]
				\tikzset{label  style/.style={font=\footnotesize}}
				\draw[->] (-4,0)--(6.1,0)  node[below] {$x$};
				\draw[->] (0,-5)--(0,4)  node[left] {$y$};
				\draw (0,0) node [above  left] {\scriptsize$O$};
				\foreach \x/\nx in  {-4/-4,-3/-3,-2/-2,-1/,1/1,2/2,3/3,4/4,5/5}
				\draw[thin] (\x,1pt)--(\x,-1pt)  node [below] {\scriptsize$\nx$};
				\foreach \y/\ny  in {-4/left,-3/left,-2/left,-1/left,1/left,2/left,3/left}
				\draw[thin]  (1pt,\y)--(-1pt,\y) node [\ny] {\scriptsize $\y$};
				\draw[dashed,thin]  (1,-5)--(1,4) node[above,rotate=90,pos=0.1]{$x=1$};
				\draw (-1,0) node[above left]{$-1$}; 
				\begin{scope}
					 \clip (-6,-5) rectangle (6,5);
					\draw[samples=200,domain=-3:0.85,smooth,variable=\x,magenta]  plot (\x,{(1*((\x)^2)+-4*(\x)+2)/(1*(\x)+-1)});
					\draw[samples=200,domain=1.1:6,smooth,variable=\x,magenta] plot  (\x,{(1*((\x)^2)+-4*(\x)+2)/(1*(\x)+-1)});
					\draw[dashed,thin]  (-5,-8)--(6,3) node[above,rotate=45,pos=0.8]{$y=x-3$};
				\end{scope}
			\end{tikzpicture}
		\end{center}
	\choice 
			 {$y=\dfrac{x^2-3}{x+1}$}
			 {\True $y=\dfrac{x^2-4x+2}{x-1}$}
			{$y=\dfrac{x^2-x}{x-1}$} 
			 {$y=\dfrac{x^2-4x+5}{x+1}$}
	\loigiai{ 
			Từ đồ thị  ta thấy:
			\begin{itemize}
				\item Tiệm cận đứng là đường  thẳng $x=1$ và tiệm cận xiên là đường thẳng $y=x-3$.
				\item Giới hạn $\lim\lim\limits_{x\to-\infty}y=-\infty; \lim\limits_{x  \rightarrow+\infty}y=+\infty$.
			\end{itemize}
			Do đó hàm số  cần tìm là $y=\dfrac{x^2-4x+2}{x-1}$.
		} 
	\end{ex}

% nguồn: Phan1_Lop_12_Dai_so_Chuong_1
\begin{ex}%[2D1H5-1]%[ma:7ea9f886]
		Cho hàm số $y = \dfrac{ax - b}{x - 1}$ có đồ thị như hình  vẽ.
		\begin{center}
			\begin{tikzpicture}[scale=0.8,>=stealth,font=\footnotesize,line join=round,line cap=round]
				\draw[->] (-3,0)  -- (5,0) node[below] {$x$};
				\draw[->] (0,-4) --  (0,3) node[left] {$y$};
				\draw[dashed] (-3,-1) -- (5,-1); 
				\node  at (-2.5,-0.7) {\footnotesize $y=-1$};
				\draw[dashed] (1,-4) --  (1,3);
				\node[left] at (1.1,2.7) {\footnotesize $x=1$}; 
				\draw[domain=-3:0.69,smooth] plot  (\x,{-1-1/(1-\x)});
				\draw[domain=1.26:4.8,smooth] plot  (\x,{-1-1/(1-\x)});
				\node at (0,0)  [below left] {\footnotesize $O$};
			\end{tikzpicture}
		\end{center}
		Khẳng định nào  dưới đây là đúng?
	\choice 
			{$b < 0  < a$}
			{$0 < b  < a$}
			{\True $b  < a < 0$}
			 {$0 < a < b$}
	 \loigiai{Tập xác định $\mathscr{D}=\mathbb{R}\setminus\{1\}$.\\
			Ta có $y = \dfrac{ax - b}{x -  1} \Rightarrow y' = \dfrac{-a + b}{(x - 1)^2}$.\\
			Từ đồ thị suy ra hàm số nghịch biến trên mỗi khoảng xác định nên $y' <  0 \,\forall x\ne 1$\\
			$$ \Rightarrow -a  + b < 0 \Leftrightarrow a > b.$$
			Mặt khác, đồ thị có tiệm cận  ngang $y = -1$ nên $ a = -1 \Rightarrow a < 0$.\\
			Vậy  $b < a < 0$.
		} 
	\end{ex}

% nguồn: Phan1_Lop_12_Dai_so_Chuong_1
\begin{ex}%[2D1H5-1]%[ma:1a51a5e4]
		Hàm số  nào dưới đây có bảng biến thiên như hình vẽ?
		\begin{center}
			\begin{tikzpicture}[line join=round,font=\footnotesize,>=stealth,scale=1]
				 \tkzTabInit[nocadre=false,lgt=1.2,espcl=2.5,deltacl=0.6]
					{$x$  /1,$y'$ /1,$y$ /3}
					{$-\infty$,$1$,$3$,$5$,$+\infty$} 
				\tkzTabLine{,-,0,+,d,+,0,-} 
				 \tkzTabVar{+/$+\infty$,-/$-1$,+D-/$+\infty$/$-\infty$,+/$-9$,-/$-\infty$}
			\end{tikzpicture}
		\end{center}
	 \choice
			 {$y=\dfrac{x^2-4x+3}{x-3}$}
			 {$y=\dfrac{-x^2-x+2}{x-3}$}
			{\True  $y=\dfrac{-x^2+x+2}{x-3}$}
			{$y=\dfrac{x^2-4x+4}{-x+3}$} 
	\loigiai{ 
			Dựa vào bảng  biến thiên, ta thấy:
			\begin{itemize}
				\item Tiệm cận đứng là đường thẳng $x=3$, các giới hạn  $\lim\limits_{x \rightarrow+\infty}=-\infty; \lim\limits_{x \rightarrow-\infty}=+\infty$.
				\item  Điểm cực đại $A(5; -9)$ và điểm cực tiểu $B(1; -1)$.
			\end{itemize}
			Do đó hàm số cần tìm là $y=\dfrac{-x^2+x+2}{x-3}$. 
		} 
	\end{ex}

% nguồn: Phan1_Lop_12_Dai_so_Chuong_1
\begin{ex}%[2D1H5-1]%[ma:a8654077]
		Cho hàm số $y = f(x) = ax^3 + bx^2 + cx  + d$ có bảng biến thiên sau
		\begin{center}
			\begin{tikzpicture}[scale=1,font=\footnotesize,line join=round,line cap=round,>=stealth]
				\tkzTabInit[nocadre=false,lgt=1.2,espcl=2.5,deltacl=0.6] 
					{$x$/0.6,$y'$/0.6,$y$/2} 
					 {$-\infty$,$-1$,$1$,$+\infty$}
				 \tkzTabLine{,+,0,-,0,+}
				 \path
					($(N13)!0.1!(N12)$) node  (A1){$-\infty$}
					($(N23)!0.8!(N22)$) node (A2){$2  $}
					 ($(N33)!0.1!(N32)$) node (A3){$-2 $}
					 ($(N43)!0.8!(N42)$) node (A4){$+\infty$};
				\foreach \x/\y  in {A1/A2,A2/A3,A3/A4}{
					\draw[-stealth] (\x)--(\y); 
				 }
			\end{tikzpicture}
		\end{center}
		Đồ thị nào  thể hiện hàm số $y=f(x)$?
	 \choice
			{\True \begin{tikzpicture}[line cap=butt,line join=miter,>=stealth,scale=0.7]
		\tikzset{declare function={xmin=-2.5; xmax=2.5;  ymin=-2.5; ymax=2.5;
			f (\x)=1*(\x)^3-3*(\x); 
			}, 
			smooth,samples=450 
		 }
		\draw[->] (xmin,0)--(xmax,0) node[shift={(-100:7pt)},font=\normalsize]{$  x $};
		\draw[->] (0,ymin)--(0,ymax)  node[shift={(190:7pt)},font=\normalsize]{$ y $};
		\fill (0,0)  node[shift={(225:7pt)},font=\normalsize]{$ O $};
		 \draw[dashed] (-1,0) node[below]{$-1$}--(-1,2)--(0,2) node[right]{$2$}
			 (1,0) node[above]{$1$}--(1,-2)--(0,-2) node[left]{$-2$};
		 \clip (xmin,ymin) rectangle (xmax,ymax);
		\draw plot[domain=xmin:xmax] (\x,{f (\x)}); 
	\end{tikzpicture}
			} 
			{\begin{tikzpicture}[line cap=butt,line join=miter,>=stealth,scale=0.7]
		\tikzset{declare function={xmin=-2.5; xmax=2; ymin=-1; ymax=4.5; 
			f  (\x)=1*(\x)^3-3*(\x)+2;
			}, 
			 smooth,samples=450
		 }
		\draw[->] (xmin,0)--(xmax,0) node[shift={(-100:7pt)},font=\normalsize]{$  x $};
		\draw[->]  (0,ymin)--(0,ymax) node[shift={(190:7pt)},font=\normalsize]{$ y $};
		\fill (0,0) node[shift={(225:7pt)},font=\normalsize]{$ O $}; 
		\draw[dashed] (-1,0) node[below]{$-1$}--(-1,4)--(0,4)  node[right]{$4$}
			(1,0)  node[below]{$1$}
			 (0,2) node[right]{$2$};
		\clip (xmin,ymin)  rectangle (xmax,ymax);
		\draw plot[domain=xmin:xmax]  (\x,{f (\x)});
	\end{tikzpicture}
			 }
			{\begin{tikzpicture}[line cap=butt,line join=miter,>=stealth,scale=0.7]
		\tikzset{declare function={xmin=-2.5; xmax=2; ymin=-4.5; ymax=1; 
			 f (\x)=-1*(\x)^3+3*(\x)-2;
			}, 
			smooth,samples=450 
		} 
		\draw[->]  (xmin,0)--(xmax,0) node[shift={(-100:7pt)},font=\normalsize]{$ x $};
		 \draw[->] (0,ymin)--(0,ymax) node[shift={(190:7pt)},font=\normalsize]{$ y $};
		\fill (0,0) node[shift={(225:7pt)},font=\normalsize]{$ O  $};
		\draw[dashed] (-1,0)  node[above]{$-1$}--(-1,-4)--(0,-4) node[right]{$-4$}
			 (1,0) node[above]{$1$}
			 (0,-2) node[right]{$-2$};
		\clip (xmin,ymin) rectangle  (xmax,ymax);
		 \draw plot[domain=xmin:xmax] (\x,{f (\x)});
	\end{tikzpicture}
			} 
			{\begin{tikzpicture}[line cap=butt,line join=miter,>=stealth,scale=0.7]
		\tikzset{declare function={xmin=-2.5;  xmax=2.5; ymin=-2.5; ymax=2.5;
			f (\x)=-1*(\x)^3+3*(\x); 
			}, 
			 smooth,samples=450
		 }
		\draw[->] (xmin,0)--(xmax,0) node[shift={(-100:7pt)},font=\normalsize]{$ x  $};
		\draw[->] (0,ymin)--(0,ymax)  node[shift={(190:7pt)},font=\normalsize]{$ y $};
		\fill (0,0) node[shift={(225:7pt)},font=\normalsize]{$ O  $};
		\draw[dashed] (-1,0)  node[above]{$-1$}--(-1,-2)--(0,-2) node[right]{$-2$}
			(1,0)  node[below]{$1$}--(1,2)--(0,2) node[left]{$2$};
		\clip (xmin,ymin) rectangle  (xmax,ymax);
		\draw plot[domain=xmin:xmax] (\x,{f  (\x)});
	\end{tikzpicture}
			} 
	\loigiai{ 
			Dựa vào  bảng biến thiên, ta thấy:\\
			Hàm số có giá trị cực đại bằng $2$ và giá trị cực tiểu bằng $-2$  nên loại
			\begin{center}
				\begin{tikzpicture}[line cap=butt,line join=miter,>=stealth,scale=0.7]
					\tikzset{declare function={xmin=-2.5;  xmax=2; ymin=-1; ymax=4.5;
						f  (\x)=1*(\x)^3-3*(\x)+2;
						 },
						smooth,samples=450 
					} 
					\draw[->] (xmin,0)--(xmax,0) node[shift={(-100:7pt)},font=\normalsize]{$ x $}; 
					 \draw[->] (0,ymin)--(0,ymax) node[shift={(190:7pt)},font=\normalsize]{$ y $};
					\fill  (0,0) node[shift={(225:7pt)},font=\normalsize]{$ O $};
					\draw[dashed] (-1,0) node[below]{$-1$}--(-1,4)--(0,4)  node[right]{$4$}
						 (1,0) node[below]{$1$}
						(0,2)  node[right]{$2$};
					 \clip (xmin,ymin) rectangle (xmax,ymax);
					\draw  plot[domain=xmin:xmax] (\x,{f (\x)});
				\end{tikzpicture}

				\hspace{2cm} 
				\begin{tikzpicture}[line cap=butt,line join=miter,>=stealth,scale=0.7]
					\tikzset{declare  function={xmin=-2.5; xmax=2; ymin=-4.5; ymax=1;
						f (\x)=-1*(\x)^3+3*(\x)-2; 
						 },
						 smooth,samples=450
					} 
					\draw[->] (xmin,0)--(xmax,0)  node[shift={(-100:7pt)},font=\normalsize]{$ x $};
					\draw[->] (0,ymin)--(0,ymax) node[shift={(190:7pt)},font=\normalsize]{$ y $}; 
					\fill (0,0) node[shift={(225:7pt)},font=\normalsize]{$ O $}; 
					\draw[dashed] (-1,0) node[above]{$-1$}--(-1,-4)--(0,-4) node[right]{$-4$} 
						(1,0) node[above]{$1$} 
						(0,-2) node[right]{$-2$}; 
					\clip  (xmin,ymin) rectangle (xmax,ymax);
					\draw plot[domain=xmin:xmax]  (\x,{f (\x)});
				\end{tikzpicture}
			\end{center}
			Khi  $x \to+\infty$ thì $y \to+\infty$ nên chỉ có chọn
			\begin{center}
				\begin{tikzpicture}[line cap=butt,line join=miter,>=stealth,scale=0.7]
					\tikzset{declare function={xmin=-2.5;  xmax=2.5; ymin=-2.5; ymax=2.5;
						f (\x)=1*(\x)^3-3*(\x); 
						 },
						 smooth,samples=450
					 }
					 \draw[->] (xmin,0)--(xmax,0) node[shift={(-100:7pt)},font=\normalsize]{$ x $};
					\draw[->] (0,ymin)--(0,ymax) node[shift={(190:7pt)},font=\normalsize]{$ y  $};
					\fill (0,0) node[shift={(225:7pt)},font=\normalsize]{$ O $}; 
					\draw[dashed] (-1,0) node[below]{$-1$}--(-1,2)--(0,2)  node[right]{$2$}
						(1,0) node[above]{$1$}--(1,-2)--(0,-2)  node[left]{$-2$};
					\clip (xmin,ymin)  rectangle (xmax,ymax);
					\draw plot[domain=xmin:xmax]  (\x,{f (\x)});
				\end{tikzpicture}
			\end{center}
		} 
	\end{ex}

% nguồn: Phan1_Lop_12_Dai_so_Chuong_1
\begin{ex}%[2D1H5-1]%[ma:c5ff4b47]
		Một trong bốn hàm số sau đây, có một hàm số mà đồ thị nhận $Oy$ là một trục đối  xứng, hàm số đó là hàm số nào?
	\choice 
			{$y=\sin  x$}
			{$y=\tan x$} 
			 {\True $y=\cos x$}
			 {$y=\cot x$}
	\loigiai{ 
			Hàm số nhận $Oy$ là một trục  đối xứng tức là hàm số chẵn. Trong các hàm số trên thì chỉ $y=\cos x$ thoả mãn.
		} 
	\end{ex}

% nguồn: Phan1_Lop_12_Dai_so_Chuong_1
\begin{ex}%[2D1H5-1]%[ma:6f186bbb]
		Hàm số nào  sau đây có đồ thị như hình vẽ?
		\begin{center}
			\begin{tikzpicture}[line cap=round,line join=round,font=\footnotesize,>=stealth,scale=1,x=0.5cm,y=0.5cm]
				\tikzset{label  style/.style={font=\footnotesize}}
				\draw[->] (-4,0)--(5,0) node[below] {$x$}; 
				 \draw[->] (0,-4)--(0,5) node[left] {$y$};
				\draw (0,0) node [above  left] {\scriptsize$O$};
				\foreach  \x/\nx in {-3/-3,-2/-2,-1/,1/1,2/2,3/3,4/4}
				\draw[thin] (\x,1pt)--(\x,-1pt) node [below]  {\scriptsize$\nx$};
				\foreach \y/\ny in {-3/left,-2/left,-1/left,1/left,2/right,3/right,4/right} 
				\draw[thin] (1pt,\y)--(-1pt,\y) node [\ny] {\scriptsize  $\y$};
				\draw[dashed,thin] (-1,-4)--(-1,5)  node[above,rotate=90,pos=0.1]{$x=-1$};
				\draw (-1,0)  node[above left]{$-1$};
				\begin{scope}
					\clip  (-6,-5) rectangle (6,5);
					\draw[samples=200,domain=-4:-1.2,smooth,variable=\x,magenta] plot (\x,{(-1*((\x)^2)+-(\x)+2)/(1*(\x)+1)}); 
					\draw[samples=200,domain=-0.9:4,smooth,variable=\x,magenta] plot  (\x,{(-1*((\x)^2)+-1*(\x)+2)/(1*(\x)+1)});
					\draw[dashed,thin] (-4,4)--(4,-4)  node[below,rotate=-45,pos=0.8]{$y=-x$};
				\end{scope}
			\end{tikzpicture}
		\end{center}
	\choice 
			 {$y=\dfrac{x^2-x+4}{x+1}$}
			{\True  $y=\dfrac{x^2-2x+3}{x+1}$}
			 {$y=\dfrac{-x^2-x+2}{x+1}$}
			 {$y=\dfrac{x^2+x-1}{x+1}$}
	 \loigiai{
			Dựa vào đồ thị, ta  thấy
			\begin{itemize}
				\item Tiệm cận đứng là đường thẳng $x=-1$, tiệm cận xiên là đường thẳng $y=-x$. 
				 \item Giới hạn $\lim\limits_{x \rightarrow-\infty}y=+\infty; \lim\limits_{x \rightarrow+\infty}y=-\infty$.
			\end{itemize}
			Do đó  hàm số cần tìm là $y=\dfrac{-x^2-x+2}{x+1}$.
		 }
	\end{ex}

% nguồn: Phan1_Lop_12_Dai_so_Chuong_1
\begin{ex}%[2D1H5-1]%[ma:64151c4b]
		Hàm số nào sau đây có bảng biến thiên  như hình vẽ?
		\begin{center}
			\begin{tikzpicture}[line join=round,font=\footnotesize,>=stealth,scale=1]
				 \tkzTabInit[nocadre=false,lgt=1.2,espcl=2.5,deltacl=0.6]
					{$x$  /1,$y'$ /1,$y$ /3}
					{$-\infty$,$-9$,$-4$,$1$,$+\infty$} 
				\tkzTabLine{,+,0,-,d,-,0,+} 
				\tkzTabVar{-/$-\infty$,+/$-20$,-D+/$-\infty$/$+\infty$,-/$0$,+/$+\infty$} 
			\end{tikzpicture}
		\end{center}
	 \choice
			{\True $y=\dfrac{x^2-2x+1}{x+4}$} 
			{$y=\dfrac{x^2-4x+2}{x+4}$} 
			{$y=\dfrac{x^2-x+2}{-x-4}$} 
			 {$y=\dfrac{x^2-3x+4}{-x-4}$}
	\loigiai{ 
			Dựa vào bảng biến thiên, ta  thấy:
			\begin{itemize}
				\item Tiệm cận đứng là đường thẳng $x=-4$, các  giới hạn $\lim\limits_{x \rightarrow+\infty}=+\infty; \lim\limits_{x \rightarrow-\infty}=-\infty$.
				\item  Điểm cực đại $A(-9; -20)$ và điểm cực tiểu $B(1; 0)$.
			\end{itemize}
			Do đó hàm số  cần tìm là $y=\dfrac{x^2-2x+1}{x+4}$.
		} 
	\end{ex}

% nguồn: Phan1_Lop_12_Dai_so_Chuong_1
\begin{ex}%[2D1H5-1]%[ma:258bb9f5]
		Hàm số nào  dưới đây có bảng biến thiên như hình bên dưới?
		\begin{center}
			\begin{tikzpicture}
				 \tkzTabInit[espcl=3,lgt=1.5]
					 {$x$/0.7,$y'$/0.7,$y$/2.1}
					{$-\infty$,$5$,$+\infty$} 
				 \tkzTabLine{,-,d,-,}
				\tkzTabVar{+/$1$,-D+/$-\infty$/$+\infty$,-/$1$} 
			\end{tikzpicture}
		\end{center}
	\choice 
			 {$y=\dfrac{x-7}{x-5}$}
			 {$y=\dfrac{x-5}{x-1}$}
			{$y=\dfrac{5x}{x-1}$} 
			{\True $y=\dfrac{x-3}{x-5}$} 
	 \loigiai{
			Theo bảng biến thiên ta có tiệm cận đứng có phương trình $x=5$, tiệm cận ngang có phương  trình $y=1$ và $y'<0, \forall x \in \mathbb{R}$ nên ta chọn hàm số $y=\dfrac{x-3}{x-5}$.
		} 
	\end{ex}

% nguồn: Phan1_Lop_12_Dai_so_Chuong_1
\begin{ex}%[2D1H5-1]%[ma:8889bed9]
		 Hàm số nào dưới đây có bảng biến thiên như hình vẽ?
		\begin{center}
			\begin{tikzpicture}[line join=round,font=\footnotesize,>=stealth,scale=1]
				\tkzTabInit[nocadre=false,lgt=1.2,espcl=2.5,deltacl=0.6] 
					{$x$ /1,$y'$ /1,$y$ /3} 
					{$-\infty$,$-10$,$-4$,$2$,$+\infty$} 
				 \tkzTabLine{,-,0,+,d,+,0,-}
				 \tkzTabVar{+/$+\infty$,-/$24$,+D-/$+\infty$/$-\infty$,+/$0$,-/$-\infty$}
			\end{tikzpicture}
		\end{center}
	\choice 
			 {$y=\dfrac{x^2-3x+4}{-x-4}$}
			{\True $y=\dfrac{x^2-4x+4}{-x-4}$} 
			{$y=\dfrac{x^2-5x+4}{x+4}$} 
			{$y=\dfrac{x^2-4x+4}{x+4}$} 
	\loigiai{ 
			Dựa vào  bảng biến thiên ta thấy:\\
			Tiệm cận đứng của đồ thị hàm số  là đường thẳng $x=4$ và $\lim\limits_{x \rightarrow+\infty}y=-\infty; \lim\limits_{x \rightarrow-\infty}y=+\infty$.\\
			Điểm cực đại $A(2;0)$ và điểm  cực tiểu $B(-10; 24)$. Do đó hàm số cần tìm là $y=\dfrac{x^2-4x+4}{-x-4}$.
		} 
	\end{ex}

% nguồn: Phan1_Lop_12_Dai_so_Chuong_1
\begin{ex}%[2D1H5-1]%[ma:3cf3e11f]
		Hình vẽ sau đây là đồ thị của một trong bốn hàm số cho ở các đáp án A,  B, C, D. Hỏi đó là hàm số nào?
		\begin{center}
			\begin{tikzpicture}[line join=round,line cap=round,>=stealth,scale=0.75,font=\footnotesize]
				 \tikzset{ethicknode/.style={outer sep=-0.5mm}}
				\def \xmin{-2.5} \def \xmax{2.5} 
				\def  \ymin{-3} \def \ymax{3}
				\def  \fx{(\x)^3-3*\x}
				 \draw[->] (\xmin,0)--(\xmax,0) node[below left] {\normalsize$x$};
				\draw[->] (0,\ymin)--(0,\ymax) node[below left] {\normalsize$y$}; 
				\draw  (0,0) node[below left] {$O$};
				\fill  (-1,0) circle (1.25pt) node[below] {$-1$};
				\fill (1,0)  circle (1.25pt) node[above] {$1$};
				 \draw[dashed,thin](1,0)--(1,-2)--(0,-2)
					(0,2)--(-1,2)--(-1,0); 
				\begin{scope}
					\clip (\xmin+0.01,\ymin+0.01)  rectangle (\xmax-0.01,\ymax-0.01);
					\draw[samples=150,domain=\xmin+0.01:\xmax-0.01,smooth,variable=\x]  plot (\x,{\fx});
				\end{scope}
			\end{tikzpicture}
		\end{center}
	 \choice
			 {$y=-x^3+3x$}
			 {$y=x^3+2x+1$}
			{\True  $y=x^3-3x$}
			 {$y=x^3+3x^2$}
	\loigiai{ 
			Từ đồ thị của hàm số $y=ax^3+bx^2+cx+d$, ta thấy hệ  số $a>0$.\\
			Mặt khác ta thấy đồ thị đi qua  điểm $O(0;0)$ nên $d=0$.\\
			Lại có đồ thị nhận $x=1$ là  điểm cực trị nên đáp án là $y=x^3-3x$.
		} 
	\end{ex}

% nguồn: Phan1_Lop_12_Dai_so_Chuong_1
\begin{ex}%[2D1H5-1]%[ma:0fd68a75]
	\immini{ 
		Đồ thị sau đây là của hàm  số nào?
	\choice 
			 {\True $y=x^3-3x+1$}
			{$y=x^3-3x-1$} 
			{$y=-x^3-3x-1$} 
			 {$y=-x^3+3x+1$}
			} 
			 {
			\begin{tikzpicture}[line cap=butt,line join=miter,>=stealth,scale=0.7]
				\tikzset{declare function={xmin=-2.5; xmax=2.5;  ymin=-2; ymax=4;
					f (\x)=1*(\x)^3-3*(\x)+1; 
					}, 
					 smooth,samples=450
				} 
				 \draw[->] (xmin,0)--(xmax,0) node[shift={(-100:7pt)},font=\normalsize]{$ x $};
				\draw[->] (0,ymin)--(0,ymax) node[shift={(190:7pt)},font=\normalsize]{$  y $};
				\fill (0,0) node[shift={(225:7pt)},font=\normalsize]{$ O  $};
				\draw[dashed]  (-1,0) node[below]{$-1$}--(-1,3)--(0,3) node[right]{$3$}
					(1,0)  node[above]{$1$}--(1,-1)--(0,-1) node[left]{$-1$}
					 (0,1) node[right]{$1$};
				\clip (xmin,ymin) rectangle (xmax,ymax); 
				\draw plot[domain=xmin:xmax] (\x,{f (\x)}); 
			\end{tikzpicture}
			 }
	\loigiai{ 
			Đồ thị hàm số trên là đồ thị hàm số bậc ba  $y=a x^3+b x^2+c x+d~(a \neq 0)$.\\
			Nhìn vào nhánh phải của đồ thị ta thấy đồ thị có hướng  đi lên suy ra $a > 0$.\\
			Đồ thị hàm  số cắt trục tung tại điểm có tung độ $y=1$. \\
			Vậy hàm số thỏa  đề là $y=x^3-3x+1$.
		} 
	\end{ex}

% nguồn: Phan1_Lop_12_Dai_so_Chuong_1
\begin{ex}%[2D1H5-1]%[ma:b268bc75]
		Bảng biến thiên  ở hình bên dưới là một trong bốn hàm số nào sau đây?
		\begin{center}
			\begin{tikzpicture}[scale=1,font=\footnotesize,line join=round,line cap=round,>=stealth]
				 \tkzTabInit[nocadre=false,lgt=1.2,espcl=2.5,deltacl=0.6]
					 {$x$/0.6,$y'$/0.6,$y$/2}
					 {$-\infty$,$0$,$2$,$+\infty$}
				\tkzTabLine{,-,0,+,0,-} 
				\path 
					($(N13)!0.8!(N12)$) node  (A1){$+\infty$}
					($(N23)!0.1!(N22)$)  node (A2){$1$}
					 ($(N33)!0.8!(N32)$) node (A3){$5$}
					 ($(N43)!0.1!(N42)$) node (A4){$-\infty$};
				\foreach \x/\y in  {A1/A2,A2/A3,A3/A4}{
					\draw[-stealth]  (\x)--(\y);
				 }
			\end{tikzpicture}
		\end{center}
	\choice 
			 {$y=-x^3-3x^2$}
			{$y=x^3-3x^2-1$} 
			{$y=x^3+2x^2+1$} 
			 {\True $y=-x^3+3x+1$}
	 \loigiai{
			Ta thấy đây là hàm số bậc ba  và $\lim\limits_{x \to-\infty}y=-\infty$ nên $a < 0$.\\
			Ta có $f(0)=1$ nên hàm số cần tìm là $y=-x^3+3x+1$. 
		} 
	\end{ex}

% nguồn: Phan1_Lop_12_Dai_so_Chuong_1
\begin{ex}%[2D1H5-1]%[ma:d5a5e4dd]
		Hàm số nào dưới đây có bảng biến  thiên như hình vẽ?
		\begin{center}
			\begin{tikzpicture}[line join=round,font=\footnotesize,>=stealth,scale=1]
				\tkzTabInit[nocadre=false,lgt=1.2,espcl=2.5,deltacl=0.6] 
					 {$x$ /0.6,$y'$ /0.6,$y$ /2}
					 {$-\infty$,$2$,$+\infty$}
				 \tkzTabLine{,+,d,+,}
				 \tkzTabVar{-/$-\infty$,+D-/$+\infty$/$-\infty$,+/$+\infty$}
			\end{tikzpicture}
		\end{center}
	 \choice
			 {$y=\dfrac{x^2-3}{x-2}$}
			{\True  $y=\dfrac{x^2-4x+2}{x-2}$}
			{$y=\dfrac{x^2-x}{x-2}$} 
			{$y=\dfrac{x^2-4x+5}{x-2}$} 
	 \loigiai{
			Dựa  vào bảng biến thiên, ta thấy:\\
			Tiệm  cận đứng là đường thẳng $x=2$ và có đạo hàm $y'>0$ với mọi $x\neq2$.\\
			 Do đó hàm số cần tìm là $y=\dfrac{x^2-4x+2}{x-2}$.
		} 
	\end{ex}

% nguồn: Phan1_Lop_12_Dai_so_Chuong_1
\begin{ex}%[2D1H5-1]%[ma:55258653]
		 Hàm số nào dưới đây có đồ thị như hình vẽ?
		\begin{center}
			\begin{tikzpicture}[line cap=round,line join=round,font=\footnotesize,>=stealth,scale=1,x=0.5cm,y=0.5cm]
				 \tikzset{label style/.style={font=\footnotesize}}
				\draw[->] (-5,0)--(7,0) node[below] {$x$}; 
				 \draw[->] (0,-5)--(0,7) node[left] {$y$};
				\draw (0,0) node [above  left] {\scriptsize$O$};
				\foreach \x/\nx in {-5/-5,-4/-4,-3/-3,-2/-2,-1/,1/1,2/2,3/3,4/4,5/5,6/6} 
				\draw[thin] (\x,1pt)--(\x,-1pt) node  [below] {\scriptsize$\nx$};
				 \foreach \y/\ny in {-4/right,-3/right,-2/right,-1/right,1/right,2/right,3/right,4/right,5/right,6/right}
				\draw[thin] (1pt,\y)--(-1pt,\y) node [\ny]  {\scriptsize $\y$};
				\draw[dashed,thin] (-1,-5)--(-1,6)  node[above,rotate=90,pos=0.1]{$x=-1$};
				\begin{scope}
					\clip (-6,-5) rectangle (6,5); 
					\draw[samples=200,domain=-5:-1.01,smooth,variable=\x,magenta]  plot (\x,{(-1*((\x)^2))/(1*(\x)+1)});
					\draw[samples=200,domain=-0.9:5.5,smooth,variable=\x,magenta]  plot (\x,{(-1*((\x)^2))/(1*(\x)+1)});
					\draw[dashed,thin]  (-5,6)--(6,-5) node[above,rotate=-45,pos=0.8]{$y=-x+1$};
				\end{scope}
			\end{tikzpicture}
		\end{center}
	\choice 
			{$y=\dfrac{x^2-x}{x+1}$} 
			 {$y=\dfrac{x^2-3x}{x+1}$}
			{$y=\dfrac{x^2+x+2}{x+1}$} 
			{\True  $y=\dfrac{-x^2}{x+1}$}
	 \loigiai{
			Dựa vào đồ thị ta thấy: 
			\begin{itemize}
				\item Tiệm cận đứng là đường thẳng $x=-1$, tiệm cận xiên  là đường thẳng $y=-x+1$.
				\item Giới hạn $\lim\limits_{x \rightarrow+\infty}y=-\infty; \lim\limits_{x  \rightarrow-\infty}y=+\infty$.
			\end{itemize}
			Do đó hàm số cần tìm là  $y=-\dfrac{x^2}{x+1}$.
		 }
	\end{ex}

% nguồn: Phan1_Lop_12_Dai_so_Chuong_1
\begin{ex}%[2D1H5-1]%[ma:e3697090]
		Tìm $m$ để  đồ thị của hàm số $y=x^3+3m x^2-2x-4$ cắt trục hoành tại $3$ điểm phân biệt có hoành độ là $x_1$, $x_2$, $x_3$ thỏa mãn $x_1+x_2=2x_3$.
	\choice 
			 {$m=0$}
			 {\True $m=1$}
			{$m=2$} 
			 {$m=3$}
	 \loigiai{
			Xét phương trình hoành độ giao điểm của $y=x^3+3m x^2-2x-4$ và trục hoành  ta có
			\[x^3+3 m x^2-2 x-4=0. \qquad (1)\] 
			Áp dụng định  lý Vi-et ta có $x_1+x_2+x_3=3x_3=-3m \Rightarrow x_3=-m$.\\
			 Thay $x_3=-m$ vào (1) ta được $2m^3+2m-4=0\Leftrightarrow m=1$.\\
			Thử lại: Thay $m=1$ vào (1) ta được  $x^3+3x^2-2x-4=0\Leftrightarrow\hoac{&x=-1\\& x=-1\pm \sqrt{5}}$, thỏa mãn điều kiện.\\
			 Vậy $m=1$ thỏa mãn điều kiện bài toán.
		 }
	\end{ex}

% nguồn: Phan1_Lop_12_Dai_so_Chuong_1
\begin{ex}%[2D1H5-1]%[ma:8dcfb462]
		Hàm số nào sau  đây có bảng biến thiên như hình dưới đây?
		\begin{center}
			\begin{tikzpicture}
				 \tkzTabInit[nocadre=false,lgt=1.2,espcl=2.5,deltacl=0.6]
					 {$x$/0.6,$y'$/0.6,$y$/2}
					{$-\infty$,$2$,$+\infty$} 
				\tkzTabLine{,-,d,-} 
				 \tkzTabVar{+/$2$,-D+/$-\infty$ / $+\infty$,-/$2$}
			\end{tikzpicture}
		\end{center}
	 \choice
			{$y  = \dfrac{2x - 7}{x - 2}$}
			{$y = \dfrac{2x + 1}{x +  2}$}
			 {\True $y = \dfrac{2x + 1}{x - 2}$}
			 {$y = \dfrac{1 - 2x}{x - 2}$}
	\loigiai{ 
			Ta có $y  = \dfrac{2x + 1}{x - 2}$.\\
			Ta có $\lim\limits_{x \to +\infty} \dfrac{2x + 1}{x - 2} = 2$,$\lim\limits_{x \to -\infty} \dfrac{2x + 1}{x - 2} = 2$, $\lim\limits_{x \to 2^+} y = +  \infty$, $\lim\limits_{x \to 2^-} y = - \infty$.\\
			Lại có $y' = \dfrac{-5}{(x - 2)^2} <  0\ \forall x\ne 2$ nên hàm số nghịch biến trên các khoảng $(-\infty;2)$ và $(2;+\infty)$.
		 }
	\end{ex}

% nguồn: Phan1_Lop_12_Dai_so_Chuong_1
\begin{ex}%[2D1H5-1]%[ma:a85f07de]
	\immini{ 
		Một trong bốn  hàm số sau đây có đồ thị như bên. Hàm số đó là hàm số nào?
	\choice 
			 {\True $y=\dfrac{-x^3+12 x}{8}-1$}
			{$y=\dfrac{-x^3+12 x}{-8}-1$} 
			{$y=\dfrac{-x^3-12 x}{8}-1$} 
			{$y=\dfrac{x^3+12  x}{8}-1$}}
			{\resizebox{!}{5cm}{\begin{tikzpicture}[scale=0.7,>=stealth]
		 \draw[cyan!50,thin] (-4.5,-4) grid (4.5,3.5);
		\draw[->] (-4,0)--(0,0) 
			node[below  right]{$O$}--(5,0) node[below]{$x$};
		\draw[->]  (0,-4) --(0,4) node[right]{$y$};
		\foreach \x in  {-4,-3,-2,-1,1,2,3,4}{
			\draw (\x,-.1)--(\x,.1)  node[above]{$\x$};
		} 
		\foreach  \y in {-3,-2,-1,1,2,3}{
			\draw  (-.1,\y)--(.1,\y) node[right]{$\y$};
		} 
		\draw[samples=100,domain=-4.1:4.2]  plot (\x,{(-(\x)^3+12*(\x))/8-1});
	\end{tikzpicture}
			} 
			 }
	\loigiai{ 
			Quan sát đồ thị ta thấy đây là đồ  thị hàm số bậc $3$ với hệ số $a<0$.\\
			Đồ thị đi  qua điểm $(2;1)$, ta thay toạ độ vào thì thấy chỉ có $y=\dfrac{-x^3+12x}{8}-1$ thoả mãn.
		 }
	\end{ex}

% nguồn: Phan1_Lop_12_Dai_so_Chuong_1
\begin{ex}%[2D1H5-1]%[ma:0d70be99]
	 \immini
		 {Đường cong ở hình vẽ bên là đồ thị của hàm số nào dưới đây?
	 \choice
			 {$y=\dfrac{x-1}{x+1}$}
			{\True $y=\dfrac{x+1}{x-1}$} 
			{$y=\dfrac{x^2-x+1}{x-1}$} 
			 {$y=x^3-3 x-1$}}
			 {
			\begin{tikzpicture}[scale=.75,font=\footnotesize,line join=round,line cap=round,>=stealth]
				\def\a{1}  \def\b{1} \def\c{1} \def\d{-1}
				\def\tcn{\a/\c}\def\tcd{-\d/\c} 
				\draw[->] (-3,0)  -- (5,0) node[below] {$x$};
				\draw[->] (0,-3) -- (0,5)  node[left] {$y$};
				\clip (-3,-3) rectangle  (5,5);
				\fill (0,0) circle (1pt) node[below right]{$O$}; 
				 \draw[dashed] (\tcd,-5)--(\tcd,5)
					 (-5,\tcn)--(5,\tcn);
				 \pgfmathsetmacro{\can}{-(\d)/(\c)}
				\fill (1,0) circle (1pt) node[below right] {$1$}; 
				\draw[samples=150,smooth,domain=-3:{\can-.1}]  plot (\x,{(\a*\x+(\b))/(\c*\x+(\d))});
				 \draw[samples=150,smooth,domain={\can+.1}:5] plot (\x,{(\a*\x+(\b))/(\c*\x+(\d))});
			\end{tikzpicture}
			} 
	\loigiai{ 
			Đồ thị của hàm số phân thức dạng $y=\dfrac{ax+b}{cx+d}$ nên loại {$y=x^3-3 x-1$} và 
			 {$y=\dfrac{x^2-x+1}{x-1}$}.\\
			Tiệm  cận đứng $x=1$ loại {$y=\dfrac{x-1}{x+1}$}.\\
		Vậy đáp án là  hàm số {$y=\dfrac{x+1}{x-1}$}.}
	\end{ex}

% nguồn: Phan1_Lop_12_Dai_so_Chuong_1
\begin{ex}%[2D1H5-1]%[ma:9e5a5b0f]
		Hàm số nào sau đây mà đồ thị có  dạng như hình vẽ bên dưới?
		\begin{center}
			\begin{tikzpicture}[scale=0.7,>=stealth,font=\footnotesize,line join=round,line cap=round]
				\draw[->] (-3,0)  -- (5,0) node[below] {$x$};
				\draw[->] (0,-3) -- (0,5)  node[left] {$y$};
				\draw[dashed] (-3,1) -- (5,1); 
				\node at (-2.3,1.2) {\footnotesize $y=1$}; 
				\draw[dashed]  (1,-3) -- (1,5);
				\node at (1.5,-2.8) {\footnotesize  $x=1$};
				\draw[domain=-3:0.75,smooth] plot (\x,{1-1/(1-\x)}); 
				 \draw[domain=1.26:5,smooth] plot (\x,{1-1/(1-\x)});
				 \node at (0,0) [below left] {\footnotesize $O$};
			\end{tikzpicture}
		\end{center}
	\choice 
			 {$y = \dfrac{x}{1 - x}$}
			{$y =  \dfrac{x + 1}{1 - x}$}
			{$y = \dfrac{x  + 1}{x - 1}$}
			{\True  $y = \dfrac{x}{x - 1}$}
	 \loigiai{
			Đồ thị trong hình có đường tiệm cận ngang $y = 1$ và  tiệm cận đứng $x = 1$, đồng thời cắt trục hoành tại điểm $M(0; 0)$.\\
			Mặt  khác ta thấy hàm số nghịch biến trên từng khoảng xác định. Do đó chỉ có hàm số $y=\dfrac{x}{x-1}$ thỏa mãn những yếu tố trên.
		} 
	\end{ex}

% nguồn: Phan1_Lop_12_Dai_so_Chuong_1
\begin{ex}%[2D1H5-1]%[ma:f22398b3]
		Hàm số nào sau đây có đồ  thị như hình vẽ?
		\begin{center}
			\begin{tikzpicture}[line cap=round,line join=round,font=\footnotesize,>=stealth,scale=1,x=0.5cm,y=0.5cm]
				\tikzset{label style/.style={font=\footnotesize}} 
				\draw[->] (-5,0)--(6,0)  node[below] {$x$};
				\draw[->] (0,-4)--(0,8)  node[left] {$y$};
				\draw (0,0) node [above  left] {\scriptsize$O$};
				\foreach \x/\nx in {-4/-4,-3/-3,-2/-2,-1/,1/1,2/2,3/3,4/4,5/5} 
				\draw[thin] (\x,1pt)--(\x,-1pt) node [below]  {\scriptsize$\nx$};
				\foreach \y/\ny in  {-3/left,-2/left,-1/left,1/right,2/right,3/left,4/left,5/left,6/left}
				\draw[thin] (1pt,\y)--(-1pt,\y) node [\ny] {\scriptsize $\y$}; 
				 \draw[dashed,thin] (1,-4)--(1,8) node[above,rotate=90,pos=0.1]{$x=1$};
				\begin{scope}
					\clip  (-5,-4) rectangle (6,8);
					\draw[samples=200,domain=-4:0.9,smooth,variable=\x,magenta] plot  (\x,{(1*((\x)^2)+(\x)-3)/(1*(\x)+-1)});
					 \draw[samples=200,domain=1.1:6,smooth,variable=\x,magenta] plot (\x,{(1*((\x)^2)+1*(\x)-3)/(1*(\x)+-1)});
					 \draw[dashed,thin] (-4,-2)--(5,7) node[above,rotate=45,pos=0.8]{$y=x+2$};
				\end{scope}
			\end{tikzpicture}
		\end{center}
	\choice 
			{$y=\dfrac{x^2+3}{x-1}$} 
			{\True  $y=\dfrac{x^2+x-3}{x-1}$}
			 {$y=\dfrac{x^2-2x+3}{-x+1}$}
			{$y=\dfrac{x^2+3}{-x+1}$} 
	 \loigiai{
			Dựa vào đồ  thị ta thấy
			\begin{itemize}
				\item Tiệm cận đứng là đường thẳng $x=1$, tiệm cận xiên là đường  thẳng $y=x+2$.
				\item Giới hạn  $\lim\limits_{x \rightarrow+\infty}y=+\infty; \lim\limits_{x \rightarrow-\infty}y=-\infty$.
			\end{itemize}
			Do đó  hàm số cần tìm là $y=\dfrac{x^2+x-3}{x-1}$.
		} 
	\end{ex}

% nguồn: Phan1_Lop_12_Dai_so_Chuong_1
\begin{ex}%[2D1H5-1]%[ma:ad8966b3]
	 \immini{
		Đường cong trong hình bên là đồ thị của hàm  số nào?
	\choice 
			 {$y=\dfrac{1}{2} x^3+\dfrac{3}{2} x^2-2x+1$}
			 {$y=x^3-3x^2+1$}
			{$y=-\dfrac{1}{2} x^3+3x^2+\dfrac{9}{2}  x+1$}
			 {\True $y=\dfrac{1}{2} x^3-3x^2+\dfrac{9}{2} x+1$}
			} 
			 {
			\begin{tikzpicture}[line cap=butt,line join=miter,>=stealth,scale=0.7]
				 \tikzset{declare function={xmin=-1.5; xmax=5; ymin=-1; ymax=4;
					f (\x)=0.5*(\x)^3-3*(\x)^2+4.5*(\x)+1;}, 
					 smooth,samples=250
				} 
				 \draw[->] (xmin,0)--(xmax,0) node[shift={(-100:7pt)},font=\normalsize]{$ x $};
				\draw[->] (0,ymin)--(0,ymax) node[shift={(190:7pt)},font=\normalsize]{$ y $}; 
				 \fill (0,0) node[shift={(-50:7pt)},font=\normalsize]{$ O $};;
				\clip  (xmin,ymin) rectangle (xmax,ymax);
				\draw plot[domain=xmin:xmax] (\x,{f  (\x)});
			\end{tikzpicture}
			 }
	 \loigiai{
			Nhánh phải của đồ thị đi lên  suy ra $a>0$.\\
			Giao điểm của  đồ thị với $Oy$ nằm trên $Ox$ suy ra $d>0$.\\
			Hàm số có hai điểm cực trị  $x_1$; $x_2$ thỏa mãn $\heva{&x_1+x_2=-\dfrac{2b}{3a} > 0\\&x_1x_2=\dfrac{c}{3a} > 0}
			\Rightarrow\heva{&\dfrac{b}{a} < 0\\&\dfrac{c}{a} > 0}  \Rightarrow b < 0$; $c> 0$.\\
			Vậy đồ thị là của hàm số $y=\dfrac{1}{2} x^3-3x^2+\dfrac{9}{2} x+1$. 
		 }
	\end{ex}

% nguồn: Phan1_Lop_12_Dai_so_Chuong_1
\begin{ex}%[2D1V5-1]%[ma:01523f21]
	\immini{ 
	Cho hàm số $y=a x^3+b x^2+c  x+d$ có đồ thị như hình vẽ bên. Mệnh đề nào dưới đây đúng?
	 \choice
		{\True $a < 0$, $b > 0$, $c  > 0$, $d < 0$}
		{$a < 0$, $b <  0$, $c > 0$, $d < 0$}
		{$a > 0$, $b < 0$, $c < 0$, $d  > 0$}
		{$a < 0$, $b  > 0$, $c < 0$, $d > 0$}
		} 
		{ 
		\begin{tikzpicture}[line cap=butt,line join=miter,>=stealth,scale=0.7]
			\tikzset{declare function={xmin=-2; xmax=2.5; ymin=-3; ymax=2.5; 
				f (\x)=-1.29*(\x)^3+1.23*(\x)^2+3.14*(\x)-1.39; 
				}, 
				smooth,samples=450 
			 }
			\draw[->] (xmin,0)--(xmax,0) node[shift={(-100:7pt)},font=\normalsize]{$ x  $};
			 \draw[->] (0,ymin)--(0,ymax) node[shift={(190:7pt)},font=\normalsize]{$ y $};
			\fill (0,0)  node[shift={(225:7pt)},font=\normalsize]{$ O $};
			 \clip (xmin,ymin) rectangle (xmax,ymax);
			\draw plot[domain=xmin:xmax]  (\x,{f (\x)});
		\end{tikzpicture}
		 }
	 \loigiai{
		Ta thấy nhánh đồ thị ngoài cùng bên phải hướng xuống suy ra $a  < 0$.\\
		Đồ thị cắt  trục tung tại điểm có tung độ âm suy ra $d < 0$.\\
		Gọi $x_1$, $x_2$ là $2$ điểm cực trị của hàm số.\\ 
		Ta có $x_1+x_2=\dfrac{-b}{3a}  > 0\Rightarrow b > 0$ và $x_1 \cdot x_2=\dfrac{c}{3a} < 0\Rightarrow c > 0$.\\
		Vậy $a < 0$, $b > 0$, $c  > 0$, $d < 0$.
	 }
\end{ex}

% nguồn: Phan1_Lop_12_Dai_so_Chuong_1
\begin{ex}%[2D1V5-1]%[ma:c38237bc]
	 \immini{
	Cho hàm số $y=a x^3+b x^2+c  x+d$ có đồ thị như hình vẽ. Chọn khẳng định đúng?
	 \choice
		{$a b > 0$, $b c <  0$, $c d < 0$}
		{\True  $a b > 0$, $b c < 0$, $c d > 0$}
		{$a b > 0$, $b c > 0$,  $c d > 0$}
		{$a b <  0$, $b c < 0$, $c d > 0$}
		} 
		 {
		\begin{tikzpicture}[line cap=butt,line join=miter,>=stealth,scale=0.7]
			 \tikzset{declare function={xmin=-3.5; xmax=2.5; ymin=-2.5; ymax=2.5;
				 f (\x)=-0.34*(\x)^3-0.5*(\x)^2+1.84*(\x)+1;
				}, 
				 smooth,samples=450
			 }
			\draw[->]  (xmin,0)--(xmax,0) node[shift={(-100:7pt)},font=\normalsize]{$ x $};
			\draw[->] (0,ymin)--(0,ymax) node[shift={(190:7pt)},font=\normalsize]{$ y $}; 
			\fill (0,0) node[shift={(225:7pt)},font=\normalsize]{$ O  $};
			\draw[dashed] (-2,0) node[above]{$-2$}--(-2,-2) 
				 (1,0) node[below]{$1$}-- (1,2);
			 \clip (xmin,ymin) rectangle (xmax,ymax);
			\draw plot[domain=xmin:xmax]  (\x,{f (\x)});
		\end{tikzpicture}
		} 
	 \loigiai{
		Dựa vào dáng điệu đồ thị,  ta suy ra $a < 0$.\\
		Đồ thị  cắt trục tung lại điểm có tung độ dương suy ra $d > 0$.\\
		Hàm  số có các điểm cực trị $x=1$ và $x=-2$ nên phương trình $y'=3a x^2+2b x+c=0$ có hai nghiệm là $x=1$ và $x=-2$. \\
		Ta có $\dfrac{-2b}{3a}=-1$ và $\dfrac{c}{3a}=-2$.\\ 
		Do đó $b < 0$ và $c  > 0$.\\
		Như vậy $a b > 0$, $b c < 0$ và $c d  > 0$.
	} 
\end{ex}

% nguồn: Phan1_Lop_12_Dai_so_Chuong_1
\begin{ex}%[2D1V5-1]%[ma:c2516224]
	 \immini{
	 Cho hàm số $y=a x^3+b x^2+c x+d$ có đồ thị như hình vẽ.
	Số lớn nhất trong các số $a$,  $b$, $c$, $d$ là
	 \choice
		 {\True $c$}
		 {$a$}
		 {$b$}
		 {$d$}
		 }
		{\begin{tikzpicture}[line cap=butt,line join=miter,>=stealth,scale=0.7]
		 \tikzset{declare function={xmin=-1; xmax=4.5; ymin=-2.5; ymax=2.5;
			f (\x)=1*(\x)^3-6*(\x)^2+9*(\x)-2; 
			 },
			 smooth,samples=450
		 }
		\draw[->] (xmin,0)--(xmax,0)  node[shift={(-100:7pt)},font=\normalsize]{$ x $};
		\draw[->] (0,ymin)--(0,ymax) node[shift={(190:7pt)},font=\normalsize]{$ y  $};
		\fill (0,0) node[shift={(225:7pt)},font=\normalsize]{$ O $}; 
		 \draw[dashed]
			(1,0) node[below]{$1$}--(1,2)--(0,2) node[left]{$2$} 
			(3,0)  node[above]{$3$}--(3,-2)--(0,-2) node[left]{$-2$}
			(2,0) node[above right]{$2$}; 
		 \clip (xmin,ymin) rectangle (xmax,ymax);
		\draw plot[domain=xmin:xmax]  (\x,{f (\x)});
	\end{tikzpicture}
		 }
	\loigiai{ 
		Ta có $y' = 3ax^2 + 2bx +  c$.\\
		Đồ thị hàm số có hai điểm cực  trị là $M(1; 2)$ và $N(3;-2)$ nên
		\allowdisplaybreaks 
		\begin{eqnarray*}
			\heva{&a \cdot 1^3 + b \cdot 1^2 + c \cdot 1 + d = 2 \\& 
			a \cdot 3^3 + b \cdot 3^2 +  c \cdot 3 + d = -2 \\&
			3a \cdot 1^2 + 2b \cdot 1 + c  = 0 \\&
			3a \cdot 3^2 + 2b \cdot  3 + c = 0}
			 &\Leftrightarrow& \heva{&a+b+c+d=2 \\&27a+9b+3c+d=-2\\&
			3a+2b+c=0\\&27a+6b+c=0}\\\\ 
			&\Leftrightarrow& \heva{&a=1\\ &b=-6\\ &c=9\\&d=-2.} 
		\end{eqnarray*}
		 Vậy $c = 9$ là số lớn nhất.
	} 
\end{ex}

% nguồn: Phan1_Lop_12_Dai_so_Chuong_1
\begin{ex}%[2D1V5-1]%[ma:77f02304]
	 \immini{
	Cho  hàm số $y=a x^3+b x^2+c x+d$ có đồ thị như hình vẽ.
	Mệnh  đề nào sau đây đúng?
	\choice 
		{$a < 0$, $b > 0$, $c >  0$, $d > 0$}
		{$a > 0$,  $b > 0$, $c < 0$, $d > 0$}
		{$a <  0$, $b < 0$, $c < 0$, $d > 0$}
		{\True $a < 0$, $b > 0$, $c < 0$,  $d > 0$}
		 }
		 {
		\begin{tikzpicture}[line cap=butt,line join=miter,>=stealth,scale=0.7]
			\tikzset{declare function={xmin=-1; xmax=4.5; ymin=-1.5; ymax=3; 
				f (\x)=-0.8*(\x)^3+4.64*(\x)^2-6.42*(\x)+1.55; 
				 },
				 smooth,samples=450
			 }
			\draw[->] (xmin,0)--(xmax,0) node[shift={(-100:7pt)},font=\normalsize]{$ x $}; 
			\draw[->] (0,ymin)--(0,ymax)  node[shift={(190:7pt)},font=\normalsize]{$ y $};
			 \fill (0,0) node[shift={(225:7pt)},font=\normalsize]{$ O $};
			 \clip (xmin,ymin) rectangle (xmax,ymax);
			\draw plot[domain=xmin:xmax] (\x,{f  (\x)});
		\end{tikzpicture}
		} 
	 \loigiai{
		Do nhánh bên phải của đồ thị  đi xuống nên $a < 0$.
		Đồ thị cắt trục tung ở phần dương nên $d  > 0$.
		Đồ thị có 2 cực trị tại hai giá trị $x$ dương nên phương trình $y'=0$ có 2 nghiệm dương phân biệt $\Rightarrow 3a x^2+2b x+c=0$ có $2$ nghiệm dương  phân biệt
		 $$
		\Rightarrow\heva{&\Delta' > 0\\&S > 0\\&P > 0} \Leftrightarrow \heva{&b ^2-3 a c > 0 \\&-\dfrac {2 b}{3 a} > 0\\&\dfrac {c}{3 a} > 0} \Rightarrow\heva{&b > 0 \\& c < 0.} 
		$$ 
	 }
\end{ex}

% nguồn: Phan1_Lop_12_Dai_so_Chuong_1
\begin{ex}%[2D1V5-1]%[ma:a714652e]
	\immini{ 
	Cho hàm số  $y=a x^3+b x^2+c x+d$ có đồ thị như hình bên. Mệnh đề nào sau đây là đúng?
	 \choice
		{$a < 0$, $b  > 0$, $c > 0$, $d > 0$}
		{$a < 0$, $b < 0$, $c < 0$, $d > 0$} 
		{$a > 0$,  $b > 0$, $c < 0$, $d > 0$}
		{\True $a < 0$, $b > 0$, $c < 0$, $d  > 0$}
		 }
		{ 
		\begin{tikzpicture}[line cap=butt,line join=miter,>=stealth,scale=0.7]
			\tikzset{declare  function={xmin=-1; xmax=5; ymin=-2; ymax=2.5;
				 f (\x)=-0.39*(\x)^3+2.55*(\x)^2-4.13*(\x)+0.97;
				}, 
				 smooth,samples=450
			 }
			\draw[->] (xmin,0)--(xmax,0) node[shift={(-100:7pt)},font=\normalsize]{$  x $};
			 \draw[->] (0,ymin)--(0,ymax) node[shift={(190:7pt)},font=\normalsize]{$ y $};
			\fill (0,0) node[shift={(225:7pt)},font=\normalsize]{$  O $};;
			\clip (xmin,ymin)  rectangle (xmax,ymax);
			\draw plot[domain=xmin:xmax] (\x,{f (\x)}); 
		\end{tikzpicture}
		} 
	 \loigiai{
		Dựa vào đồ  thị suy ra $a < 0$.\\
		Giao điểm của đồ thị với trục  $Oy$ suy ra $d > 0$.\\
		Dựa vào cực trị ta có $x_1 \cdot x_2=\dfrac{c}{3a}  > 0\Rightarrow c < 0\Rightarrow x_1+x_2=\dfrac{-2b}{3a} > 0\Rightarrow b > 0$.
	} 
\end{ex}

% nguồn: Phan1_Lop_12_Dai_so_Chuong_1
\begin{ex}%[2D1V5-1]%[ma:1f8a8962]
	 \immini{
	Cho hàm số $y=f(x)=a x^3+b x^2+c x+d$ có đồ thị như hình vẽ bên. Mệnh đề nào sau đây đúng? 
	\choice 
		 {$a > 0$, $b > 0$, $c > 0$, $d > 0$}
		 {\True $a > 0$, $b < 0$, $c > 0$, $d > 0$}
		{$a  < 0$, $b < 0$, $c > 0$, $d < 0$}
		{$a > 0$, $b > 0$, $c < 0$, $d > 0$} 
		 }
		 {
		\begin{tikzpicture}[line cap=butt,line join=miter,>=stealth,scale=0.7]
			\tikzset{declare function={xmin=-1.5;  xmax=5; ymin=-1; ymax=4;
				f  (\x)=0.5*(\x)^3-3*(\x)^2+4.5*(\x)+1;
				}, 
				smooth,samples=450 
			 }
			\draw[->]  (xmin,0)--(xmax,0) node[shift={(-100:7pt)},font=\normalsize]{$ x $};
			\draw[->] (0,ymin)--(0,ymax) node[shift={(190:7pt)},font=\normalsize]{$ y $}; 
			\fill (0,0) node[shift={(-50:7pt)},font=\normalsize]{$ O  $};;
			\clip (xmin,ymin)  rectangle (xmax,ymax);
			\draw plot[domain=xmin:xmax] (\x,{f (\x)}); 
		\end{tikzpicture}
		} 
	 \loigiai{
		Dựa vào dạng đồ thị hàm bậc ba ta có $a  > 0$.\\
		Đồ thị cắt trục tung tại  điểm $M(0; 1)$ suy ra $d=1> 0$.\\
		Hàm số  có hai điểm cực trị dương suy ra $\heva{&-\dfrac{b}{a} > 0\\
		&\dfrac{c}{a} >  0} \Leftrightarrow\heva{&b < 0\\&c > 0.}$\\
		Vậy $a > 0$, $b < 0$, $c > 0$, $d > 0$. 
	} 
\end{ex}

% nguồn: Phan1_Lop_12_Dai_so_Chuong_1
\begin{ex}%[2D1V5-1]%[ma:6d070242]
	 \immini{
	Cho hàm số $y=a x^3+b x^2+c x+d$ có đồ thị là đường cong trong  hình vẽ. Trong các hệ số $a$, $b$, $c$, $d$ có bao nhiêu số âm?
	 \choice
		 {$3$}
		{$0$} 
		{$1$} 
		{\True  $2$}
		 }
		{ 
		\begin{tikzpicture}[line cap=butt,line join=miter,>=stealth,scale=0.7]
			\tikzset{declare function={xmin=-1;  xmax=4.7; ymin=-1; ymax=4;
				 f (\x)=-0.36*(\x)^3+1.94*(\x)^2-2.18*(\x)+1.84;
				 },
				smooth,samples=450 
			 }
			\draw[->] (xmin,0)--(xmax,0) node[shift={(-100:7pt)},font=\normalsize]{$ x $}; 
			\draw[->] (0,ymin)--(0,ymax)  node[shift={(190:7pt)},font=\normalsize]{$ y $};
			\fill (0,0) node[shift={(225:7pt)},font=\normalsize]{$ O $};; 
			\clip (xmin,ymin)  rectangle (xmax,ymax);
			\draw plot[domain=xmin:xmax] (\x,{f  (\x)});
		\end{tikzpicture}
		} 
	 \loigiai{
		Dựa vào hình dạng đồ thị: đồ thị hàm bậc ba có hệ số $a < 0$, đồ thị cắt trục tung tại điểm có  tung độ dương nên $d > 0$.\\
		Ta có $y'=3a  x^2+2b x+c$. \\
		Đồ thị  có hai điểm cực trị cùng nằm bên phải trục tung nên $y'=0$ có hai nghiệm dương phân biệt $x_1$, $x_2$. \\
		Ta có $\heva{&x_1+x_2=-\dfrac{2b}{3a} >  0\\&x_1\cdot x_2=\dfrac{c}{3a} > 0} \Rightarrow\heva{&b > 0\\&c < 0.}$\\
		Vậy  $a < 0$, $c < 0$.
	 }
\end{ex}

% nguồn: Phan1_Lop_12_Dai_so_Chuong_1
\begin{ex}%[2D1V5-1]%[ma:abbc0197]
	 \immini{
	Cho hàm số $y=a x^3+b x^2+c x+d$  có đồ thị như hình bên. Khẳng định nào sau đây đúng?
	 \choice
		{\True $a > 0$, $b < 0$, $c < 0$, $d >  0$}
		{$a  > 0$, $b > 0$, $c < 0$, $d > 0$}
		{$a >  0$, $b < 0$, $c > 0$, $d > 0$}
		{$a <  0$, $b < 0$, $c < 0$, $d > 0$}
		 }
		{ 
		\begin{tikzpicture}[line cap=butt,line join=miter,>=stealth,scale=0.7]
			 \tikzset{declare function={xmin=-2.5; xmax=2.5; ymin=-2; ymax=4;
				f (\x)=1*(\x)^3-3*(\x)+1; 
				}, 
				 smooth,samples=450
			 }
			\draw[->] (xmin,0)--(xmax,0) node[shift={(-100:7pt)},font=\normalsize]{$ x $}; 
			\draw[->] (0,ymin)--(0,ymax) node[shift={(190:7pt)},font=\normalsize]{$  y $};
			 \fill (0,0) node[shift={(225:7pt)},font=\normalsize]{$ O $};
			\draw[dashed] (-1,0)  node[below]{$-1$}--(-1,3)--(0,3) node[right]{$3$}
				 (1,0) node[above]{$1$}--(1,-1)--(0,-1) node[left]{$-1$}
				 (0,1) node[right]{$1$};
			\clip (xmin,ymin) rectangle (xmax,ymax); 
			\draw plot[domain=xmin:xmax] (\x,{f  (\x)});
		\end{tikzpicture}
		 }
	 \loigiai{
		Quan sát đồ thị hàm số ta thấy $a > 0$.\\ 
		Đồ thị hàm số cắt trục tung tại điểm có tung độ dương suy ra $d >  0$.\\
		Hàm số có  hai điểm cực trị $x_1$; $x_2$ thỏa mãn $\heva{&x_1+x_2=-\dfrac{2b}{3a} > 0\\&x_1x_2=\dfrac{c}{3a} < 0}
		\Rightarrow\heva{&\dfrac{b}{a} < 0\\&\dfrac{c}{a}  < 0} \Rightarrow b < 0$; $c < 0$.
	 }
\end{ex}

% nguồn: Phan1_Lop_12_Dai_so_Chuong_1
\begin{ex}%[2D1V5-1]%[ma:f0f62187]
	\immini{ 
	Cho hàm số $y=a x^3+b x^2+c  x+d$ có đồ thị như hình vẽ.
	 Mệnh đề nào dưới đây đúng?
	 \choice
		{$a < 0$, $b <  0$, $c < 0$, $d > 0$}
		{$a > 0$,  $b < 0$, $c > 0$, $d < 0$}
		{\True $a  > 0$, $b < 0$, $c < 0$, $d > 0$}
		{$a >  0$, $b < 0$, $c > 0$, $d > 0$}
		 }
		 {
		\begin{tikzpicture}[line cap=butt,line join=miter,>=stealth,scale=0.7]
			\tikzset{declare  function={xmin=-2; xmax=4; ymin=-2; ymax=4;
				f  (\x)=0.63*(\x)^3-1.25*(\x)^2-1.75*(\x)+2.52;
				}, 
				 smooth,samples=450
			} 
			\draw[->] (xmin,0)--(xmax,0) node[shift={(-100:7pt)},font=\normalsize]{$ x $}; 
			 \draw[->] (0,ymin)--(0,ymax) node[shift={(190:7pt)},font=\normalsize]{$ y $};
			\fill (0,0) node[shift={(225:7pt)},font=\normalsize]{$ O $}; 
			\draw[dashed] (-.5,0)--(-.5,3)--(0,3) 
				(1.84,0)--(1.84,-1.02)--(0,-1.02); 
			\clip (xmin,ymin)  rectangle (xmax,ymax);
			 \draw plot[domain=xmin:xmax] (\x,{f (\x)});
		\end{tikzpicture}
		} 
	 \loigiai{
	\immini{Ta  có $\lim\limits_{x \to+\infty} y=+\infty \Rightarrow a > 0$.\\
		Đồ  thị hàm số cắt $Oy$ tại điểm có tọa độ $(0; d)$ suy ra $d > 0$.\\
		Ta có  $y'=3a x^2+2b x+c$.\\
		 Gọi $x_1$, $x_2$ là các điểm cực trị của hàm số.\\
		Dựa vào đồ thị ta có $\heva{&x_1+x_2=-\dfrac{2b}{3a} > 0\\& x_1 x_2=\dfrac{c}{3a} < 0}$, mà $a > 0$, suy ra  $b < 0$, $c < 0$.\\
		Vậy  $a > 0$, $b < 0$, $c < 0$, $d > 0$.
		 }
		{\begin{tikzpicture}[line cap=butt,line join=miter,>=stealth,scale=0.8]
		 \tikzset{declare function={xmin=-2; xmax=4; ymin=-2; ymax=4;
			 f (\x)=0.63*(\x)^3-1.25*(\x)^2-1.75*(\x)+2.52;
			}, 
			smooth,samples=450 
		 }
		 \draw[->] (xmin,0)--(xmax,0) node[shift={(-100:7pt)},font=\normalsize]{$ x $};
		\draw[->] (0,ymin)--(0,ymax) node[shift={(190:7pt)},font=\normalsize]{$ y $}; 
		\fill (0,0) node[shift={(225:7pt)},font=\normalsize]{$  O $};
		\draw[dashed] (-.5,0) node[below]{$x_1$}--(-.5,3)--(0,3) 
			(1.84,0) node[above]{$x_2$}--(1.84,-1.02)--(0,-1.02); 
		\clip  (xmin,ymin) rectangle (xmax,ymax);
		\draw plot[domain=xmin:xmax]  (\x,{f (\x)});
	\end{tikzpicture}
		} 
	} 
\end{ex}

% nguồn: Phan1_Lop_12_Dai_so_Chuong_1
\begin{ex}%[2D1V5-1]%[ma:ffa19f67]
	\immini{ 
	Cho hàm số  $y=a x^3+b x^2+c x+d$ có đồ thị như hình vẽ. Mệnh đề nào dưới đây đúng?
	 \choice
		{$a >  0$, $b > 0$, $c > 0$, $d > 0$}
		{$a > 0$, $b > 0$, $c < 0$,  $d > 0$}
		{\True $a > 0$, $b < 0$, $c <  0$, $d > 0$}
		{$a > 0$, $b < 0$, $c  < 0$, $d < 0$}
		} 
		{\begin{tikzpicture}[line cap=butt,line join=miter,>=stealth,scale=0.7]
		\tikzset{declare function={xmin=-3; xmax=2.5; ymin=-1;  ymax=3.5;
			 f (\x)=0.33*(\x)^3+0.25*(\x)^2-1.08*(\x)+1.5;
			 },
			smooth,samples=450 
		} 
		\draw[->]  (xmin,0)--(xmax,0) node[shift={(-100:7pt)},font=\normalsize]{$ x $};
		\draw[->] (0,ymin)--(0,ymax) node[shift={(190:7pt)},font=\normalsize]{$ y $}; 
		 \fill (0,0) node[shift={(-50:7pt)},font=\normalsize]{$ O $};;
		\clip (xmin,ymin)  rectangle (xmax,ymax);
		\draw plot[domain=xmin:xmax] (\x,{f (\x)}); 
	\end{tikzpicture}
		} 
	\loigiai{ 
		Đồ thị hàm số đã cho là đồ thị hàm số $y=a x^3+b x^2+c x+d$ với $a > 0$.\\ 
		Đồ thị hàm số cắt $Oy$ tại điểm  $(0, d)$ với $d > 0$.\\
		Hàm số có hai cực trị tại $x_1 < 0$, $x_2 > 0$. Suy ra $x_1x_2 = \dfrac{c}{3a} <  0 \Rightarrow c < 0$.\\
		Ta có $y' = 3ax^2 +  2bx + c$.\\
		Do $x_1 + x_2 = -\dfrac{2b}{3a} < 0$ nên  $b > 0$.
	} 
\end{ex}

% nguồn: Phan1_Lop_12_Dai_so_Chuong_1
\begin{ex}%[2D1V5-1]%[ma:e0118f08]
	 \immini{
	Cho hàm số $y=a x^3+b x^2+c x+d$ có đồ thị như hình vẽ. Mệnh đề nào sau đây  đúng?
	 \choice
		{$a < 0$, $b  > 0$, $c > 0$, $d > 0$}
		{$a < 0$, $b < 0$,  $c < 0$, $d > 0$}
		{$a > 0$, $b >  0$, $c < 0$, $d > 0$}
		{\True  $a < 0$, $b > 0$, $c < 0$, $d > 0$}
		} 
		{ 
		\begin{tikzpicture}[line cap=butt,line join=miter,>=stealth,scale=0.7]
			\tikzset{declare function={xmin=-1; xmax=4.5; ymin=-1.5; ymax=3; 
				f  (\x)=-0.8*(\x)^3+4.64*(\x)^2-6.42*(\x)+1.55;
				}, 
				 smooth,samples=450
			} 
			\draw[->]  (xmin,0)--(xmax,0) node[shift={(-100:7pt)},font=\normalsize]{$ x $};
			\draw[->] (0,ymin)--(0,ymax) node[shift={(190:7pt)},font=\normalsize]{$  y $};
			\fill  (0,0) node[shift={(225:7pt)},font=\normalsize]{$ O $};
			\clip (xmin,ymin) rectangle (xmax,ymax); 
			\draw plot[domain=xmin:xmax] (\x,{f (\x)}); 
		\end{tikzpicture}
		} 
	\loigiai{ 
		Ta  có $a < 0$ và đồ thị cắt $Oy$ tại điểm có tung độ dương nên $d > 0$.\\
		Mặt khác $y'=3a x^2+2b x+c$;  $y''=6a x+2b$ và từ đồ thị ta thấy đồ thị hàm số có hai điểm cực trị có hoành độ dương và điểm uốn có hoành độ dương.\\
		Khi đó $\heva{&a c > 0\\&-\dfrac{b}{3a} > 0} \Rightarrow\heva{&b > 0\\&c < 0}$ do  $a < 0$.
	} 
\end{ex}

% nguồn: Phan1_Lop_12_Dai_so_Chuong_1
\begin{ex}%[2D1V5-1]%[ma:d8ba8f58]
	 \immini{
	Cho hàm số $y=a x^3+b  x^2+c x+d$ có đồ thị như hình vẽ bên. Mệnh đề nào dưới đây đúng?
	\choice 
		{\True  $a < 0$, $b > 0$, $c > 0$, $d < 0$}
		{$a < 0$, $b < 0$, $c > 0$, $d < 0$} 
		{$a > 0$,  $b < 0$, $c < 0$, $d > 0$}
		{$a < 0$, $b >  0$, $c < 0$, $d < 0$}
		 }
		{\begin{tikzpicture}[line cap=butt,line join=miter,>=stealth,scale=0.7]
		\tikzset{declare function={xmin=-2; xmax=2.5;  ymin=-3; ymax=2.5;
			 f (\x)=-1.29*(\x)^3+1.23*(\x)^2+3.14*(\x)-1.39;
			}, 
			smooth,samples=450 
		 }
		\draw[->] (xmin,0)--(xmax,0) node[shift={(-100:7pt)},font=\normalsize]{$ x  $};
		\draw[->]  (0,ymin)--(0,ymax) node[shift={(190:7pt)},font=\normalsize]{$ y $};
		\fill (0,0) node[shift={(225:7pt)},font=\normalsize]{$ O $}; 
		\clip (xmin,ymin) rectangle (xmax,ymax); 
		\draw  plot[domain=xmin:xmax] (\x,{f (\x)});
	\end{tikzpicture}
		} 
	 \loigiai{
		Ta có $y' = 3ax^2 + 2bx  + c$.\\
		 Nhìn vào đồ thị ta có phần bên phải của đồ thị đi xuống nên $a < 0$.\\
		Giao điểm với trục tung  nằm phía dưới điểm $O$ nên $d < 0$.\\
		Hai điểm cực trị nằm về hai phía đối  với trục tung nên suy ra phương trình $y'=0$ có hai nghiệm trái dấu $\Leftrightarrow a c < 0\Leftrightarrow c > 0$. \\
		Điểm uốn lệch phải  so với trục tung nên $-\dfrac{b}{3a} > 0\Leftrightarrow b > 0$.\\
		Vậy  suy ra $a < 0$, $b > 0$, $c > 0$, $d < 0$.
	} 
\end{ex}

% nguồn: Phan1_Lop_12_Dai_so_Chuong_1
\begin{ex}%[2D1V5-1]%[ma:f5150cbc]
	\immini{ 
	Cho hàm số $y=a x^3+b x^2+c x+d$ $(a \neq  0)$ có đồ thị như hình vẽ.
	Mệnh đề nào  sau đây đúng?
	\choice 
		{\True $a > 0$, $b < 0$, $c < 0$, $d  > 0$}
		{$a > 0$, $b > 0$, $c > 0$, $d >  0$}
		{$a > 0$, $b < 0$, $c > 0$,  $d < 0$}
		{$a > 0$, $b >  0$, $c < 0$, $d > 0$}
		} 
		{ 
		\begin{tikzpicture}[line cap=butt,line join=miter,>=stealth,scale=0.7]
			\tikzset{declare  function={xmin=-2; xmax=4; ymin=-2; ymax=4;
				f  (\x)=0.63*(\x)^3-1.25*(\x)^2-1.75*(\x)+2.52;
				}, 
				 smooth,samples=450
			 }
			\draw[->] (xmin,0)--(xmax,0)  node[shift={(-100:7pt)},font=\normalsize]{$ x $};
			\draw[->] (0,ymin)--(0,ymax) node[shift={(190:7pt)},font=\normalsize]{$ y $}; 
			 \fill (0,0) node[shift={(225:7pt)},font=\normalsize]{$ O $};
			\draw[dashed] (-.5,0)--(-.5,3)--(0,3) 
				(1.84,0)--(1.84,-1.02)--(0,-1.02); 
			\clip (xmin,ymin)  rectangle (xmax,ymax);
			\draw plot[domain=xmin:xmax]  (\x,{f (\x)});
		\end{tikzpicture}
		 }
	 \loigiai{
		 Ta có $y'=3a x^2+2b x+c$.\\
		Theo  hình vẽ\\
		Đồ thị cắt trục tung tại điểm $(0, d)$ nằm phía trên  trục hoành nên $d > 0$.\\
		Hàm số  có hai cực trị trái dấu nên $a c < 0$ mà $a > 0$, do đó $c < 0$.\\
		Điểm uốn của đồ thị có hoành độ  dương nên $\dfrac{x_1+x_2}{2}=-\dfrac{2b}{6a} > 0\Leftrightarrow a b < 0$.\\
		 Do $a > 0$ nên $b < 0$.
	} 
\end{ex}

% nguồn: Phan1_Lop_12_Dai_so_Chuong_1
\begin{ex}%[2D1V5-1]%[ma:9674b96a]
	 \immini{
	Hàm số $y=a x^3+b x^2+c x+d$ có đồ thị như  hình vẽ bên. Mệnh đề nào sau đây là đúng?
	 \choice
		{$a < 0$, $b < 0$, $c < 0$, $d < 0$} 
		{$a > 0$, $b > 0$,  $c > 0$, $d < 0$}
		{$a > 0$, $b > 0$, $c  < 0$, $d > 0$}
		{\True $a > 0$, $b < 0$, $c < 0$, $d >  0$}
		 }
		{ 
		\begin{tikzpicture}[line cap=butt,line join=miter,>=stealth,scale=0.7]
			 \tikzset{declare function={xmin=-1.5; xmax=3.5; ymin=-3; ymax=3;
				f (\x)=0.93*(\x)^3-1.94*(\x)^2-1.88*(\x)+1.95; 
				 },
				 smooth,samples=450
			} 
			\draw[->] (xmin,0)--(xmax,0)  node[shift={(-100:7pt)},font=\normalsize]{$ x $};
			 \draw[->] (0,ymin)--(0,ymax) node[shift={(190:7pt)},font=\normalsize]{$ y $};
			\fill (0,0)  node[shift={(225:7pt)},font=\normalsize]{$ O $};
			\fill (1,-0.94) circle (1pt); \fill (-1,0.96) circle (1pt); 
			 \draw[dashed] (-1,0) node[below,xshift=3pt]{$-1$}--(-1,0.96)
				(1,0) node[above]{$1$}--(1,-0.94); 
			\clip (xmin,ymin)  rectangle (xmax,ymax);
			\draw plot[domain=xmin:xmax] (\x,{f  (\x)});
		\end{tikzpicture}
		 }
	\loigiai{ 
		Nhìn vào nhánh phải của đồ thị ta thấy đồ thị có hướng đi lên suy ra $a  > 0$.\\
		Nhìn vào giao điểm của đồ thị với trục tung ta thấy đồ thị cắt trục tung tại điểm có tung độ dương suy ra $d >  0$.\\
		Ta có  $y'=3a x^2+2b x+c$.\\
		Hàm số đã  cho có hai điểm cực trị $x_1$, $x_2$ với $x_1 \cdot x_2 < 0\Leftrightarrow \dfrac{c}{3a} < 0\Leftrightarrow c < 0$ (vì $a > 0$).\\
		Vì $-1< x_1 < 0$ và $x_2 > 1$ nên $x_1+x_2 > 0\Leftrightarrow \dfrac{-2b}{3a} > 0\Leftrightarrow-2b >  0\Leftrightarrow b < 0$ (vì $a > 0$).\\
		Vậy $a > 0$,  $b < 0$, $c < 0$, $d > 0$.
	 }
\end{ex}


%%%%%%%%%%%%%%%% ĐÚNG/SAI %%%%%%%%%%%%%%%%

% nguồn: Phan1_Lop_12_Dai_so_Chuong_1
\begin{ex}%[2D1N5-1]%[ma:3f5f7b7d]
	\immini{ 
	Cho hàm số $y=\dfrac{(a-1)x+b}{(c-1)x+d}$, $d<0$  có đồ thị như hình vẽ.
	 \choiceTF
		{$a>1, b>0, c<1$} 
		{$a>1, b<0, c>1$} 
		{$a<1, b>0, c<1$} 
		{\True  $a>1, b>0, c>1$}
		 }{
		\begin{tikzpicture}[scale=1,font=\footnotesize,line join=round,line cap=round,>=stealth]
			\def\a{1} \def\b{1} \def\c{2} \def\d{-1} 
			\def\xmin{-2} \def\xmax{3} 
			 \def\ymin{-2} \def\ymax{3}
			\draw[->]  (\xmin,0)--(\xmax,0) node [below]{$x$};
			\draw[->] (0,\ymin)--(0,\ymax)  node [left]{$y$};
			\node at (0,0)  [above left]{$O$};
			\clip (\xmin+0.1,\ymin+0.1) rectangle (\xmax-0.1,\ymax-0.1); 
			\draw[smooth,samples=300,domain=\xmin:(-\d/\c-0.1)]  plot (\x,{(\a*(\x)+\b)/(\c*(\x)+\d)});
			\draw[smooth,samples=300,domain=(-\d/\c+0.1:\xmax)] plot  (\x,{(\a*(\x)+\b)/(\c*(\x)+\d)});
			 \draw[dashed] (-\d/\c,\ymin)--(-\d/\c,\ymax);
			\draw[dashed]  (\xmin,\a/\c)--(\xmax,\a/\c);
		\end{tikzpicture}
		} 
	 \loigiai{
		Theo bài  ra, đường tiệm cận đứng của đồ thị hàm số là đường thẳng là $x=-\dfrac{d}{c-1}$.\\
		Đường tiệm cận ngang của  đồ thị hàm số là đường thẳng $y=\dfrac{a-1}{c-1}$.\\
		Nhìn đồ thị ta  thấy: $x=-\dfrac{d}{c-1}>0$ mà $d<0\Rightarrow c-1>0\Rightarrow c>1$.\\
		$y=\dfrac{a-1}{c-1}>0\Rightarrow a-1>0\Rightarrow  a>1$.\\
		Đồ thị cắt trục  tung tại điểm có tung độ bằng $\dfrac{b}{d}<0\Rightarrow b>0$.
		\begin{itemchoice}
			 \itemch
			 \itemch
			 \itemch
			\itemch 
		\end{itemchoice}
	} 
\end{ex}

% nguồn: Phan1_Lop_12_Dai_so_Chuong_1
\begin{ex}%[2D1N5-1]%[ma:0afefa06]
	\immini{ 
	Hình vẽ sau là  đồ thị của hàm số $y=\dfrac{ax+b}{cx+d}$.
	 \choiceTF
		{\True $ad>0$} 
		{\True  $bd<0$}
		 {$ab>0$}
		{$ad-bc<0$} 
		}{ 
		\begin{tikzpicture}[>=stealth',scale=0.5]
			 \def\xmin{-5.5} \def\xmax{4}
			\def\ymin{-2.5}  \def\ymax{6.5}
			\draw[domain=-5.5:4] plot  (\x,{2});
			\draw[domain=-2.5:6.5,variable=\t]  plot ({-1},{(\t)});
			\draw[domain=-5.5:-1.67] plot (\x,{(2*(\x)-1)/((\x)+1)}); 
			 \draw[domain=-0.33:4] plot (\x,{(2*(\x)-1)/((\x)+1)});
			\begin{scriptsize}
				\draw[->](\xmin,0)--(\xmax,0); \draw (\xmax-0.1,0) node[below]{$x$}; 
				\draw[->](0,\ymin)--(0,\ymax); \draw  (0,\ymax-0.2) node[right]{$y$};
				\draw  node [above left]{$O$};
			\end{scriptsize}
		\end{tikzpicture}
		 }
	 \loigiai{
		Dựa vào  đồ thị hàm số ta thấy
		\begin{itemchoice}
			\itemch Đồ thị hàm số có tiệm cận ngang là đường thẳng $y=\dfrac{a}{c}>0$ và tiệm cận đứng là đường thẳng $x=-\dfrac{d}{c}<0.$ Ta  có $\heva{& \dfrac{a}{c}>0\\&\dfrac{d}{c}>0}\Rightarrow\dfrac{ad}{c^2}>0\Rightarrow ad>0$.
			\itemch Đồ thị hàm số cắt trục $Oy$ tại điểm có tung độ nhỏ hơn $0\Rightarrow\dfrac{b}{d}<0\Rightarrow  bd<0.$
			\itemch  Từ $ad>0$ và $bd<0$ suy ra $adbd<0\Rightarrow ab<0$.
			\itemch Hàm số đã cho là hàm đồng biến trên mỗi  khoảng xác định nên
			$$y'=\dfrac{ad-bc}{(cx+d)^2}>0\Rightarrow  ad-bc>0.$$
		\end{itemchoice}
	} 
\end{ex}

% nguồn: Phan1_Lop_12_Dai_so_Chuong_1
\begin{ex}%[2D1N5-1]%[ma:6a7ab0bb]
	Cho hàm  số $y=\dfrac{x+a}{bx+c}$ với $a, b, c\in \mathbb{Z}$ có đồ thị như hình vẽ dưới đây
	\begin{center}
		\begin{tikzpicture}[scale=0.7,font=\footnotesize,line join=round,line cap=round,>=stealth]
			\def\xmin{-3}; \def\ymin{-3}; \def\xmax{4};  \def\ymax{4};
			\coordinate (O) at (0,0); 
			 \draw[->] (\xmin,0)--(\xmax,0) node[below]{$x$};
			\draw[->] (0,\ymin)--(0,\ymax) node[left]{$y$}; 
			 \fill (O) node[below left]{$O$} circle (1pt);
			\clip ({\xmin-0.1},{\ymin-0.1}) rectangle ({\xmax+0.1},{\ymax+0.1}); 
			\foreach \x  in {2}{
				 \fill (\x,0) node[below]{$\x$} circle (1pt);
			} 
			 \foreach \y in {2}{
				\fill (0,\y)  node[left]{$\y$} circle (1pt);
			 }
			 \draw[samples=100] plot[domain=\xmin:0.99](\x,{(\x-2)/(\x-1)});
			 \draw[samples=100] plot[domain=1.01:\xmax](\x,{(\x-2)/(\x-1)});
			 \draw[dashed] (1.0,\ymin)--(1.0,\ymax);
			\draw[dashed] (\xmin,1.0)--(\xmax,1.0); 
			\draw (1,-0.3) node[left]{$1$} 
				 (-0.2,1) node[above]{$1$};
		\end{tikzpicture}
	\end{center}
	\choiceTF 
		{\True Đồ  thị hàm số có tiệm cận đứng $x=1$}
		{Đồ thị hàm số  có tiệm cận ngang $y=0$}
		{Hàm số đồng biến trên $\mathbb{R}$} 
		 {\True $T=a-3b-2c=-3$}
	 \loigiai{
		\begin{itemchoice}
			 \itemch
			Đồ thị hàm số  có tiệm cận đứng $x=1$.
			 \itemch
			Đồ thị hàm số  có tiệm cận ngang $y=0$.
			 \itemch
			Hàm số đồng biến trên các  khoảng $(-\infty;1)$ và $(1;+\infty)$.
			\itemch 
			Đồ thị hàm số có tiệm cận ngang  $y=1$ nên $\dfrac{1}{b}=1\Rightarrow b=1$.\\
			 Đồ thị hàm số có tiệm cận đứng $x=1$ nên $-\dfrac{c}{b}=1$ mà $b=1\Rightarrow c=-1$.\\
			Đồ thị hàm số cắt trục tung tại điểm  $(0;2)$ nên $\dfrac{a}{c}=2$ mà $c=-1\Rightarrow a=-2$.\\
			Vậy $T=a-3b-2c=-2-3\cdot 1-2\cdot  (-1)=-3$.
		\end{itemchoice}
	} 
\end{ex}

% nguồn: Phan1_Lop_12_Dai_so_Chuong_1
\begin{ex}%[2D1H5-1]%[2D1H2-2]%[ma:b5b25c2e]
	\immini{Cho hàm số $y = f(x)  = x^3 + ax^2 + bx + c$ có đồ thị như hình.
	\choiceTF 
		{\True Hàm số $y = f(x)$ có  hai điểm cực trị là~$0$ và $2$}
		{\True  Giá trị $b$ bằng $0$}
		{Giá trị $c  = -2$}
		{$f(x)  = x^3 - 6x^2 + 2$}
		}{\resizebox{!}{4cm}{\begin{tikzpicture}[line join=round,line cap=round,>=stealth]
		\draw[->] (-2.6,0)--(3.6,0)  node[below left]{$x$};
		\draw[->] (0,-3.1)--(0,3.1)  node[below left]{$y$};
		 \draw (0,0) node [below left]{$O$};
		\foreach \x/\nx in {2/2} 
		\draw[thin]  (\x,1pt)--(\x,-1pt) node [above]{$\nx$};
		 \foreach \y/\ny in {2/2}
		 \draw[thin] (1pt,\y)--(-1pt,\y) node [above left]{$\ny$};
		 \draw[dashed,thin](2,0)--(2,-2);
		\begin{scope}
			\clip  (-2.5,-3) rectangle (3.5,3);
			\draw[samples=100,domain=-2.5:3.5,smooth,variable=\x]  plot (\x,{1*((\x)^3)+-3*((\x)^2)+0*(\x)+2});
		\end{scope}
		\path (2,2) node{$y=f  (x)$};
	\end{tikzpicture}
		}} 
	\loigiai{\begin{itemchoice}
		\itemch  Hàm số $y = f(x)$ có điểm cực tiểu là~$x = 2$, điểm cực đại là~$x = 0$.
		\itemch Ta có $f'(x) = 3x^2 + 2ax + b$. Vì $0$,  $2$ là~hai nghiệm của phương trình $f'(x) = 0$ nên $b = 0$, $a = -3$.
		\itemch Vì đồ thị hàm số đi qua điểm có tọa độ $(0;  2)$ nên $c = 2$.
		\itemch Ta có $a=-3$, $b=0$, $c=2$ suy ra $f(x) = x^3 - 3x^2 + 2$. 
	\end{itemchoice}
 }
\end{ex}

% nguồn: Phan1_Lop_12_Dai_so_Chuong_1
\begin{ex}%[2D1H5-1]%[ma:d38d7484]
	Cho hàm số $y=f(x)=ax^3+bx^2+cx+d$ có đồ thị như hình vẽ  dưới đây
	\begin{center}
		\begin{tikzpicture}[scale=0.8,font=\footnotesize,line join=round,line cap=round,>=stealth]
			\draw[->]  (-2.5,0)--(3.5,0) node[below right] {$x$};
			\draw[->] (0,-2)--(0,4) node[left]  {$y$};
			\draw[domain=-2.1:2.1,samples=200] plot (\x,{\x*\x*\x  - 3*\x + 1});
			 \draw[fill] (-1,3) circle (1.3pt);
			\draw[fill] (1,-1) circle (1.3pt); 
			\draw[fill] (0,-1)  circle (1.3pt);
			 \draw[fill] (-1,0) circle (1.3pt);
			\draw[fill] (0,1) circle (1.3pt)  node[right]{$1$};
			\draw[fill]  (1,0) circle (1pt) node[above]{$1$};
			\draw[fill] (0,3) circle  (1pt) node[right]{$3$};
			\draw[dashed] (-1,0) -- (-1,3) --  (0,3);
			\draw[dashed]  (1,0) -- (1,-1) -- (0,-1);
			\foreach \x in {-1} \draw  (\x,0.1)--(\x,-0.1) node[below] {\x};
			\foreach \y in {-1} \draw (0.1,\y)--(-0.1,\y)  node[left] {\y};
		\end{tikzpicture}
	\end{center}
	 \choiceTF
		{\True Đồ thị hàm số cắt trục  tung tại điểm có tọa độ $(0;1)$}
		{\True Đường thẳng đi qua điểm $(0;1)$ luôn cắt đồ thị tại ba điểm phân biệt  có hoành độ lập thành một cấp số cộng}
		{\True  $a - b + c + d = -1$}
		{Đồ  thị hàm số đi qua điểm $(3;18)$}
	 \loigiai{
		\begin{itemchoice}
			\itemch Đồ  thị có hai điểm cực trị là $(-1;3)$ và $(1;-1)$,có tâm đối xứng là $(0;1)$. Do đó, đồ thị cắt trục tung tại $(0;1)$.
			\itemch Do $(0;1)$ là tâm đối xứng của đồ  thị nên đường thẳng qua điểm đó luôn cắt đồ thị tại ba điểm phân biệt $I$, $A$, $B$ với $I$ là trung điểm của $A$, $B$, suy ra $x_A+x_B=2x_I$.\\
			Do đó hoành độ $3$ điểm này  lập thành cấp số cộng.
			\itemch Ta  có:
			 $$
			 \heva{&
			 f(-1) = 3\\
			&f(1) =  -1\\
			&f'(1) =  0\\
			&f'(-1)  = 0}
			 \Leftrightarrow
			\heva{& 
			 -a + b - c + d = 3\\
			&a + b + c + d  = -1\\
			&3a +  2b + c = 0\\
			&3a -  2b + c = 0}
			\Leftrightarrow \heva{&a = 1\\& b = 0\\&c = -3\\&d = 1.} 
			$$ 
			Vậy $a -  b + c + d = 1 - 0 - 3 + 1 = -1$.
			\itemch Hàm số tìm được là  $y = x^3 - 3x + 1$, ta có
			 $$
			f(3) = 27 - 9 + 1  = 19 \ne 18.
			 $$
			Vậy đồ  thị không đi qua điểm $(3;18)$.
		\end{itemchoice}
	 }
\end{ex}

% nguồn: Phan1_Lop_12_Dai_so_Chuong_1
\begin{ex}%[2D1H5-1]%[ma:43d9abff]
	Cho hàm  số $y=\dfrac{mx-1}{2x+m}$ có đồ thị là $\left(C_m \right)$ với $m$ là tham số
	\choiceTF 
		{\True Khi $m=2$ thì đồ thị hàm số  có đường tiệm cận ngang $y=1$}
		{Khi $m=2$  thì giao điểm của các đường tiệm cận có toạ độ $I\left(1;-1\right)$}
		{\True Đường tiệm cận đứng của đồ  thị hàm số đi qua điểm $A\left(-1; \sqrt{2} \right)$ thì $m=2$}
		{\True Với mọi giá trị của tham số $m$ thì hàm số luôn đồng biến trên mỗi khoảng  xác định của nó}
	\loigiai{ 
		\begin{itemchoice}
			 \itemch
			Khi $m=2$ thì đồ thị  hàm số có đường tiệm cận ngang $y=1$.
			\itemch 
			Khi $m=2$ thì giao điểm của các đường  tiệm cận có toạ độ $I\left(-1; 1\right)$.
			 \itemch
			Đồ thị hàm số có tiệm cận đứng khi và chỉ khi $m\left(-\dfrac{m}{2} \right)-1\ne 0\Leftrightarrow m^2+2\ne 0$  đúng với mọi giá trị $m$.\\
			Khi đó, TCĐ là đường thẳng $x=-\dfrac{m}{2}$. Vì TCĐ đi qua điểm $A(-1; \sqrt{2})$ nên ta có $-\dfrac{m}{2}=-1\Leftrightarrow m=2$. 
			 \itemch
			Ta có $y'=\dfrac{m^2+2}{\left(2x+m\right)^2}  > 0, \forall x\ne-\dfrac{m}{2}$.\\
			Vậy hàm số đồng biến trên từng khoảng xác định  của nó.
		\end{itemchoice}
	} 
\end{ex}

% nguồn: Phan1_Lop_12_Dai_so_Chuong_1
\begin{ex}%[2D1H5-1]%[ma:d4a079e4]
	Cho hàm  số $f(x)=\dfrac{ax+b}{cx+d}$ với $a, b, c, d\in \mathbb{R}$ có đồ thị hàm số $y=f'(x)$ nhận $x=-1$ làm tiệm cận đứng như hình vẽ bên.
	\begin{center}
		\begin{tikzpicture}[scale=0.8,font=\footnotesize,line join=round,line cap=round,>=stealth]
			\def\xmin{-4}; \def\ymin{-1}; \def\xmax{3}; \def\ymax{6}; 
			 \draw[->] (\xmin,0) -- (\xmax,0) node[below]{$x$};
			\draw[->] (0,\ymin) -- (0,\ymax) node[right]{$y$}; 
			\fill (0,0)  node[below left]{$O$} circle (1.5pt);
			\clip ({\xmin-0.1},{\ymin-0.1})  rectangle ({\xmax+0.1},{\ymax+0.1});
			\draw[dashed,red]  (-1,\ymin) -- (-1,\ymax);
			\draw[blue,samples=100,domain=\xmin:-1.1]  plot (\x,{3/((\x+1)*(\x+1))});
			\draw[blue,samples=100,domain=-0.9:\xmax] plot  (\x,{3/((\x+1)*(\x+1))});
			\fill (0,3) node[left]{$3$}  circle (1.5pt);
			\fill (-1,0) circle  (1.5pt);
			 \draw (-1.3,0) node[below]{$-1$};
		\end{tikzpicture}
	\end{center}
	Biết rằng giá trị lớn nhất của hàm số $y=f(x)$ trên  đoạn $\left[-3;-2\right]$ bằng $8$.
	 \choiceTF
		{\True $f'(0)=3$} 
		{Hàm số $y=f(x)$ nghịch  biến trên khoảng $\left(-1;+\infty \right)$}
		 {Giá trị của $f\left(-3\right)$ bằng $8$}
		{\True Giá trị của $f(2)$ bằng $4$} 
	\loigiai{ 
		\begin{itemchoice}
			\itemch 
			 Theo hình vẽ, đồ thị $f'(x)$ qua điểm $\left(0; 3\right)$ nên $f'(0)=3$.
			 \itemch
			Do $f'(x) > 0, \forall x\ne-1$ nên hàm số đã cho $f(x)$ đồng biến trên các  khoảng $\left(-\infty;-1\right)$ và $\left(-1;+\infty \right)$.
			\itemch 
			Do $f'(x)> 0, \forall x\ne-1\Rightarrow \max\limits_{\left[-3;-2\right]} f(x)=f\left(-2\right)=8$ suy ra $f\left(-3\right)\ne 8$. 
			 \itemch
			 Ta có $f'(x)=\dfrac{ad-bc}{\left(cx+d\right)^2}$.\\
			Đồ thị hàm số $y=f'(x)$ đi qua điểm $\left(0; 3\right)$ nên  $f'(0)=3\Leftrightarrow \dfrac{ad-bc}{d^2}=3$.\\
			Mặt khác,  đồ thị hàm số $y=f'(x)$ có tiệm cận đứng $x=-1$ nên $-c+d=0$.\\
			 Vì $f'(x)>0, \forall x\ne 1\Rightarrow \max\limits_{\left[-3;-2\right]} f(x)=f\left(-2\right)=8\Leftrightarrow \dfrac{-2a+b}{-2c+d}=8$.\\
			Vậy ta có hệ phương trình $\heva{&ad-bc=3d^2 \\ 
			 &-c+d=0 \\
			&b-2a=8\left(d-2c\right) 
			} \Leftrightarrow \heva{&c=d  \\
			&a-b=3d \\ 
			&b-2a=-8d 
			 } \Leftrightarrow \heva{&a=5d \\
			&b=2d  \\
			 &c=d.
			}$\\ 
			Từ đó suy ra $f(x)=\dfrac{5dx+2d}{dx+d}=\dfrac{5x+2}{x+1} \Rightarrow  f(2)=4$.
		\end{itemchoice}
	} 
\end{ex}

% nguồn: Phan1_Lop_12_Dai_so_Chuong_1
\begin{ex}%[2D1H5-1]%[ma:f7953c3b]
	Cho hàm số $y=\dfrac{ax+b}{x+c}$ ($a$, $b$, $c\in\mathbb{Z}$) có đồ thị như hình dưới 
	\begin{center}
		\begin{tikzpicture}[scale=1,font=\footnotesize,line join=round,line cap=round,>=stealth]
			\draw[->] (-3,0)--(0,0)  node[below left]{$O$}--(4,0) node[below left]{$x$};
			\draw[->] (0,-3.5)--(0,3) node[below  left]{$y$};
			\clip (-3,-3.5)  rectangle (4,3);
			\draw[smooth,samples=100]plot[domain=-3:4](\x,{(-1*(\x)+2)/((\x)-1)}) 
				(-3,-1)--(4,-1); 
			\foreach \x/\y in  {-2/0,-1/0,1/0,2/0,3/0}{\node[below] at (\x,\y) {$\x$};};
			\foreach \x/\y in {0/-3,0/-2,0/-1,0/1,0/2}{\node[left] at  (\x,\y) {$\y$};};
		\end{tikzpicture}
	\end{center}
	Khi đó: 
	\choiceTF 
		{\True Đường thẳng $x=1$ là đường tiệm cận đứng  của đồ thị hàm số}
		{Hàm số đồng  biến trên từng khoảng xác định}
		{\True Hàm số đạt giá trị nhỏ nhất  bằng $-2$ trên đoạn $[-2;0]$}
		{Các số $a$, $b$, $c$  đều âm}
	\loigiai{ 
		\begin{itemchoice}
			\itemch Đúng. Vì từ đồ thị ta thấy $x=1$ là đường tiệm cận đứng của đồ thị hàm  số.
			\itemch Sai. Vì từ đồ thị ta thấy đồ thị đi xuống nên hàm số  nghịch biến trên từng khoảng xác định.
			\itemch Đúng. Vì hàm số nghịch biến trên đoạn  $[-2;0]$ nên $\min\limits_{[-2;0]}y=y(0)=-2$.
			\itemch Sai.  Vì
			\begin{itemize}
				\item Từ  dồ thị ta có TCĐ là đường thẳng $x=1\Rightarrow 1+c=0\Rightarrow c=-1$
				 \item TCN là đường thẳng $y=-1\Rightarrow a=-1$
				\item Đồ  thị cắt trục $Oy$ tại điểm $(0;-2)\Rightarrow-2=\dfrac{a.0+b}{0+c}\Rightarrow b=-2c=2$.
			\end{itemize}
		\end{itemchoice}
	} 
\end{ex}

% nguồn: Phan1_Lop_12_Dai_so_Chuong_1
\begin{ex}%[2D1H5-1]%[ma:b6319406]
	 Cho hàm số $\mathrm{f}(x)=2x+\dfrac{2}{x-1}$.
	 \renewcommand{\dotEX}{}
	\choiceTF 
		{$f'(x)=1-\dfrac{1}{(x-1)^2}$.} 
		{\True Phương trình $\mathrm{f}'(x)=0$ có hai nghiệm phân biệt là $0$ và  $2$.}
		{\True Hàm số đã cho có bảng  biến thiên như sau
		\begin{center}
			\resizebox{0.7\linewidth}{!}{\begin{tikzpicture}[>=stealth]
		\tkzTabInit[nocadre,lgt=1.2,espcl=3] 
			{$x$ /1,$y'$  /1,$y$ /2.5}
			{$-\infty$,$0$,$1$,$2$,$+\infty$} 
		\tkzTabLine{,+,,-,,-,,+,} 
		\tkzTabVar{-/$-\infty$,+/ $-2$,-D+/ $-\infty$/$+\infty$,-/  $6$,+/ $+\infty$}
	\end{tikzpicture}
			 }
	\end{center}}
	{Hàm số đã cho có đồ thị như  sau
	\begin{center}
		\resizebox{!}{5cm}{\begin{tikzpicture}[>=stealth]
		\draw[cyan!50,thin] (-4.5,-4) grid  (5,5);
		 \draw[->] (-4,0)--(0,0)
			node[below right]{$O$}--(5,0) node[below]{$x$}; 
		\draw[->] (0,-4) --(0,5.5)  node[right]{$y$};
		\foreach \x  in {-4,-3,-2,-1,1,2,3,4}{
			 \draw (\x,-.1)--(\x,.1) node[above]{$\x$};
		 }
		\foreach \y in  {-4,-3,-2,-1,1,2,3,4,5}{
			 \draw (-.1,\y)--(.1,\y) node[right]{$\y$};
		 }
		\draw[samples=100,domain=-4.1:0.89]  plot (\x,{1/2*(\x)+1/(2*(\x-1))});
		\draw[samples=100,domain=1.11:4.5]  plot (\x,{1/2*(\x)+1/(2*(\x-1))});
		\draw  [domain=-4.5:5,dashed] plot (\x,{1/2*(\x)});
		\draw  [dashed] (1,-4.1)--(1,5.1);
		\draw[fill] (1,0.5) circle (1pt); 
	\end{tikzpicture}
		} 
	\end{center}
	} 
	\loigiai{ 
		\begin{itemchoice}
			\itemch  Sai. Vì $f'(x)=2-2\dfrac{1}{(x-1)^2}$.
			\itemch Đúng. Ta dễ dàng đạo hàm và giải phương trình $f'(x)=0$ và có $2$ nghiệm  phân biệt là $0$ và $2$.
			\itemch Đúng. Vì ta có thể suy luận được tiệm cận  đứng là $x=1$ và $f'(x)=0$ và có $2$ nghiệm phân biệt là $0$ và $2$.
			\itemch Sai. Vì thay $x=2$  vào hàm số $\mathrm{f}(x)=2x+\dfrac{2}{x-1}$ ta có giá trị là $6$.
		\end{itemchoice}
	} 
\end{ex}

% nguồn: Phan1_Lop_12_Dai_so_Chuong_1
\begin{ex}%[2D1H5-1]%[ma:a0513258]
	Cho  hàm số $f(x)=\dfrac{ax-1}{bx+c}$ với $a, b, c\in \mathbb{R}$ có bảng biến thiên như hình vẽ dưới đây
	\begin{center}
		\begin{tikzpicture}[scale=1,font=\footnotesize]
			 \tkzTabInit[nocadre=false,lgt=1.2,espcl=2.5,deltacl=0.7]
				{$x$/0.8,$f'(x)$/0.6,$f (x)$/2} 
				 {$-\infty$,$3$,$+\infty$};
			\tkzTabLine{,-,d,-,}; 
			\tkzTabVar{+/$\dfrac{1}{2}$,-D+/$-\infty$/$+\infty$,-/$\dfrac{1}{2}$}; 
		\end{tikzpicture}
	\end{center}
	\def\dotEX{} 
	\choiceTF 
		{\True Hàm  số nghịch biến trên khoảng $\left(-\infty; \dfrac{1}{2} \right)$.}
		 {Đồ thị hàm số có tiệm cận đứng $x=\dfrac{1}{2}$.}
		{\True  Đồ thị giao với trục hoành tại điểm có hoành độ nhỏ hơn $3$.}
		{\True  $\hoac{&b>\dfrac{2}{3}\\&b<0.}$}
	\loigiai{ 
		\begin{itemchoice}
			 \itemch
			Hàm số nghịch biến trên khoảng $\left(-\infty;3\right)$ nên nghịch biến  trên khoảng $\left(-\infty;\dfrac{1}{2}\right)$.
			 \itemch
			Đồ thị hàm số có tiệm  cận đứng $x=3$.
			 \itemch
			Đồ thị giao với trục hoành tại điểm thuộc nhánh  trái của đồ thị, suy ra hoành độ giao điểm này nhỏ hơn $3$.
			\itemch 
			Từ bảng biến thiên, ta thấy đồ thị có tiệm cận  đứng $x=-\dfrac{c}{b}=3$ và tiệm cận ngang $y=\dfrac{a}{b}=\dfrac{1}{2}$.\\
			Suy ra  $\heva{&\dfrac{a}{b}=\dfrac{1}{2}\\&-\dfrac{c}{b}=3}\quad (1)$.\\
			Ta có  $y'=\dfrac{ac+b}{\left(bx+c\right)^2}<0, \forall x\ne-\dfrac{c}{b} \Leftrightarrow ac+b<0\quad (2)$.\\
			Từ (1) và (2) suy ra $\dfrac{b}{2}\cdot (-3b)+b<0\Leftrightarrow \hoac{&b>\dfrac{2}{3}\\&b<0.}$ 
		\end{itemchoice}
	 }
\end{ex}

% nguồn: Phan1_Lop_12_Dai_so_Chuong_1
\begin{ex}%[2D1H5-1]%[ma:8f8084a3]
	\immini{Cho hàm số $y = \dfrac{ax^2 + bx + c}{mx + n}$ có đồ thị như hình  vẽ bên
	\choiceTF 
		{Hàm số đã cho nghịch  biến trên khoảng $\left(-2; 0 \right)$}
		{\True Đồ thị của hàm  số đã cho có tiệm cận xiên $y = x+1$}
		{Gọi $A$, $B$ là~hai cực trị của hàm  số đã cho, diện tích tam giác $OAB$ bằng $8$ (với $O$ là~gốc toạ độ)}
		{\True Một trục đối xứng của đồ thị đã cho là~$d \colon y = \left(x +  1\right) \tan{\dfrac{3\pi}{8}}$}
		}{\resizebox{!}{5cm}{\begin{tikzpicture}[line join = round,line cap=round,>=stealth]
		 \draw[->]
			(-4,0) --  (3,0) node[below]{$x$};
		\draw[->] 
			(0,-4) -- (0,3.5) node[right]{$y$}; 
		\begin{scope}
			 \clip (-5,-4) rectangle (3,3.5);
			 \draw (-1,-4) -- (-1,4);
			 \draw[samples=100,domain=-0.9:3,smooth,variable=\x] plot (\x,{((\x)^2+2*(\x)+2)/((\x)+1)});
			 \draw[samples=100,domain=-4:-1.1,smooth,variable=\x] plot (\x,{((\x)^2+2*(\x)+2)/((\x)+1)});
			\draw[samples=100,domain=-4:3.5,smooth,variable=\x]  plot (\x,{(\x)+1});
		\end{scope}
		\fill 
			(0,0) circle  (1pt) node[below right]{$O$}
			(-2,-2) circle (1pt) 
			(0,2) circle  (1pt)
			(-2,0) circle  (1pt) node[above left]{$-2$}
			(-1,0) circle (1pt)  node[above left]{$-1$}
			(0,-2) circle (1pt)  node[right]{$-2$}
			(0,1) circle  (1pt) node[right]{$1$}
			(0,2)  circle (1pt) node[below left]{$2$};
		 \draw[dashed]
			(-2,0) -- (-2,-2) --  (0,-2);
	\end{tikzpicture}
		}} 
	\loigiai{\begin{itemchoice}
		\itemch \textbf{Sai}.\\ 
		Dựa vào đồ thị hàm số ta thấy hàm số nghịch biến trên $\left(-2;  -1 \right)$ và $\left(-1; 0 \right)$.
		\itemch \textbf{Đúng}.\\ 
		Gọi phương trình đường tiệm cận xiên của đồ thị hàm  số là~$d\colon y = a_1x + b_1 \left(a_1 \neq 0 \right)$.\\
		Ta thấy $d$ đi qua các điểm có tọa độ $\left(0;1 \right)$, $\left(-1; 0 \right)$ nên $\heva{&1 = b\\&-a +  b = 0} \Leftrightarrow \heva{&a = 1\\&b=1.}$\\
		Tiệm cận xiên của đồ thị hàm số có phương trình là~$y =  x+1$.
		\itemch  \textbf{Sai}.\\
		Ta có $2$ điểm cực trị của hàm số  đã cho có tọa độ là~$A\left(0;2 \right)$, $B\left(-2; -2 \right)$.\\
		Ta có $AO = 2$,  $BO = 2\sqrt{2}$ và $AB = 2\sqrt{5}$.\\
		Nửa chu vi của $\triangle ABO$ bằng $p  = \dfrac{2 + 2\sqrt{2}+ 2\sqrt{5}}{2}$.\\
		Suy ra $S_{OAB}  = \sqrt{\mathrm{P}\left(p - 2\right) \left(p-2\sqrt{2}\right)\left(p - 2\sqrt{5}\right)} = 2$.
		 \itemch \textbf{Đúng}.\\
		Từ đô thị hàm số ta thấy tiệm cận đứng của đồ thị hàm số  là~$\Delta \colon x = -1$.\\
		Trục đối xứng của đồ thị hàm số đã cho chính là~tia phân giác của góc tạo bởi tiệm cận xiên $d \colon x - y + 1 =  0$ và tiệm cận đứng $\Delta \colon x + 1 = 0$.\\
		Phương  trình đường phân giác góc tạo bởi $\left(d, \Delta \right)$ là
		\begin{align*}
			\dfrac{x-y+1}{\sqrt{2}} = \pm \left(x+1 \right) \Leftrightarrow \hoac{&x - y+1 = \sqrt{2}\left(x+1 \right)\\&x-y+1=-\sqrt{2}\left(x+1 \right)} \Leftrightarrow \hoac{&y = \left(1-\sqrt{2} \right)\left(x+1 \right)\\&y = \left(1+\sqrt{2} \right) \left(x+1 \right) =  \left(x+1 \right) \tan{\dfrac{3\pi}{8}}.}
		\end{align*}
	\end{itemchoice}
} 
\end{ex}

% nguồn: Phan1_Lop_12_Dai_so_Chuong_1
\begin{ex}%[2D1H5-1]%[ma:8ae3a43b]
	Cho hàm số $y=f(x)=ax^3+bx^2+cx+d$ có đồ thị như hình vẽ dưới  đây
	\begin{center}
		\begin{tikzpicture}[scale=0.8,font=\footnotesize,line join=round,line cap=round,>=stealth]
			\draw[->] (-2.5,0)--(2,0) node[below right] {$x$}; 
			\draw[->] (0,-1)--(0,3.5)  node[left] {$y$};
			\draw[domain=-2.3:0.8,samples=200] plot (\x,{\x*\x*\x  +2*\x*\x + 1});
			\draw[fill]  (-2,1) circle (1.3pt);
			 \draw[fill] (-1,2) circle (1.3pt);
			\draw[fill] (0,1)  circle (1.3pt);
			\draw[fill] (-2,0) circle (1.3pt); 
			\draw[fill] (-1,0)  circle (1.3pt);
			 \draw[fill] (0,2) circle (1.3pt) node[right]{$2$};
			 \draw[fill] (0,0) circle (1.3pt) node[below right]{$O$};
			\draw[dashed]  (-2,0) -- (-2,1) -- (0,1);
			\draw[dashed] (-1,0) -- (-1,2)  -- (0,2);
			\draw[dashed] (0,0) --  (0,1);
			\foreach \x in {-2,-1} \draw (\x,0.1)--(\x,-0.1) node[below]  {\x};
			\foreach \y  in {1} \draw (0.1,\y)--(-0.1,\y) node[below right] {\y};
		\end{tikzpicture}
	\end{center}
	 \choiceTF
		{Hàm số đạt cực tiểu tại  $x = 1$}
		{\True Đồ thị hàm số  cắt trục $Oy$ tại điểm có tọa độ $(0;1)$}
		{Hàm số  đồng biến trên khoảng $(-\infty; -1)$}
		{$2a + 3b +  c = 9$}
	 \loigiai{
		\begin{itemchoice}
			\itemch Hàm số đạt cực tiểu tại $x = 0$, giá trị cực tiểu  là $y = 1$.
			\itemch Đồ thị  hàm số cắt trục tung tại điểm có tọa độ $(0;1)$.
			\itemch Hàm số đồng biến trên khoảng $(-\infty; x_0)$ với $-2  < x_0 < -1$.
			\itemch Đồ thị đi qua các điểm $(-2;1)$, $(-1;2)$, $(0;1)$ và đạt cực  trị tại $x = 0$ nên ta có hệ phương trình
			$$\heva{&f(-2)=1\\&f(-1)=2\\&f(0)=1\\&f'(0)=0}\Leftrightarrow 
			 \heva{
			&-8a + 4b - 2c + d  = 1\\&
			-a + b  - c + d = 2\\&
			d  = 1\\&
			 c = 0}
			\Leftrightarrow \heva{&a = 1\\& b = 2\\  &c = 0\\& d = 1} \Rightarrow 2a + 3b + c = 2 + 6 = 8.
			 $$
		\end{itemchoice}
	} 
\end{ex}

% nguồn: Phan1_Lop_12_Dai_so_Chuong_1
\begin{ex}%[2D1H5-1]%[ma:82c189b2]
	Cho hàm số $y=ax^3+bx^2+cx+d$ có bảng biến thiên như hình vẽ  dưới đây
	\begin{center}
		\begin{tikzpicture}[scale=1,font=\footnotesize]
			 \tkzTabInit[nocadre=false,lgt=1.2,espcl=2.5,deltacl=0.6]
				{$x$/1,$y'$/1,$y$/2} 
				 {$-\infty$,$-2$,$0$,$+\infty$}
			\tkzTabLine{,-,z,+,z,-,} 
			 \tkzTabVar{+/$+\infty$,-/0,+/4,-/$-\infty$}
		\end{tikzpicture}
	\end{center}
	 \choiceTF
		{Hàm số đạt giá trị  lớn nhất là $4$}
		{\True Đường thẳng $y = 2$ cắt đồ thị hàm số tại ba điểm phân  biệt}
		{\True Trong bốn hệ số $a$, $b$, $c$, $d$ có đúng hai  số âm}
		{\True Đồ thị  hàm số đi qua điểm $(-4;20)$}
	\loigiai{ 
		\begin{itemchoice}
			\itemch Hàm số không có giá trị lớn nhất  trên $\mathbb{R}$.
			\itemch Ta có bảng biến thiên 
			\begin{center}
				\begin{tikzpicture}[scale=1,font=\footnotesize]
					\tkzTabInit[nocadre=false,lgt=1.2,espcl=2.5,deltacl=0.6] 
						 {$x$/1,$y'$/1,$y$/2}
						 {$-\infty$,$-2$,$0$,$+\infty$}
					 \tkzTabLine{,-,z,+,z,-,}
					\tkzTabVar{+/$+\infty$,-/0,+/4,-/$-\infty$} 
					\draw  (2.2,-3) -- (9.2,-3);
					\node[above right] at (8.2,-3) {$y=2$}; 
				\end{tikzpicture}
			\end{center}
			Kẻ đường thẳng $y = 2$ đi qua điểm $(0;2)$ và song song  với trục $Ox$, đường này cắt đồ thị tại ba điểm phân biệt.
			\itemch Từ bảng  biến thiên ta có
			$$ 
			\heva{& f(-2) = 0\\ 
			&f(0)  = 4\\
			 &f'(-2) = 0\\
			&f'(0) =  0}
			\Leftrightarrow\heva{&-8a+4b-2c+d=0\\&d=4\\&12a-4b+c=0\\&c=0} \Leftrightarrow \heva{&a = -1\\& b = -3\\&c = 0\\ &d  = 4.}
			$$ 
			Vậy có đúng hai số âm  trong bốn hệ số.
			\itemch Với $a = -1,\ b = -3,\ c = 0,\  d = 4$, ta có hàm số $y = -x^3 - 3x^2 + 4$. Thay $x = -4$ ta được
			 $$
			y = -(-4)^3 - 3(-4)^2 + 4 =  64 - 48 + 4 = 20.
			 $$
			Nên đồ thị đi qua điểm $(-4;20)$. 
		\end{itemchoice}
	 }
\end{ex}

% nguồn: Phan1_Lop_12_Dai_so_Chuong_1
\begin{ex}%[2D1H5-1]%[ma:eb422168]
	 \immini{
	Cho  hàm số $y=\dfrac{x-2}{2x+1}$.
	\choiceTF 
		 {\True Tập xác định của hàm số là $\mathbb{R}\setminus \left\{-\dfrac{1}{2}\right\}$}
		{Đạo hàm của hàm  số là $y'=\dfrac{3}{(2x+1)^2}$}
		{Các đường  tiệm cận của hàm số là $x=\dfrac{1}{2}$, $y=-\dfrac{1}{2}$}
		{\True Đồ thị của hàm số có  dạng như hình bên}
		} 
		{\resizebox{!}{5.5cm}{\begin{tikzpicture}[line join=round,line cap=round,>=stealth]
		 \draw[->] (-4.1,0)--(5,0) node[below left] {$x$};
		\draw[->] (0,-4.1)--(0,5.1) node[below right] {$y$}; 
		\fill (0,0) circle (1pt)  node [below right] {$O$};
		\foreach \x/\gx  in {-3/-3,-2/-2,-1/-1,2/2,3/3,4/4}
		\draw (\x,1pt)--(\x,-1pt)  node [below] {$\gx$};
		\foreach \y/\gy  in {-3/-3,-2/-2,1/1,2/2,3/3,4/4}
		\draw (1pt,\y)--(-1pt,\y)  node [right] {$\gy$};
		 \draw[dashed,thin] (-0.49,-4)--(-0.49,5);
		\begin{scope}
			\clip (-5,-4) rectangle (5,5); 
			\draw[samples=100,domain=-4:-0.51,smooth,variable=\x] plot (\x,{((\x)+-2)/(2*(\x)+1)}); 
			 \draw[samples=100,domain=-0.49:4,smooth,variable=\x] plot (\x,{((\x)+-2)/(2*(\x)+1)});
			 \draw[dashed,thin] (-4,1/2)--(4,1/2);
		\end{scope}
	\end{tikzpicture}
		 }
		 }
	\loigiai{ 
	 \immini{
		\begin{itemchoice}
			\itemch Tập xác định của hàm  số là $\mathbb{R}\setminus \left\{-\dfrac{1}{2}\right\}$.
			\itemch Đạo  hàm của hàm số là $y'=\dfrac{5}{(2x+1)^2}$.
			\itemch Các đường tiệm cận  của hàm số là $x=-\dfrac{1}{2}$, $y=\dfrac{1}{2}$.
			\itemch  Đồ thị của hàm số có dạng như hình bên.
		\end{itemchoice}
		}{\resizebox{!}{5.5cm}{\begin{tikzpicture}[line join=round,line cap=round,>=stealth]
		\draw[->] (-4.1,0)--(5,0) node[below  left] {$x$};
		 \draw[->] (0,-4.1)--(0,5.1) node[below right] {$y$};
		\fill (0,0)  circle (1pt) node [below right] {$O$};
		\foreach \x/\gx in {-3/-3,-2/-2,-1/-1,2/2,3/3,4/4} 
		\draw (\x,1pt)--(\x,-1pt) node [below] {$\gx$}; 
		\foreach \y/\gy in  {-3/-3,-2/-2,1/1,2/2,3/3,4/4}
		\draw (1pt,\y)--(-1pt,\y) node [right] {$\gy$}; 
		 \draw[dashed,thin] (-0.49,-4)--(-0.49,5);
		\begin{scope}
			\clip (-5,-4) rectangle (5,5); 
			\draw[samples=100,domain=-4:-0.51,smooth,variable=\x] plot  (\x,{((\x)+-2)/(2*(\x)+1)});
			\draw[samples=100,domain=-0.49:4,smooth,variable=\x]  plot (\x,{((\x)+-2)/(2*(\x)+1)});
			\draw[dashed,thin]  (-4,1/2)--(4,1/2);
		\end{scope}
	\end{tikzpicture}
		 }}
	 }
\end{ex}

% nguồn: Phan1_Lop_12_Dai_so_Chuong_1
\begin{ex}%[2D1H5-1]%[ma:40e1543c]
	 Cho hàm số $f(x)=2 x^3-3x^2+1$.
	 \renewcommand{\dotEX}{}
	\choiceTF 
		{\True $f'(x)=6 x^2-6 x.$} 
		{\True Phương trình $f'(x)=0$ có hai nghiệm phân biệt  là 0 và 1.}
		{Hàm số đã cho có  bảng biến thiên như sau
		\begin{center}
			\begin{tikzpicture}
				 \tkzTabInit[nocadre,deltacl=0.6]{$x$/0.7,$f'(x)$/0.7,$f (x)$/2.5}{$-\infty$,$0$,$1$,$+\infty$}
				\tkzTabLine{,+,0,-,0,+,} 
				\tkzTabVar{-/$+\infty$,+/$1$,-/$0$,+/$-\infty$} 
			\end{tikzpicture}
	\end{center}}
	{\True Hàm số đã cho có đồ thị  như sau
	\begin{center}
		 \resizebox{!}{5cm}{
		\begin{tikzpicture}[>=stealth]
			\draw[cyan!50,thin] (-3,-3) grid (4,4); 
			\draw[->]  (-3,0)--(0,0) node[below right]{$O$}--(4,0) node[below]{$x$};
			 \draw[->] (0,-3) --(0,4) node[right]{$y$};
			 \foreach \x in {-2,-1,1,2,3}{
				\draw (\x,-.1)--(\x,.1) node[above]{$\x$}; 
			} 
			\foreach  \y in {-2,-1,1,2,3}{
				 \draw (-.1,\y)--(.1,\y) node[right]{$\y$};
			 }
			\draw[samples=100,domain=-0.9:1.9] plot  (\x,{2*(\x)^3-3*(\x)^2+1});
		\end{tikzpicture}
		} 
	\end{center}}
	 \loigiai{
		\begin{itemchoice}
			 \itemch Đúng. Ta dễ dàng đạo hàm để có kết quả $f'(x)=6 x^2-6 x$.
			\itemch Đúng. Ta dễ dàng đạo hàm và giải phương trình $f'(x)=0$ và có $2$ nghiệm. 
			 \itemch Sai. Vì
			\itemch Đúng. Vì căn cứ vào $2$ nghiệm  khi $f'(x)=0$ và hệ số $a>0$ của hàm số bậc ba.
		\end{itemchoice}
	} 
\end{ex}

% nguồn: Phan1_Lop_12_Dai_so_Chuong_1
\begin{ex}%[2D1H5-1]%[ma:18f50d8b]
	Cho hàm số $y=f(x)=ax^3+bx^2+cx+d$ có đồ thị như hình vẽ  dưới đây
	\begin{center}
		\begin{tikzpicture}[scale=0.8,font=\footnotesize,line join=round,line cap=round,>=stealth]
			\draw[->] (-2.5,0)--(3.5,0) node[below right] {$x$}; 
			\draw[->] (0,-2)--(0,4) node[left]  {$y$};
			\draw[domain=-2.1:2.1,samples=200] plot (\x,{\x*\x*\x  - 3*\x + 1});
			 \draw (-1,3) circle (1.3pt);
			\draw  (1,-1) circle (1.3pt);
			\draw (0,-1) circle (1.3pt); 
			\draw (-1,0)  circle (1.3pt);
			\draw  (0,1) circle (1.3pt) node[right]{$1$};
			 \draw (1,0) circle (1pt) node[above]{$1$};
			\draw (0,3)  circle (1pt) node[right]{$3$};
			\draw[dashed]  (-1,0) -- (-1,3) -- (0,3);
			\draw[dashed] (1,0) -- (1,-1) -- (0,-1); 
			\foreach \x in {-1} \draw (\x,0.1)--(\x,-0.1) node[below] {\x}; 
			\foreach \y in {-1} \draw (0.1,\y)--(-0.1,\y) node[left]  {\y};
		\end{tikzpicture}
	\end{center}
	 \choiceTF
		{\True Đồ thị hàm số cắt trục tung tại điểm có tọa  độ $(0;1)$}
		{\True Đường thẳng đi qua điểm $(0;1)$ luôn cắt đồ  thị tại ba điểm phân biệt có hoành độ lập thành một cấp số cộng}
		{\True $a - b + c + d =  -1$}
		{Đồ thị hàm số đi qua điểm $(3;18)$} 
	\loigiai{ 
		\begin{itemchoice}
			\itemch Đồ thị có hai điểm cực trị là $(-1;3)$ và $(1;-1)$,có  tâm đối xứng là $(0;1)$. Do đó, đồ thị cắt trục tung tại $(0;1)$.
			\itemch Do $(0;1)$ là tâm đối xứng của đồ thị nên đường thẳng qua điểm đó luôn cắt đồ thị tại ba điểm phân biệt $I$, $A$, $B$ với $I$ là trung  điểm của $A$, $B$, suy ra $x_A+x_B=2x_I$.\\
			Do đó hoành  độ $3$ điểm này lập thành cấp số cộng.
			\itemch  Ta có:
			$$ 
			\heva{& 
			 f(-1) = 3\\
			&f(1)  = -1\\
			&f'(1) = 0\\ 
			 &f'(-1) = 0}
			\Leftrightarrow 
			\heva{& 
			-a + b - c + d = 3\\ 
			&a + b  + c + d = -1\\
			&3a + 2b + c = 0\\ 
			&3a - 2b + c =  0}
			\Leftrightarrow \heva{&a = 1\\& b = 0\\&c = -3\\&d =  1.}
			 $$
			Vậy $a - b + c + d  = 1 - 0 - 3 + 1 = -1$.
			\itemch Hàm số tìm được là $y = x^3 - 3x + 1$,  ta có
			 $$
			f(3) = 27 - 9  + 1 = 19 \ne 18.
			$$ 
			Vậy đồ thị không  đi qua điểm $(3;18)$.
		\end{itemchoice}
	 }
\end{ex}

% nguồn: Phan1_Lop_12_Dai_so_Chuong_1
\begin{ex}%[2D1H5-1]%[ma:81ad9572]
	Cho một hàm số có đồ  thị như sau
	\begin{center}
		\begin{tikzpicture}[line join=round,line cap=round,>=stealth]
			\tikzset{ethicknode/.style={scale=0.9}} 
			 \draw[->] (-3.1,0)--(4.1,0) node[below left] {$x$};
			\draw[->] (0,-3.1)--(0,4.1)  node[below left] {$y$};
			\draw (0,0)  node [below left] {$O$};
			\foreach \x/\nx in {1/1,-1/-1} 
			\draw[thin] (\x,1pt)--(\x,-1pt) node [below left]  {$\nx$};
			\foreach \y/\ny in {1/1,-1/-1} 
			\draw[thin] (1pt,\y)--(-1pt,\y)  node [below left] {$\ny$};
			\draw[dashed,thin]  (1.01,-3)--(1.01,4);
			\begin{scope}
				\clip (-3,-3) rectangle (4,4); 
				\draw[samples=200,domain=-3:0.99,smooth,variable=\x] plot (\x,{(1*(\x)+1)/(1*(\x)+-1)}); 
				 \draw[samples=200,domain=1.01:4,smooth,variable=\x] plot (\x,{(1*(\x)+1)/(1*(\x)+-1)});
				\draw[dashed,thin]  (-3,1/1)--(4,1/1);
			\end{scope}
		\end{tikzpicture}
	\end{center}
	\choiceTF 
		{Đây là đồ thị của một hàm số đồng biến trên  mỗi khoảng xác định}
		{\True Đồ thị có tiệm cận đứng là đường thẳng $x=1$} 
		{Đồ thị có tiệm  cận ngang là đường thẳng $y=0$}
		{\True Đường cong trên là đồ thị  của hàm số $y=\dfrac{x+1}{x-1}$}
	\loigiai{ 
		\begin{itemchoice}
			\itemch Vì nhìn từ trái sang  phải đồ thị hướng xuống nên hàm số nghịch biến trên mỗi khoảng xác định.
			\itemch 
			\itemch Vì nhìn vào đồ  thị ta thấy tiệm cận ngang là đường thẳng $y=1.$
			 \itemch
		\end{itemchoice}
	 }
\end{ex}

% nguồn: Phan1_Lop_12_Dai_so_Chuong_1
\begin{ex}%[2D1H5-1]%[ma:c5bda32b]
	\immini{Cho hàm số  $y=ax^3+bx^2+cx+d$ có đồ thị như hình bên
	 \choiceTF
		{Hàm số có hai  điểm cực trị trái dấu}
		{\True Tổng giá trị cực đại  và cực tiểu là số âm}
		{Phương trình $y' = 0$ có ba  nghiệm phân biệt}
		{\True Trong các hệ số $a$, $b$, $c$, $d$  có $2$ hệ số dương}}
		{\begin{tikzpicture}[scale=0.4,font=\footnotesize,line join=round,line cap=round,>=stealth]
		\draw[->] (-3.2,0)--(5.2,0)  node[below] {$x$};
		 \draw[->] (0,-6.5)--(0,3.2) node[left] {$y$};
		 \draw[domain=-0.6:4.59,samples=200] plot (\x,{(\x-2)*(\x-2)*(\x-2) - 5*(\x-2)-2});
		\draw[fill] (0,0) circle (1pt) node[below  left] {$O$};
	\end{tikzpicture}
		} 
	\loigiai{ 
		Ta có $y=ax^3+bx^2+cx+d$, $y'=3ax^2+2bx+c$.\\ 
		Gọi $x_1$, $x_2$ lần lượt là điểm cực  đại, cực tiểu của hàm số $\Rightarrow y_1$, $y_2$ lần lượt là giá trị cực đại, cực tiểu của hàm số.
		\begin{itemchoice}
			\itemch Hàm số có hai  điểm cực trị cùng dấu.
			\itemch  Từ đồ thị ta thấy $y_1>0$, $y_2<0$ mà $-y_1>y_2\Rightarrow y_1+y_2<0$.
			\itemch Từ hình vẽ ta thấy đồ thị có hai điểm  cực trị.\\
			Do đó phương trình $y' = 3ax^2 + 2bx + c = 0$ có hai nghiệm phân biệt $x_1$, $x_2$. 
			\itemch Dựa vào đồ  thị ta thấy
			\begin{itemize}
				\item Nhánh bên phải  hướng lên $\Rightarrow a>0$.
				\item Đồ thị hàm số đi qua  gốc tọa độ $O(0;0)\Rightarrow d=0$.\\
				$\heva{&x_1+x_2=\dfrac{-2b}{3a}>0\\&x_1\cdot x_2=\dfrac{c}{3a}>0}\Rightarrow\heva{& \dfrac{b}{a}<0\\&c>0}\Rightarrow\heva{& b<0\\&c>0.}$\\ 
				Do đó $a > 0$, $b < 0$, $c > 0$, $d = 0$ $\Rightarrow$  Trong các hệ số $a$, $b$, $c$, $d$ có $2$ hệ số dương.
			\end{itemize}
		\end{itemchoice}
	 }
\end{ex}

% nguồn: Phan1_Lop_12_Dai_so_Chuong_1
\begin{ex}%[2D1V5-1]%[ma:7e298f30]
	Cho $a$, $b$, $c\in  \mathbb{R}$, $a\ne 0$, hàm số $f(x)=ax+b+\dfrac{c}{x+d}$ có đồ thị như hình bên.
	\begin{center}
		\resizebox{!}{5cm}{\begin{tikzpicture}[line join=round,line cap=round,>=stealth]
		\draw[cyan!50,thin](-3.5,-3.5) grid  (4.5,4.5);
		 \draw[->] (-3.5,0)--(4.5,0) node[below]{$x$};
		\draw[->]  (0,-3.5)--(0,4.5) node[left]{$y$};
		\fill (0,0) circle (1pt) node[above  left]{$O$};
		\fill (0,-1) circle  (1pt) node[left]{$-1$};
		\fill (0,3) circle (1pt) node[left]{$3$}; 
		 \fill (3,0) circle (1pt) node[below]{$3$};
		\fill (2,0) circle (1pt) node[below  right]{$2$};
		\fill (1,0) circle (1pt) node[below right]{$1$}; 
		\begin{scope}
			\clip (-3.5,-3.5)  rectangle (4.5,4.5);
			 \draw[domain=-3.5:1.1,smooth,samples=100] plot (\x,{(\x)^2-3*(\x)+3)/(\x-2)});
			 \draw[domain=1.15:4.5,smooth,samples=100] plot (\x,{(\x)^2-3*(\x)+3)/(\x-2)});
			 \draw[domain=-3.5:4.5,smooth] plot (\x,{(\x)-1});
		\end{scope}
		\draw  [dashed] (3,0)--(3,3)--(0,3);
	\end{tikzpicture}
		} 
	\end{center}
	 \choiceTF
		{\True  $d=-2$}
		{$a-\dfrac{c}{(x+d)^2}>0,\forall x \in  (1;2)\cup (2;3)$}
		{\True $a=1$, $b=-1$} 
		{$c=-2$} 
	\loigiai{ 
		\begin{itemchoice}
			\itemch Đúng.  Vì đồ thị hàm số có tiệm cận đứng $x=2$ nên $d=-2$.
			\itemch Sai. Vì $y'=a-\dfrac{c}{(x+d)^2}<0,\forall x  \in (1;2)\cup (2;3)$ do đồ thị đi xuống trên các khoảng $(1;2)$ và $(2;3)$.
			\itemch Đúng. Vì đồ thị hàm số có tiệm cận xiên $y=ax+b$ đi  qua điểm $A(1;0)$ và $B(0;-1)$.\\
			Suy ra $\heva{&a+b=0\\&b=-1} \Leftrightarrow \heva{&a=1\\&b=-1.}$ 
			\itemch Sai. Vì hàm số $f(x)=x-1+\dfrac{c}{x-2}$ có  đạo hàm $f'(x)=1-\dfrac{c}{(x-2)^2}$.\\
			Từ  đồ thị hàm số ta có $f'(1)=f'(3)=0 \Leftrightarrow 1-c=0 \Leftrightarrow c=1$.
		\end{itemchoice}
	} 
\end{ex}

% nguồn: Phan1_Lop_12_Dai_so_Chuong_1
\begin{ex}%[2D1V5-1]%[ma:2ee9f316]
	\immini{Cho hàm số $f(x)=x^3+ax^2+bx+c$, có đồ  thị như hình vẽ.
	 \choiceTF
		 {\True Hàm số $y=f(x)$ có hai điểm cực trị là $0$ và $2$}
		 {\True Giá trị $b$ bằng $0$}
		{Giá trị $c$  bằng $-2$}
		 {$f(x)=x^3-6x^2+2$}}
		{\resizebox{!}{3.5cm}{\begin{tikzpicture}[line cap=butt,line join=miter,>=stealth,scale=0.55]
		 \draw[->] (-2,0)--(4,0)
			node[shift={(-100:7pt)},font=\normalsize]{$x$}; 
		\draw[->]  (0,-3)--(0,3.5)
			node[shift={(180:7pt)},font=\normalsize]{$y$}; 
		\draw  (5pt,0) |- (0,5pt)
			(0,0) node[shift={(220:11pt)},font=\normalsize]{$O$}; 
		\tikzset{declare  function={f (\x)=(\x)^3-3*(\x)^2+2;}}
		\begin{scope}
			\clip (-2,-3)  rectangle (4,3.5);
			\draw[samples=100]  plot[domain=-2.5:4] (\x,{f (\x)});
		\end{scope}
		\draw[dashed,fill] 
			 (0,2) circle (1pt) node[shift=(150:0.25)]{$2$}
			(2,0) circle  (1pt) node[shift=(90:0.2)]{$2$}--
			(2,-2)  circle (1pt);
	\end{tikzpicture}
		 }}
	 \loigiaiDS{
		\begin{itemchoice}
			 \itemch {\bf Đúng}.\\
			 Từ đồ thị hàm số, suy ra hàm số $y=f(x)$ có hai điểm cực trị là $x_{\text{ CĐ}}=0$ và $x_{\text{ CT}}=2$.
			\itemch  {\bf Sai}.\\
			Ta có $f'(x)=3x^2+2ax+b$.\\ 
			Vì $x=0$ là điểm cực  đại của hàm số nên $f'(0)=0$. Suy ra $b=0$.
			 \itemch {\bf Sai}.\\
			Vì đồ thị hàm số đi qua điểm có tọa độ $(0;2)$ nên $f(0)=2$. Suy ra $c=2$. 
			\itemch {\bf Sai}.\\ 
			Ta có $\heva{&f'(0)=0 \\&f'(2)=0  \\&f(0)=2} \Leftrightarrow \heva{&a=-3\\&b=0\\&c=2.}$\\
			Vậy  hàm số cần tìm là $f(x)=x^3-3x^2+2$.
		\end{itemchoice}
	 }
\end{ex}

% nguồn: Phan1_Lop_12_Dai_so_Chuong_1
\begin{ex}%[2D1V5-1]%[ma:641a4539]
	Cho hàm số $y=\dfrac{x-1}{x+1}$ có đồ  thị $\left( C \right)$.
	 \choiceTF
		{\True Đạo hàm  hàm số là $y'=\dfrac{2}{\left( x+1 \right)^2}$}
		{\True Tiệm cận ngang của đồ thị hàm số  là $y=1$}
		{Tâm đối xứng của đồ thị hàm số  là $\left( 1;-1 \right)$}
		{Với mọi $M\in \left( C \right)$, tích khoảng cách từ $M$  đến các đường tiệm cận bằng $3$}
	 \loigiai{
		\begin{itemchoice}
			 \itemch {\bf Đúng.}\\
			Đạo hàm hàm  số là $y'=\dfrac{2}{\left( x+1 \right)^2}$.
			\itemch {\bf Đúng.}\\ 
			Tiệm cận ngang của đồ thị hàm số là $y=1$. 
			 \itemch {\bf Sai.}\\
			Tâm đối xứng của đồ thị hàm số  $\left( -1;1 \right)$.
			\itemch {\bf  Sai.}\\
			 Gọi $M\left( x_0;y_0 \right)\in \left( C \right)\Rightarrow M\left( x_0;1-\dfrac{2}{x_0+1} \right)$.\\
			 Khoảng cách từ $M$ đến tiệm cận đứng: $d_1=\left| x_0+1 \right|$.\\
			Khoảng cách từ $M$ đến tiệm cận ngang: $d_2=\left| y_0-1 \right|=\left| 1-\dfrac{2}{x_0+1}-1 \right|=\dfrac{2}{\left| x_0+1 \right|}$.\\ 
			Vậy $d_1\cdot  d_2=2$.
	\end{itemchoice}}
\end{ex}

% nguồn: Phan1_Lop_12_Dai_so_Chuong_1
\begin{ex}%[2D1V5-1]%[ma:e3163d6d]
	Cho hàm số $ y=f(x)=ax^3+bx^2+cx+d$ có đồ thị  như hình dưới đây:
	\begin{center}
		\begin{tikzpicture}[x=1cm,y=1cm,>=stealth]
			\def\aa{1} \def\bb{2} \def\cc{0} \def\dd{1} 
			\draw[->] (-2.46,0)--(1.5,0) node[below right]{$x$}; 
			\draw[->] (0,-1)--(0,3.5) node[left]{$y$}; 
			\fill (0,0) node[below left]{$O$}; 
			\draw[samples=150,smooth,domain=-2.36:0.9] plot (\x,{\aa*(\x)^3+\bb*(\x)^2+\cc*\x+\dd}); 
			\path  (-2,0) coordinate[label=below:$-2$](A)
				(-1,0)  coordinate[label=below:$-1$](B)
				(0,1) coordinate[label=right:$1$](C) 
				(0,2) coordinate[label=right:$2$](D) 
				(-2,1)  coordinate (E)
				(-1,2) coordinate (F); 
			\draw [dashed] (A)--(E)--(C) 
				 (B)--(F)--(D);
			\foreach \diem in  {A,B,C,D,E,F} \fill (\diem) circle (1pt);
		\end{tikzpicture}
	\end{center}
	\choiceTF 
		{\True Trong $4$  số $a$,$ b$, $c$, $d$ có ba số dương}
		 {Hàm số đồng biến trên khoảng $(-\infty;-1)$}
		{\True Đồ thị hàm số cắt $Oy$ tại điểm có toạ độ  $(0;1)$}
		{Hàm số  đạt cực tiểu tại $x=1$}
	 \loigiai{
		\begin{itemchoice}
			 \itemch \textbf{Đúng.}\\
			Ta  có$y'=f'(x)=3ax^2+2bx+c$.\\
			Hàm số đạt cực  tiểu tại $x=0$ nên $c=0$.\\
			Đồ thị đi qua điểm $(0;1)$ nên  $d=1$.\\
			Đồ thị đi qua các điểm $(-1;2)$; $(-2;1)$ nên $-a+b+1=2$ và $-8a+4b+1=1$.\\ 
			Ta  có hệ phương trình $\heva{&-a+b=1 \\ & -8a+4b=0 }$ $\Leftrightarrow \heva{&a=1 \\ & b=2 }$.\\
			 Vậy $ a=1,b=2,c=0,d=1$. Trong 4 số $ a,b,c,d$ có ba số dương.
			\itemch \textbf{Sai.}\\ 
			Trên khoảng $(-\infty;-1)$ hàm số vừa đồng biến vừa nghịch  biến.\\
			 \itemch \textbf{Đúng.}\\
			\itemch  \textbf{Sai.}\\
			Hàm số đạt cực tiểu  tại $x=0$.
		\end{itemchoice}
	 }
\end{ex}

% nguồn: Phan1_Lop_12_Dai_so_Chuong_1
\begin{ex}%[2D1V5-1]%[ma:55c34dbd]
	Cho hàm số $y = ax^3 + bx^2 + cx + d$,($a, b, c, d$ là~các số thực) có đồ thị  như hình vẽ bên dưới.
	\begin{center}
		\begin{tikzpicture}[line join=round,line cap=round,>=stealth,scale=0.7]
			\tikzset{ethicknode/.style={scale=0.9}} 
			\draw[->] (-2.1,0)--(6.1,0) node[below left]{$x$}; 
			\draw[->] (0,-1.1)--(0,7.1)  node[below left]{$y$};
			\draw (0,0) node  [below left]{$O$};
			\foreach  \x/\nx in {-1/-1,1/1,2/2,3/3,4/4,5/5}
			\draw[thin] (\x,1pt)--(\x,-1pt) node  [below]{$\nx$};
			\foreach  \y/\ny in {-1/-1,1/1,2/2,3/3,4/4,5/5}
			\draw[thin] (1pt,\y)--(-1pt,\y) node [left]{$\ny$}; 
			\draw[dashed,thin](1,0)--(1,5)--(0,5); 
			 \draw[dashed,thin](3,0)--(3,1)--(0,1);
			\begin{scope}
				\clip (-2,-1) rectangle (6,7); 
				 \draw[samples=100,domain=-1:5,smooth,variable=\x] plot (\x,{1*((\x)^3)+-6*((\x)^2)+9*(\x)+1});
			\end{scope}
		\end{tikzpicture}
	\end{center}
	\choiceTF 
		{Hàm số đã cho nghịch biến trên khoảng $\left(3;  +\infty\right)$}
		{\True  Tâm đối xứng của đồ thị hàm số đã cho là~điểm $I\left(2; 3\right)$}
		{Giá trị nhỏ nhất của  hàm số đã cho trên khoảng $\left(1; +\infty\right)$ bằng $3$}
		{$a, b, c$ là~các số  thực dương}
	\loigiai{ 
		\begin{itemchoice}
			\itemch Từ đồ thị ta thấy hàm  số đồng biến trên khoảng $\left(3; +\infty\right)$.
			\itemch Từ đồ thị ta thấy tâm đối xứng của đồ  thị hàm số là~điểm $I\left(2; 3\right)$.
			\itemch  Giá trị nhỏ nhất của hàm số trên khoảng $\left(1; +\infty\right)$ là~$1$.
			\itemch Từ  đồ thị ta thấy
			\begin{itemize}
				\item  $\lim\limits_{x\to +\infty}=+\infty$ nên $a > 0$.
				\item Do đồ thị đi qua điểm $\left(0;1\right)$ nên $  d=1>0$.
				 \item Ta có $y'=3ax^2+2bx+c$. \\
				Do hàm số có hai điểm cực trị $x_1$, $x_2$ thỏa  mãn
				$$\heva{&x_1+x_2>0\\&x_1 \cdot  x_2>0.}$$
				 Nên
				$$\heva{&-\dfrac{2b}{3a}>0\\&\dfrac{c}{3a}>0}\Rightarrow \heva{&b<0\\&c>0.}$$ 
			\end{itemize}
		\end{itemchoice}
	} 
\end{ex}


%%%%%%%%%%%%%%%% TRẢ LỜI NGẮN %%%%%%%%%%%%%%%%

% nguồn: Phan2_Lop_12_Dai_so_Chuong_1
\begin{bt}%[2D1N5-1]%[ma:f474fc24]
	\immini{ 
	Hình bên là đồ thị của  hàm số $y=\dfrac{ax+b}{x+c}$ (với $a,b,c\in\mathbb{R}$). Khi đó $ab-c$ bằng bao nhiêu?
	\shortans{-1} 
	}{ 
	\begin{tikzpicture}[>=stealth,scale=.8]
		\def\xmin{-2}\def\xmax{4}\def\ymin{-3}\def\ymax{1.5} 
		\draw[->](\xmin,0)--(\xmax,0) node  [shift={(-90:.2)},scale=1]{\scriptsize $x$};
		\draw[->](0,\ymin)--(0,\ymax)  node [below left] {\scriptsize $y$};
		\clip (\xmin,\ymin) rectangle (\xmax,\ymax); 
		 \foreach \p/\g in {1/135,2/45}\draw (\p,0) node[shift={(\g:.2)},scale=1]{\scriptsize $\p$}--+(0,.05)--+(0,-.05);
		 \foreach \p/\g in {-2/180,-1/135}\draw (0,\p) node[shift={(\g:.3)},scale=1]{\scriptsize $\p$}--+(.05,0)--+(-.05,0);
		\foreach \p/\g in {0/-135}\draw  (0,\p) node[shift={(\g:.3)},scale=1]{\scriptsize $O$};
		\draw[smooth,samples=100] plot[domain=\xmin:0.99](\x,{(-\x+2)/(\x-1)})  node[left]{\scriptsize};
		\draw[blue,smooth,samples=100]  plot[domain=1.01:\xmax](\x,{(-\x+2)/(\x-1)}) node[left]{\scriptsize};
		\draw[dashed,smooth,samples=100]  plot[domain=\xmin:\xmax](\x,{-1}) node[left]{\scriptsize};
		\draw[dashed,blue,smooth,samples=100](1,\ymin)--(1,\ymax); 
	\end{tikzpicture}
	 }
	\loigiai{ 
		 Đồ thị có tiệm cận đứng là đường thẳng $x=1$. Suy ra $-c=1 \Rightarrow c=-1$.\\
		 Đồ thị có tiệm cận ngang là đường thẳng $y=-1\Rightarrow a=-1$.\\
		 Mặt khác, $y(0)=-2\Leftrightarrow \dfrac bc=-2\Rightarrow b=2$.\\
		 Vậy, $ab-c=-2-(-1)=-1$.
	} 
\end{bt}

% nguồn: Phan2_Lop_12_Dai_so_Chuong_1
\begin{bt}%[2D1N5-1]%[ma:e7c6205d]
	 \immini{
	Cho hàm số $ y=\dfrac{ax+2}{cx+d}$ có đồ thị như hình vẽ. Tích $acd$ bằng bao  nhiêu?
	\shortans{-3} 
	 }{
	\begin{tikzpicture}[>=stealth,scale=.5]
		\def\xmin{-4}\def\xmax{10}\def\ymin{-3}\def\ymax{5} 
		\draw[->](\xmin,0)--(\xmax,0) node  [shift={(-90:.2)},scale=1]{\scriptsize $x$};
		 \draw[->](0,\ymin)--(0,\ymax) node [below right] {\scriptsize $y$};
		\clip (\xmin,\ymin) rectangle (\xmax,\ymax); 
		\foreach \p/\g/\a in  {0/-45/O,3/-45/3,-2/-90/-2}\draw (\p,0) node[shift={(\g:.2)},scale=1]{\scriptsize $\a$}--+(0,.05)--+(0,-.05);
		\foreach \p/\g in  {-1/180,1/120}\draw (0,\p) node[shift={(\g:.3)},scale=1]{\scriptsize $\p$}--+(.05,0)--+(-.05,0);
		 \draw[smooth,samples=100] plot[domain=\xmin:2.99](\x,{((\x)+2)/(\x-3)}) node[left]{\scriptsize};
		\draw[smooth,samples=100] plot[domain=3.01:\xmax](\x,{((\x)+2)/(\x-3)})  node[left]{\scriptsize};
		 \draw (\xmin,1)--(\xmax,1)
			 (3,\ymin)--(3,\ymax);
	\end{tikzpicture}
	} 
	\loigiai{ 
		Nhìn vào  đồ thị hàm số, ta thấy:
		\begin{itemize}
			\item Đồ thị hàm  số cắt $Ox$ tại điểm có hoành độ bằng $-2$, nên $\dfrac{-2}{a}=-2$, do đó $ a=1$.
			\item Đồ thị hàm số có đường tiệm cận ngang là  đường thẳng $ y=\dfrac{a}{c}=1$, mà $ a=1$, do đó $ c=1$.
			\item Đồ thị hàm số có đường tiệm cận đứng là đường thẳng $ x=-\dfrac{d}{c}=3$,  mà $ c=1$, do đó $ d=-3$.
		\end{itemize}
		Vậy $acd=1\cdot1\cdot(-3)=-3$. 
	} 
\end{bt}

% nguồn: Phan2_Lop_12_Dai_so_Chuong_1
\begin{bt}%[2D1H5-1]%[ma:d820318f]
	Cho hàm số $y=\dfrac{ax^2+bx+c}{mx+n}$ ($a>0, m\ne 0$) có đồ  thị như hình vẽ bên.
	\begin{center}
		\begin{tikzpicture}[line cap=butt,line join=miter,>=stealth,scale=1,line width=1pt,font=\tiny]
			\tikzset{declare  function={xmin=-2; xmax=5; ymin=-2; ymax=6;},
				 smooth,samples=450}
			\draw[->] (xmin,0)--(xmax,0) node[shift={(0:7pt)}]{$ x $}; 
			\draw[->] (0,ymin)--(0,ymax) node[shift={(90:7pt)}]{$ y $}; 
			\fill (0,0) node[shift={(45:5pt)}]{$ O  $};
			\clip (xmin,ymin-1) rectangle (xmax,ymax); 
			\draw[samples=250]  plot[domain=-2:0.8] (\x,{(\x)^2/((\x)-1)+1});
			\draw[samples=250] plot[domain=1.2:4] (\x,{(\x)^2/((\x)-1)+1}); 
			\draw  plot [domain=-2:4] (\x,{(\x)+2});
			\draw  (1,ymin)--(1,ymax);
			 \fill (0,2) node[left]{$2$}
				(0,3)  node[left]{$3$}
				(1,0) node[below  right]{$1$};
			\draw[dashed] (0,3)--(1,3); 
		\end{tikzpicture}
	\end{center}
	Hỏi trong các số $b,c,m,n$ có tất cả  bao nhiêu số dương?
	 \shortans{3}
	\loigiai{ 
		Ta có $y=\dfrac{ax^2+bx+c}{mx+n}$  ($a>0, m\ne 0$);\,\, $\mathbb{D}=\mathbb{R}\setminus\{-\dfrac{n}{m}\}$\\
		Tiệm cận xiên $y=x+2$.\\ 
		$\lim\limits_{x\to\infty}\dfrac{ax^2+bx+c}{mx^2+nx}=\dfrac{a}{m}=1\Rightarrow  a=m>0$.\\
		 $\lim\limits_{x\to\infty}[\dfrac{ax^2+bx+c}{mx+n}-x]=\lim\limits_{x\to\infty}\dfrac{(b+n)x+c}{mx+n}=\dfrac{b+n}{m}=2$.\\
		$\Rightarrow b+n=2m\Leftrightarrow  b-m=2m\Leftrightarrow b=3m=3a>0$.\\
	Tiệm cận đứng $x=-\dfrac{n}{m}=1\Leftrightarrow n=-m=-a<0$. Đồ  thị hàm số cắt trục tung tại $y_0=\dfrac{c}{n}>0\Rightarrow c<0$.}
\end{bt}

% nguồn: Phan2_Lop_12_Dai_so_Chuong_1
\begin{bt}%[2D1H5-1]%[ma:97ec07b8]
	\immini{Biết rằng hàm số $y=ax+b+\dfrac{c}{x+d}$ ($a$, $b$, $c$, $d\in \mathbb{R}$) có đồ thị như hình vẽ. Giá trị của $a+b+c+d$ là bao  nhiêu?}
	{\resizebox{!}{5cm}{\begin{tikzpicture}[line join=round,line cap=round,>=stealth]
		 \draw[->] (-3.5,0)--(4.5,0) node[below]{$x$};
		\draw[->] (0,-3.5)--(0,4.5) node[left]{$y$}; 
		\fill (0,0) circle (1pt)  node[above left]{$O$};
		\fill (0,-2) circle (1pt)  node[right]{$-2$};
		\fill (0,2) circle (1pt) node[left]{$2$}; 
		\fill (0,1) circle (1pt)  node[left]{$1$};
		\fill (1,0) circle (1pt)  node[below]{$1$};
		 \fill (-1,0) circle (1pt) node[above left]{$-1$};
		\begin{scope}
			\clip (-3.5,-3.5) rectangle (4.5,4.5); 
			\draw[domain=-3.5:1.1,smooth,samples=100] plot (\x,{(\x)^2+1)/(\x)}); 
			\draw[domain=1.15:4.5,smooth,samples=100] plot (\x,{(\x)^2+1)/(\x)}); 
			 \draw[domain=-3.5:4.5,smooth,dashed] plot (\x,{(\x)});
		\end{scope}
		\draw  [dashed] (1,0)--(1,2)--(0,2)
			 (1,1)--(0,1)
			(-1,0)--(-1,-2)--(0,-2); 
	\end{tikzpicture}
	 }
	} 
	 \shortans{2}
	 \loigiai{
		Đồ thị hàm số có tiệm  cận đứng $x=0$ nên $d=0$.\\
		Đồ thị hàm số  có tiệm cận xiên $y=x$ nên $a=1$, $b=0$.\\
		Đồ thị hàm  số $y=x+\dfrac{c}{x}$ đi qua điểm $A(1;2)$ nên $2=1+c \Leftrightarrow c=1$.\\
		Suy ra  $a+b+c+d=2$.
	 }
\end{bt}

% nguồn: Phan2_Lop_12_Dai_so_Chuong_1
\begin{bt}%[2D1H5-1]%[ma:5213597f]
	Cho hàm số $y=\dfrac{ax+2}{x+b}$ ($a$, $b \in  \mathbb{R}$) có bảng biến thiên như hình bên. Giá trị của $a+b$ là bao nhiêu?
	\begin{center}
		\begin{tikzpicture}
			\tkzTabInit[nocadre,lgt=1.3,espcl=3,deltacl=0.7] 
				 {$x$/1,$f’(x)$/1,$f (x)$/2}
				{$-\infty$,$-1$,$+\infty$} 
			 \tkzTabLine{,-,d,-,}
			 \tkzTabVar{+/$2$,-D+/$-\infty$/$+\infty$,-/$2$}
		\end{tikzpicture}
	\end{center}
	\shortans{3} 
	 \loigiai{
		Đồ thị hàm số có tiệm cận đứng $x=-1$ nên  $b=1$.\\
		 Đồ thị hàm số có tiệm cận ngang $y=2$ nên $a=2$.\\
		 Suy ra $a+b=3$.
	} 
\end{bt}

% nguồn: Phan2_Lop_12_Dai_so_Chuong_1
\begin{bt}%[2D1H5-1]%[ma:64ef6e1d]
	\immini{Biết rằng hàm số $y=\dfrac{ax+b}{x-1}$ có đồ thị như hình vẽ. Giá trị của  $a+b$ là bao nhiêu?
	\shortans{3}}{\resizebox{!}{5cm}{\begin{tikzpicture}[>=stealth]
		 \def\a{1}
		 \def\b{2}
		 \def\c{1}
		 \def\d{-1}
		 \draw[->] (-5,0) -- (5,0) node[below] {$x$};
		\draw[->] (0,-5) -- (0,5)  node[left] {$y$};
		\draw  (0,0) node[above left]{$O$};
		 \draw[dashed] (1,-5)--(1,5)
			(-5,1)--(5,1); 
		\clip (-5,-5)  rectangle (5,5);
		 \pgfmathsetmacro{\can}{-(\d)/(\c)}
		 \draw[samples=100,smooth,domain=-5:{\can-.1}] plot (\x,{(\a*\x+(\b))/(\c*\x+(\d))});
		\draw[samples=100,smooth,domain={\can+.1}:5]  plot (\x,{(\a*\x+(\b))/(\c*\x+(\d))});
		\draw (1,0)  node[below right]{$1$} circle (1pt);
		\draw  (0,1) node[above left]{$1$} circle (1pt);
		\draw (0,-2) node[below left]{$-2$} circle (1pt); 
		 \draw (-2,0) node[below]{$-2$} circle (1pt);
	\end{tikzpicture}
	 }
	} 
	 \loigiai{
		Đồ thị hàm số đi qua $C(0;-2)$ nên  $b=2$.\\
		Đồ thị hàm  số có tiệm cận đứng $y=1$ nên $a=1$.\\
		 Vậy $a+b=3$.
	} 
\end{bt}

% nguồn: Phan2_Lop_12_Dai_so_Chuong_1
\begin{bt}%[2D1H5-1]%[ma:70a62a83]
	Cho hàm  số $y=\dfrac{ax+2}{bx-2}$ với $a,b\in \mathbb{Z}$ có đồ thị như hình bên dưới. Tính giá trị của biểu thức $T=a+b$.
	\begin{center}
		\begin{tikzpicture}[line join=round,line cap=round,>=stealth]
			\tikzset{ethicknode/.style={scale=0.9}} 
			\draw[->]  (-2.6,0)--(6.1,0) node[below left] {$x$};
			\draw[->] (0,-2.1)--(0,6.1) node[below left]  {$y$};
			\draw (0,0) node  [below left] {$O$};
			\foreach  \x/\nx in {-2/-2,2/2}
			\draw[thin] (\x,1pt)--(\x,-1pt) node [below]  {$\nx$};
			\foreach \y/\ny  in {-1/-1,1/1}
			 \draw[thin] (1pt,\y)--(-1pt,\y) node [left] {$\ny$};
			 \draw[dashed,thin] (2.01,-2)--(2.01,6);
			\draw[dashed,thin]  (2.01,-2)--(2.01,6);
			\begin{scope}
				\clip (-2.5,-2)  rectangle (6,6);
				\draw[samples=200,domain=-2.5:1.99,smooth,variable=\x] plot  (\x,{(1*(\x)+2)/(1*(\x)+-2)});
				\draw[samples=200,domain=2.01:8,smooth,variable=\x] plot (\x,{(1*(\x)+2)/(1*(\x)+-2)}); 
				\draw[dashed,thin] (-2.5,1/1)--(8,1/1); 
			\end{scope}
		\end{tikzpicture}
	\end{center}
	 \shortans{2}
	\loigiai{ 
		 Đường tiệm cận đứng có phương trình là $x=2$ thì $\dfrac{2}{b}=2\Rightarrow b=1$.\\
		Đường  tiệm cận ngang có phương trình là $y=1$ thì $\dfrac{a}{1}=1\Rightarrow a=1$.\\
		Vậy $T=a+b=1+1=2$. 
	} 
\end{bt}

% nguồn: Phan2_Lop_12_Dai_so_Chuong_1
\begin{bt}%[2D1H5-1]%[ma:22e262e8]
	 \immini{Biết rằng hàm số $y=ax^3+bx^2+c$ ($a, b, c\in \mathbb{R}$) có đồ thị như hình vẽ. Giá trị của $a+b+c$ là bao nhiêu?
	\shortans{0}} 
	{\resizebox{!}{5cm}{\begin{tikzpicture}[line join=round,line cap=round,>=stealth]
		\draw[->]  (-2.5,0)--(4,0) node[below]{$x$};
		 \draw[->] (0,-3)--(0,3) node[left]{$y$};
		\begin{scope}
			\clip (-2.5,-3)  rectangle (3,4);
			\draw[smooth,domain=-3:3,samples=100] 
				plot (\x,{(\x)^3-3*(\x)^2+2}); 
		\end{scope}
		 \draw[dashed] (2,0)--(2,-2)--(-1,-2)--(-1,0);
		\fill (-1,0)  circle (1pt) node[above left] {$-1$}
			(2,0) circle (1pt)  node[above] {$2$};
		\fill (0,-2) circle (1pt) node[below right]  {$-2$}
			 (0,2) circle (1pt) node[above right] {$2$};
		\fill  (0,0) circle (1pt) node[below left]{$O$};
	\end{tikzpicture}
	} 
	} 
	 \loigiai{
		 Đồ thị hàm số đi qua $A(0;2)$ nên $c=2$.\\
		Đồ thị hàm số $y=ax^3+bx^2+2$ đi qua $A(-1;-2)$ và $B(2;-2)$ nên ta  có
		 \[\heva{&-a+b+2=-2\\&8a+4b+2=-2} \Leftrightarrow \heva{&a=1\\&b=-3.}\]
		Suy ra  $a+b+c=1-3+2=0$.
	 }
\end{bt}

% nguồn: Phan2_Lop_12_Dai_so_Chuong_1
\begin{bt}%[2D1H5-1]%[ma:da32c66d]
	Cho hàm số $y=ax^3+bx^2+cx+d$ ($a$, $b$, $c$, $d \in \mathbb{R}$) có bảng biến thiên như hình bên. Giá trị  của $a+b+c+d$ là bao nhiêu?
	\begin{center}
		\begin{tikzpicture}
			\tkzTabInit[nocadre,lgt=1.2,espcl=3,deltacl=0.7] 
				{$x$ /1,$f’(x)$ /1,$f (x)$  /2}
				 {$-\infty$,$-1$,$1$,$+\infty$}
			\tkzTabLine{,+,0,-,0,+,} 
			 \tkzTabVar{-/$-\infty$,+/$4$,-/$-4$,+/$+\infty$}
		\end{tikzpicture}
	\end{center}
	 \shortans{6,75}
	 \loigiai{
		Do đồ thị hàm số đi qua $A(1;-4)$ nên $a+b+c+d=-4$. 
	 }
\end{bt}

% nguồn: Phan2_Lop_12_Dai_so_Chuong_1
\begin{bt}%[2D1H5-1]%[ma:fb2e9eff]
	Đồ thị trong hình  bên dưới là của hàm số $y=ax+b+\dfrac{1}{x+c}$ ($a,b,c \in \mathbb{R}$).
	Khi đó tổng  $a+b+c$ bằng bao nhiêu?
	\begin{center}\begin{tikzpicture}[line cap=butt,line join=miter,>=stealth,scale=1,font=\tiny]
		\tikzset{declare function={xmin=-2;  xmax=5; ymin=-3; ymax=6;},
			smooth,samples=450} 
		\draw[->] (xmin,0)--(xmax,0)  node[shift={(0:7pt)}]{$ x $};
		\draw[->] (0,ymin)--(0,ymax)  node[shift={(90:7pt)}]{$ y $};
		\fill (0,0)  node[shift={(45:5pt)}]{$ O $};
		\clip (xmin,ymin-1)  rectangle (xmax,ymax);
		\draw[blue,samples=250] plot[domain=-2:0.8]  (\x,{(\x)^2/((\x)-1)});
		 \draw[blue,samples=250] plot[domain=1.2:4] (\x,{(\x)^2/((\x)-1)});
		\draw[red] plot [domain=-2:4] (\x,{(\x)+1}); 
		\draw[red!60]  (1,ymin)--(1,ymax);
		 \fill (0,1) node[left]{$1$}
			(0,2) node[left]{$2$} 
			 (1,0) node[below right]{$1$};
		\draw[dashed]  (0,2)--(1,2);
	\end{tikzpicture}
	\end{center}
	\shortans{3} 
	 \loigiai{
		Ta có $y=ax+b$ là đường tiệm cận xiên của đồ  thị hàm số. Từ đồ thị ta suy ra được $y=x+1$ là tiệm cận xiên nên $a=1$, $b=1$.\\
	$x=1$ là tiệm cận đứng của đồ thị hàm số  nên $c=-1$. Vậy $a+b+c=3$.}
\end{bt}

% nguồn: Phan2_Lop_12_Dai_so_Chuong_1
\begin{bt}%[2D1V5-1]%[ma:40428485]
	\immini{ 
	Cho hàm số $y=f(x)=ax^3+bx^2+cx+d$ có đồ thị $(C)$ như  hình vẽ. Khi đó $f(5)$ bằng
	\shortans{64} 
	}{ 
	\begin{tikzpicture}[>=stealth]
		 \draw[->,line width = 1pt] (-2,0)--(0,0) node[below left]{$O$}--(3,0) node[below]{$x$};
		\draw[->,line width = 1pt] (0,-2)  --(0,3) node[right]{$y$};
		\foreach \x in  {-1,1,2}{
			\draw (\x,0) node[below]{$\x$} circle  (1pt);
			\draw (0,\x) node[left]{$\x$} circle (1pt); 
		} 
		\draw  [domain=-0.2:2.3,samples=100]
			plot (\x,{(\x)^3-3*(\x)^2+3*(\x)-1}) node[right]{$(C)$}; 
	\end{tikzpicture}
	} 
	\loigiai{ 
		Gọi hàm số bậc $3$ cần tìm có dạng là $y=ax^3+bx^2+cx+d  \Rightarrow y' = 3ax^2+ 2bx + c$. Từ đồ thị ta có
		$$\heva{&d=-1\\&a+b+c+d=0\\&3a+2b+c=0\\&b^2-3ac=0}\Leftrightarrow\heva{&a=1\\&b=-3\\&c=3\\&d=-1}\Rightarrow y=x^3-3x^2+3x-1.$$ 
		Suy  ra $f(5)=64$.
	 }
\end{bt}

% nguồn: Phan2_Lop_12_Dai_so_Chuong_1
\begin{bt}%[2D1V5-1]%[ma:0d0b12cd]
	Có bao nhiêu giá trị nguyên dương của m sao cho đường thẳng $y=x+m$ cắt đồ thị hàm  số $y=\dfrac{2x-1}{x+1}$ tại hai điểm phân biệt $A$, $B$ và $AB\le 4$.
	 \shortans{1}
	 \loigiai{
		Xét phương trình hoành  độ giao điểm
		\begin{eqnarray*}
			& &\dfrac{2x-1}{x+1}=x+m\\ 
			 &\Leftrightarrow& \heva{&x\ne-1\\&\left(x+1\right)\left(x+m\right)=2x-1
			 }\\
			 &\Leftrightarrow& \heva{&x\ne-1\\&x^2+\left(m-1\right)x+m+1=0.\quad (*)
			 }
		\end{eqnarray*}
		Đường thẳng $y=x+m$ cắt đồ thị hàm số $y=\dfrac{2x-1}{x+1}$ tại hai điểm phân biệt $A$, $B$ khi và chỉ khi phương trình (*) có $2$ nghiệm phân biệt khác $-1$.  Ta có
		\begin{eqnarray*}
			&\Leftrightarrow&  \heva{&\Delta > 0\\ &\left(-1\right)^2+\left(m-1\right).\left(-1\right)+m+1\ne 0
			 }\\
			&\Leftrightarrow&  \heva{&\left(m-1\right)^2-4\left(m+1\right)>0\\
			&3\ne  0}\\
			&\Leftrightarrow& m^2-6m-3>0\\ 
			&\Leftrightarrow& \hoac{&m >  3+2\sqrt{3} \\
			&m  < 3-2\sqrt{3}.}
		\end{eqnarray*}
		Giao điểm của $2$ đồ thị hàm số lần lượt là  $A\left(x_1;x_1+m\right)$ và $B\left(x_2;x_2+m\right)$, trong đó $x_1, x_2$ là $2$ nghiệm phân biệt của phương trình (*).\\
		Theo định lý Vi-et, ta  có $\heva{&x_1+x_2=1-m\\&x_1x_2=m+1}$.\\
		 Khi đó
		\begin{eqnarray*}
			 &&AB=\sqrt{2\left(x_1-x_2\right)^2} \le 4\\
			&\Leftrightarrow& \left(x_1-x_2\right)^2\le 8\\ 
			 &\Leftrightarrow& \left(x_1+x_2\right)^2-4x_1x_2\le 8\\
			 &\Leftrightarrow& \left(1-m\right)^2-4\left(m+1\right)\le 8\\
			&\Leftrightarrow& m^2-6m-11\le  0\\
			&\Leftrightarrow& 3-2\sqrt{5}  \le m\le 3+2\sqrt{5}.
		\end{eqnarray*}
		Kết  hợp với điều kiện có nghiệm $\Leftrightarrow \hoac{&m > 3+2\sqrt{3}\\&m < 3-2\sqrt{3}}\Rightarrow \hoac{&3+2\sqrt{3} < m\le 3+2\sqrt{5} \\&3-2\sqrt{5} \le m < 3-2\sqrt{3}.}$\\
		 Suy ra, một giá trị nguyên dương của $m$ là $7$.
	 }
\end{bt}

% nguồn: Phan2_Lop_12_Dai_so_Chuong_1
\begin{bt}%[2D1V5-1]%[ma:7dd89546]
	Cho hàm số $y=\dfrac{\left(2m-1\right)x-m}{x+m}$ ($m\ne 0$) có đồ thị $\left(C_m \right)$. Biết rằng tồn  tại duy nhất một đường thẳng $(d)$ có phương trình $y=ax+b$ sao cho $\left(C_m \right)$ luôn tiếp xúc với $(d)$. Tính giá trị của $a+b$.
	 \shortans{1}
	 \loigiai{
		 Điều kiện: $x\neq-m$.\\
		Phương  trình hoành độ giao điểm của $\left(C\right)$ và $(d)$ như sau
		\[\dfrac{\left(2m-1\right)x-m}{x+m}=ax+b \Leftrightarrow ax^2+\left(am+b-2m+1\right)x+bm+m=0\,  (1).\]
		Do $(d)$ luôn tiếp xúc với $\left(C\right)$ nên phương trình $(1)$  luôn có nghiệm kép khác $-m$ với mọi $m\neq0$.\\
		Suy ra $a\ne 0$ và 
		\begin{align*}
			\Delta'&=\left(am+b-2m+1\right)^2-4am\left(b+1\right)=0,\forall  m\ne 0\\
			 &\Leftrightarrow \left(a-2\right)^2m^2+\left[2\left(a-2\right)\left(b+1\right)-4a\left(b+1\right)\right]m+\left(b+1\right)^2=0,\forall m\ne 0\\
			 &\Leftrightarrow\heva{&\left(a-2\right)^2=0\\
			&2\left(a-2\right)\left(b+1\right)-4a\left(b+1\right)=0\\ 
			&\left(b+1\right)^2=0 
			} \Leftrightarrow \heva{&a=2 \\  &b=-1.}
		\end{align*}
		 Vậy $a+b=1$.
	 }
\end{bt}

% nguồn: Phan2_Lop_12_Dai_so_Chuong_1
\begin{bt}%[2D1V5-1]%[ma:bca03100]
	\immini{Cho hàm số $y=ax^3+bx^2+cx+d \, (a,b,c,d \in \mathbb{R})$ có đồ thị  là đường cong trong hình bên. Có bao~nhiêu số dương trong các số $a,b,c,d$?}
	 {
	\begin{tikzpicture}[line join=round,line cap=round,>=stealth,scale=0.6]
		\tikzset{ethicknode/.style={scale=0.9}} 
		\draw[->] (-4.1,0)--(2.1,0) node[below left]  {$x$};
		 \draw[->] (0,-2.1)--(0,5.1) node[below left] {$y$};
		\draw (0,0)  node [below left] {$O$};
		\begin{scope}
			\clip  (-4,-2) rectangle (2,5);
			\draw[samples=100,domain=-4:2,smooth,variable=\x]  plot (\x,{-1*((\x+1.5)^3)+0*((\x+1.5)^2)+3*(\x+1.5)+2.5});
		\end{scope}
	\end{tikzpicture}
	 }
	\shortans{1} 
	 \loigiai{
		 \textbf{Phương pháp:}\\
		Từ đồ thị xác định dấu của hệ số $a$,  $b$, $c$, $d$.\\
		 \textbf{Cách giải:}\\
		Ta  có $y'=3ax^2+2bx+c$.\\
		Dựa vào đồ thị  ta thấy $a<0$.\\
		Hàm số có $2$ cực trị âm nên $$\heva{&\Delta_{y'} > 0 \\& S < 0 \\& P > 0 } \Leftrightarrow \heva{& (2b)^2-4(3a)c > 0 \\& -\dfrac{2b}{3a} < 0 \\& \dfrac{c}{3a} > 0 } \Rightarrow \heva{& 4b^2-12ac > 0 \\&  \dfrac{b}{a} > 0 \\& \dfrac{c}{a} < 0 } \Rightarrow \heva{& b < 0 \\&c > 0. }$$
		Đồ thị cắt trục  $Oy$ tại điểm $(0;d)$ nên $d>0$.\\
		Vậy có đúng $1$ số dương trong các số $a$, $b$, $c$,  $d$.
	 }
\end{bt}

% nguồn: Phan2_Lop_12_Dai_so_Chuong_1
\begin{bt}%[2D1V5-1]%[ma:863cf2f4]
	\immini{ 
	Cho hàm số bậc ba $y=ax^3+bx^2+cx+d$ có đồ thị như hình vẽ. Tìm giá trị nhỏ nhất của biểu thức $P=a^2+c^2+b+1$. (Kết quả làm tròn đến  hàng phần trăm).
	}{ 
	\begin{tikzpicture}[>=stealth,font=\footnotesize,xscale=.8,scale=.8]
		\draw[->] (-1.5,0)--(3,0)  node[below]{$x$};
		\draw[->] (0,-1)--(0,3) node[left]{$y$}; 
		\draw  plot[smooth,samples=100,domain=-.25:2.25] (\x,{(\x-1)^3+1})
			(1-.5,1)--(1+.5,1); 
		\fill (1,1)  circle (1.5pt)
			(0,0)  node[below left]{$O$}circle (1.5pt);
	\end{tikzpicture}
	 }
	\loigiai{ 
		 Ta có $y'=3ax^2+2bx+c$. Từ đồ thị suy ra $d=0$, $b^2-3ac=0\Rightarrow ac=\dfrac{b^2}{3}$.\\
		Do đó, $P\ge 2ac+b+1=\dfrac{2b^2}{3}+b+1=\dfrac{2}{3}\left(b+\dfrac{3}{4}\right)^2+\dfrac{5}{8}\ge  \dfrac{5}{8}$.\\
		Vậy $\min P=\dfrac{5}{8}\approx 0{,}63$ tương  ứng với $a=c=\dfrac{\sqrt{3}}{4},b=-\dfrac{3}{4}$.
	 \shortans{0,63}
	} 
\end{bt}

% nguồn: Phan2_Lop_12_Dai_so_Chuong_1
\begin{bt}%[2D1V5-1]%[ma:54db2a5f]
	Tập tất cả các giá trị thực của tham số $m$ để đường thẳng $d\colon y=x-2m$ cắt đồ thị hàm số $y=\dfrac{x-3}{x+1}$  tại hai điểm phân biệt có hoành độ dương là $\left(a; b\right)$. Tính giá trị biểu thức $a+b$.
	\shortans{2,5} 
	\loigiai{ 
		Điều kiện:  $x\neq-1$.\\
		Xét phương trình hoành độ giao  điểm
		\begin{eqnarray*}
			 & &\dfrac{x-3}{x+1}=x-2m\\
			 &\Leftrightarrow& x-3=\left(x+1\right)\left(x-2m\right)\\
			&\Leftrightarrow&  x-3=x^2-2mx+x-2m\\
			&\Leftrightarrow& x^2-2mx+3-2m=0\quad  (1).
		\end{eqnarray*}
		Đường  thẳng $d\colon y=x-2m$ cắt đồ thị hàm số $y=\dfrac{x-3}{x+1}$ tại hai điểm phân biệt có hoành độ dương $\Leftrightarrow$ (1) có hai nghiệm dương phân biệt.
		\[\Leftrightarrow \heva{&\Delta '=m^2-\left(3-2m\right)>0\\&S=2m>0\\&P=3-2m>0}\Leftrightarrow \heva{&m^2+2m-3> 0\\&m > 0\\&m < \dfrac{3}{2}} \heva{&\hoac{&m <-3\\&m > 1}\\&m > 0\\&m < \dfrac{3}{2}} \Leftrightarrow 1< m < \dfrac{3}{2}.\] 
		Vậy  với $1< m < \dfrac{3}{2}$ thì đường thẳng $d\colon y=x-2m$ cắt đồ thị hàm số $y=\dfrac{x-3}{x+1}$ tại hai điểm phân biệt có hoành độ dương nên $a+b=\dfrac{5}{2}=2{,}5$.
	} 
\end{bt}

% nguồn: Phan2_Lop_12_Dai_so_Chuong_1
\begin{bt}%[2D1V5-1]%[ma:94958a81]
	Cho hàm số $y=\dfrac{x+2}{x+1} \left(C\right)$ và đường thẳng $(d)\colon y=x+m$. Có bao nhiêu giá trị nguyên $m$ thuộc khoảng $\left(-10;10\right)$ để đường thẳng  $(d)$ cắt đồ thị $\left(C\right)$ tại hai điểm về hai phía của trục hoành?
	 \shortans{11}
	 \loigiai{
		Phương trình hoành độ giao điểm của $(C)$ và $(d)$ là 
		\begin{eqnarray*}
			& &\dfrac{x+2}{x+1}=x+m\quad (1)\\ 
			(1)&\Leftrightarrow& \heva{&x\ne-1\\&f(x)=x^2+mx+m-2=0\quad  (2).}
		\end{eqnarray*}
		Đường thẳng $(d)$ cắt $\left(C\right)$ tại hai  điểm phân biệt $\Leftrightarrow$ Phương trình (2) có hai nghiệm phân biệt khác $-1$
		$\Leftrightarrow \heva{&\Delta=m^2-4m+8>0\\&f\left(-1\right)=1-m+m-2\ne 0}\Leftrightarrow \forall  m\in \mathbb{R}$.\\
		Gọi $x_1; x_2$ là hai  nghiệm phương trình (2).\\
		Theo định lý Vi-et ta có  $\heva{&x_1+x_2=-m\\&x_1\cdot x_2=m-2.}$\\
		Đường thẳng $(d)$ cắt đồ thị $\left(C\right)$ tại hai  điểm về hai phía của trục hoành
		\begin{eqnarray*}
			&\Leftrightarrow& y_1\cdot  y_2< 0\\
			&\Leftrightarrow& \left(x_1+m\right)\left(x_2+m\right)  < 0\\
			 &\Leftrightarrow& x_1\cdot x_2+m\left(x_1+x_2\right)+m^2< 0\\
			&\Leftrightarrow& m-2-m^2+m^2< 0\Leftrightarrow m  < 2.
		\end{eqnarray*}
		Do $m\in \mathbb{Z}$  và $m\in \left(-10;10\right)\Rightarrow m\in \left\{-9;-8;\ldots;-1;0;1\right\}$.\\
		Vậy có $11$ giá  trị $m$.
	 }
\end{bt}

% nguồn: Phan2_Lop_12_Dai_so_Chuong_1
\begin{bt}%[2D1V5-1]%[ma:076ee1e1]
	Cho hàm số $y=\dfrac{x+m}{x-1}$ có đồ thị là đường cong $\left(H\right)$ và đường thẳng $\Delta$ có phương trình  $y=x+1$. Số giá trị nguyên của tham số $m$ nhỏ hơn $10$ để đường thẳng $\Delta$ cắt đường cong $\left(H\right)$ tại hai điểm phân biệt nằm về hai nhánh của đồ thị bằng bao nhiêu?
	 \shortans{10}
	 \loigiai{
		Phương trình hoành độ giao điểm $\dfrac{x+m}{x-1}=x+1\Leftrightarrow \heva{&x\ne 1\\  &f(x)=x^2-x-m-1=0\quad (*).}$\\
		Đường thẳng  $\Delta$ cắt đường cong $\left(H\right)$ tại hai điểm phân biệt khi
		\[\heva{&f(1)\ne 0\\&\Delta  _{\left(*\right)} > 0
		} \Leftrightarrow 4m+5>  0\Leftrightarrow m >-\dfrac{5}{4}.
		\] 
		Gọi $x_1, x_2$ là hai nghiệm của phương trình (*), ta  có $\heva{&x_1+x_2=1\\ &x_1x_2=-m-1.}$\\
		Do $\Delta$ cắt $(H)$ tại hai điểm phân biệt nằm  về hai nhánh của đồ thị nên ta có
		\begin{eqnarray*}
			&&\left(x_1-1\right)\left(x_2-1\right)  < 0\\
			&\Leftrightarrow& x_1  x_2-\left(x_1+x_2 \right)+1< 0\\
			&\Leftrightarrow&-m-1< 0\Leftrightarrow m >-1. 
		\end{eqnarray*}
		Vì $m$ nguyên và  $m<10$ nên ta có $m\in \left\{0;1;2;3;4;5;6;7;8;9\right\}$.\\
		Vậy có $10$ giá  trị nguyên của tham số $m$ thỏa mãn điều kiện đề bài.
	} 
\end{bt}


%%%%%%%%%%%%%%%% TỰ LUẬN %%%%%%%%%%%%%%%%

% nguồn: Phan2_Lop_12_Dai_so_Chuong_1
\begin{bt}%[2D1N5-1]%[ma:32c2eb93]
	 \immini{
	Cho hàm  số $y=ax^3+bx^2+cx+d$ có đồ thị như hình bên. Xác định dấu của các hệ số $a$, $b$, $c$ và $d$.
	}{ 
	\begin{tikzpicture}[scale=1,font=\footnotesize,line join=round,line cap=round,>=stealth]
		 \def\a{0.5}
		\def\xmin{-1.3}  \def\xmax{4}
		\def\ymin{-2.5} \def\ymax{2} 
		 \draw[->] (\xmin,0)--(\xmax,0) node [below]{$x$};
		\draw[->] (0,\ymin)--(0,\ymax)  node [left]{$y$};
		\node at  (0,0) [below left]{$O$};
		 \clip (\xmin+0.1,\ymin+0.1) rectangle (\xmax-0.1,\ymax-0.1);
		 \draw[smooth,samples=300] plot (\x,{\a*((\x+0.8)*(\x-0.6)*(\x-3))});
	\end{tikzpicture}
	} 
	 \loigiai{
		Từ dáng điệu của đồ thị ta có $\lim\limits_{x\to+\infty} y=+\infty;\lim\limits_{x\to-\infty} y=-\infty\Rightarrow  a>0$.\\
		Đồ thị hàm số cắt trục tung tại một điểm có tung độ dương nên  $d>0$.\\
		Ta  có $y'=3ax^2+2bx+c$.\\
		Mặt khác dựa vào đồ thị ta  thấy phương trình $y'=0$ có hai nghiệm trái dấu và tổng hai nghiệm này luôn dương nên $\heva{&ac<0\\&-\dfrac{2b}{3a}>0}\Rightarrow \heva{&c<0\\&b<0.}$\\
		Vậy  $a>0$, $b<0$, $c<0$ và $d>0$.
	} 
\end{bt}

% nguồn: Phan2_Lop_12_Dai_so_Chuong_1
\begin{bt}%[2D1H5-1]%[ma:392f9e4b]
	Khảo sát sự biến thiên và vẽ đồ thị  các hàm số sau
	\begin{enumerate}
		\item  $y=x^3-3x^2+1$.
		\item $y=-2x^3-3x^2+1$. 
		\item  $y=x^3+3x^2+3x+2$.
		 \item $y=x^3-3x^2+4x-2$.
	\end{enumerate}
	 \loigiai{
		\begin{enumerate}
			\item  $y=x^3-3x^2+1$.\\
			 Tập xác định $\mathscr{D}=\mathbb{R}$.\\
			Sự  biến thiên:\\
			 $y'=3x^2-6x$.\\
			$y' = 0 \Leftrightarrow\hoac{&x=0\\&x=2}$ và $\lim\limits_{x  \to-\infty}y=-\infty$; $\lim\limits_{x \to+\infty}y=+\infty$.\\
			Bảng biến thiên 
			\begin{center}
				\begin{tikzpicture}[scale=1,font=\footnotesize,line join=round,line cap=round,>=stealth]
					\tkzTabInit[nocadre=false,lgt=1.2,espcl=2.5,deltacl=0.6] 
						 {$x$/0.6,$y'$/0.6,$y$/2}
						 {$-\infty$,$0$,$2$,$+\infty$}
					\tkzTabLine{,+,0,-,0,+} 
					 \path
						($(N13)!0.1!(N12)$) node  (A1){$-\infty$}
						($(N23)!0.8!(N22)$) node  (A2){$1$}
						($(N33)!0.1!(N32)$) node (A3){$-3  $}
						($(N43)!0.8!(N42)$)  node (A4){$+\infty$};
					\foreach \x/\y in  {A1/A2,A2/A3,A3/A4}{
						\draw[-stealth]  (\x)--(\y);
					 }
				\end{tikzpicture}
			\end{center}
		 \immini{
			Từ bảng biến thiên ta thấy hàm số nghịch biến trên khoảng $(0;  2)$ và đồng biến trên mỗi khoảng $(-\infty; 0)$ và $(2;+\infty)$.\\
			Hàm số đạt cực đại tại $x=0$; $y_{\text{ CĐ}}=1$ và đạt cực tiểu tại  $x=2$; $y_{\text{ CT}}=-3$.\\
			Đồ thị hàm số đi qua các điểm $(2;-3)$;  $(-1;-3)$; $(3; 1)$.\\
			Đồ thị nhận điểm  $I(1;-1)$ làm tâm đối xứng.
			 }
			 {
			\begin{tikzpicture}[line cap=butt,line join=miter,>=stealth,scale=0.7]
				\tikzset{declare  function={xmin=-1.5; xmax=4; ymin=-4; ymax=2;
					 f (\x)=1*(\x)^3-3*(\x)^2+1;
					}, 
					smooth,samples=450 
				} 
				\draw[->] (xmin,0)--(xmax,0) node[shift={(-100:7pt)},font=\normalsize]{$ x $}; 
				\draw[->]  (0,ymin)--(0,ymax) node[shift={(190:7pt)},font=\normalsize]{$ y $};
				\fill (0,0)  node[shift={(225:7pt)},font=\normalsize]{$ O $};
				\draw[dashed] (-1,0) node[above]{$-1$}--(-1,-3)--(0,-3) node[below left]{$-3$}--(2,-3)--(2,0)  node[above]{$2$}
					(1,0) node[above]{$1$}--(1,-1)--(0,-1) 
					(0,-1)  node[left,xshift=4pt] {$-1$}
					(3,0) node[below]{$3$}--(3,1)--(0,1) node[above  left]{$1$};
				\clip  (xmin,ymin) rectangle (xmax,ymax);
				\draw  plot[domain=xmin:xmax] (\x,{f (\x)});
			\end{tikzpicture}
			} 
			\item  $y=-2x^3-3x^2+1$.\\
			 Tập xác định $\mathscr{D}=\mathbb{R}$.\\
			 Sự biến thiên:\\
			$y'=-6x^2-6x$.\\ 
			 $y'=0\Leftrightarrow\hoac{&x=0\\&x=-1}$ và $\lim\limits_{x \to-\infty}y=-\infty$; $\lim\limits_{x \to+\infty}y=+\infty$.\\
			Bảng  biến thiên
			\begin{center}
				\begin{tikzpicture}[scale=1,font=\footnotesize,line join=round,line cap=round,>=stealth]
					\tkzTabInit[nocadre=false,lgt=1.2,espcl=2.5,deltacl=0.6] 
						 {$x$/0.6,$y'$/0.6,$y$/2}
						{$-\infty$,$-1$,$0$,$+\infty$} 
					 \tkzTabLine{,-,0,+,0,-}
					\path 
						($(N13)!0.8!(N12)$) node  (A1){$+\infty$}
						($(N23)!0.1!(N22)$) node  (A2){$0$}
						($(N33)!0.8!(N32)$) node  (A3){$1$}
						 ($(N43)!0.1!(N42)$) node (A4){$-\infty$};
					\foreach \x/\y  in {A1/A2,A2/A3,A3/A4}{
						\draw[-stealth] (\x)--(\y); 
					} 
				\end{tikzpicture}
			\end{center}
		\immini{ 
			Từ bảng biến thiên ta  thấy hàm số đồng biến trên khoảng $(-1; 0)$ và nghịch biến trên mỗi khoảng ($-\infty;-1$) và ($0;+\infty$).\\
			Hàm số đạt cực đại tại $x=0$; $y_{\text{ CĐ}}=1$ và đạt cực tiểu tại $x=-1$; $y_{\text{ CT}}=0$.\\ 
			Đồ thị hàm số đi  qua các điểm $(1;-4)$; $(-2; 5)$.\\
			Đồ  thị nhận điểm $I\left(\dfrac{1}{2};-\dfrac{1}{2}\right)$ làm tâm đối xứng.
			} 
			{ 
			\begin{tikzpicture}[line cap=butt,line join=miter,>=stealth]
				 \tikzset{declare function={xmin=-2; xmax=2; ymin=-1; ymax=2;
					 f (\x)=-2*(\x)^3-3*(\x)^2+1;
					 },
					smooth,samples=450 
				} 
				 \draw[->] (xmin,0)--(xmax,0) node[shift={(-100:7pt)},font=\normalsize]{$ x $};
				\draw[->]  (0,ymin)--(0,ymax) node[shift={(190:7pt)},font=\normalsize]{$ y $};
				\fill (0,0) node[shift={(225:7pt)},font=\normalsize]{$ O  $};
				\draw[dashed]  (-1,0) node[below]{$-1$}
					(0,1) node[above left]{$1$} 
					 (1,0) node[above]{$1$};
				\clip (xmin,ymin) rectangle  (xmax,ymax);
				\draw plot[domain=xmin:xmax] (\x,{f  (\x)});
			\end{tikzpicture}
			 }
			 \item $y=x^3+3x^2+3x+2$.\\
			Tập  xác định $\mathscr{D}=\mathbb{R}$.\\
			Sự  biến thiên:\\
			 $y'=3x^2+6x+3$.\\
			$y'=0\Leftrightarrow x=-1$ và $\lim\limits_{x \to-\infty}y=-\infty$; $\lim\limits_{x  \to+\infty}y=+\infty$.\\
			Bảng  biến thiên:
			\begin{center}
				\begin{tikzpicture}[scale=1,font=\footnotesize,line join=round,line cap=round,>=stealth]
					\tkzTabInit[nocadre=false,lgt=1.2,espcl=2.5,deltacl=0.6] 
						{$x$/0.6,$y'$/0.6,$y$/2} 
						 {$-\infty$,$-1$,$+\infty$}
					 \tkzTabLine{,+,0,+}
					 \path
						($(N13)!0.1!(N12)$)  node (A1){$-\infty$}
						 ($(N23)!0.45!(N22)$) node (A2){$1$}
						($(N33)!0.8!(N32)$) node (A3){$+\infty$}; 
					\foreach \x/\y  in {A1/A2,A2/A3}{
						\draw[-stealth]  (\x)--(\y);
					 }
				\end{tikzpicture}
			\end{center}
		 \immini{
			Từ  bảng biến thiên ta thấy hàm số đồng biến trên $\mathbb{R}$ và hàm số không có cực trị.\\
			Đồ thị nhận điểm $I(-1; 1)$ làm tâm đối xứng. 
			} 
			 {
			\begin{tikzpicture}[line cap=butt,line join=miter,>=stealth]
				\tikzset{declare function={xmin=-2.5; xmax=2; ymin=-2;  ymax=4;
					 f (\x)=1*(\x)^3+3*(\x)^2+3*(\x)+2;
					}, 
					smooth,samples=450 
				 }
				\draw[->] (xmin,0)--(xmax,0) node[shift={(-100:7pt)},font=\normalsize]{$ x  $};
				\draw[->] (0,ymin)--(0,ymax)  node[shift={(190:7pt)},font=\normalsize]{$ y $};
				\fill  (0,0) node[shift={(225:7pt)},font=\normalsize]{$ O $};
				\draw[dashed]  (-1,0) node[below]{$-1$}--(-1,1)--(0,1) node[right]{$1$}
					 (1,0) node[below]{$1$};
				\clip (xmin,ymin)  rectangle (xmax,ymax);
				 \draw plot[domain=xmin:xmax] (\x,{f (\x)});
			\end{tikzpicture}
			 }
			\item  $y=x^3-3x^2+4x-2$.\\
			Tập xác định  $\mathscr{D}=\mathbb{R}$.\\
			Sự  biến thiên: \\
			$y'=3x^2-6x+4> 0\forall x \in \mathbb{R}$ và $\lim\limits_{x \to-\infty}y=-\infty$;  $\lim\limits_{x \to+\infty}y=+\infty$.\\
			Bảng  biến thiên:
			\begin{center}
				\begin{tikzpicture}[scale=1,font=\footnotesize,line join=round,line cap=round,>=stealth]
					 \tkzTabInit[nocadre=false,lgt=1.2,espcl=2.5,deltacl=0.6]
						{$x$/0.6,$y'$/0.6,$y$/2} 
						 {$-\infty$,$+\infty$}
					 \tkzTabLine{,+}
					\path 
						($(N13)!0.1!(N12)$) node  (A1){$-\infty$}
						 ($(N23)!0.8!(N22)$) node (A2){$+\infty$};
					\foreach \x/\y in {A1/A2}{ 
						\draw[-stealth] (\x)--(\y); 
					 }
				\end{tikzpicture}
			\end{center}
		\immini{ 
			Từ bảng biến thiên ta thấy hàm số đồng biến trên $\mathbb{R}$  và hàm số không có cực trị.\\
			Đồ thị hàm số đi  qua các điểm $(2; 2)$; $(0;-2)$; $(1; 0)$.\\
			Đồ thị nhận điểm $I(1; 0)$ làm  tâm đối xứng.
			}{ 
			\begin{tikzpicture}[line cap=butt,line join=miter,>=stealth]
				\tikzset{declare  function={xmin=-1.5; xmax=3.5; ymin=-2.5; ymax=2.5;
					f (\x)=1*(\x)^3-3*(\x)^2+4*(\x)-2; 
					}, 
					 smooth,samples=450
				} 
				\draw[->] (xmin,0)--(xmax,0) node[shift={(-100:7pt)},font=\normalsize]{$  x $};
				\draw[->] (0,ymin)--(0,ymax) node[shift={(190:7pt)},font=\normalsize]{$ y $}; 
				\fill (0,0) node[shift={(225:7pt)},font=\normalsize]{$ O $}; 
				\draw[dashed] (2,0) node[below]{$2$}--(2,2)--(0,2)  node[left]{$2$}
					(1,0)  node[above]{$1$}
					(0,-2) node[left]{$-2$}; 
				\clip (xmin,ymin) rectangle  (xmax,ymax);
				\draw plot[domain=xmin:xmax] (\x,{f (\x)}); 
			\end{tikzpicture}
			 }
		\end{enumerate}
	 }
\end{bt}

% nguồn: Phan2_Lop_12_Dai_so_Chuong_1
\begin{bt}%[2D1H5-1]%[ma:ec575586]
	Khảo sát sự biến thiên và vẽ đồ thị của các hàm  số sau
	\begin{multicols}{2}
		\begin{enumerate}
			\item $y=\dfrac{x^2-x-1}{x-2}$; 
			\item  $y=\dfrac{x^2+x-2}{x+1}$;
			\item $y=\dfrac{2x^2-x+4}{x-1}$; 
			\item $y=\dfrac{x^2+2x+1}{x+3}$. 
		\end{enumerate}
	\end{multicols}
	\loigiai{ 
		\begin{enumerate}
			\item  Hàm số $y=\dfrac{x^2-x-1}{x-2}$;\\
			Tập  xác định: $D=\mathbb{R}\setminus\{2\}$.\\
			Sự biến thiên:  $y=x+1+\dfrac{1}{x-2}$.\\
			 Ta có $y'=1-\dfrac{1}{(x-2)^2}=\dfrac{x^2-4x+3}{(x-2)^2}$.\\
			 Khi đó $y'=0\Leftrightarrow \dfrac{x^2-4x+3}{(x-2)^2}=0\Leftrightarrow\hoac{&x=1\\&x=3.}$\\
			 Ta có
			\begin{itemize}
				\item $\lim\limits_{x\to-\infty}y=\lim\limits_{x\to-\infty}\dfrac{x^2-x-1}{x-2}=-\infty;  \lim\limits_{x\to+\infty}y=\lim\limits_{x\to+\infty}\dfrac{x^2-x-1}{x-2}=+\infty$.
				\item $\lim\limits_{x\to2^-}y=\lim\limits_{x\to2^-}\left(x+1+\dfrac{1}{x-2}\right)=-\infty;  \lim\limits_{x\to2^+}y=\lim\limits_{x\to2^+}\left(x+1+\dfrac{1}{x-2}\right)=+\infty$.
				\item $\lim\limits_{x\to+\infty}\left[y-(x+1)\right]=\lim\limits_{x\to+\infty}\dfrac{1}{x-2}=0; \lim\limits_{x\to-\infty}\left[y-(x+1)\right]=\lim\limits_{x\to-\infty}\dfrac{1}{x-2}=0$. 
			\end{itemize}
			Do đó, đồ thị hàm số có tiệm cận đứng là đường thẳng $x=2$, tiệm cận xiên  là đường thẳng $y=x+1$.\\
			 Bảng biến thiên
			\begin{center}
				\begin{tikzpicture}
				\tikzset{double style/.append style={double distance=1.5pt}}
				 \tkzTabInit[nocadre=false,lgt=1.2,espcl=2.5,deltacl=0.6]
					 {$x$ /0.6,$y'$ /0.6,$y$ /2}
					 {$-\infty$,$1$,$2$,$3$,$+\infty$}
				\tkzTabLine{,+,$0$,-,d,-,$0$,+,} 
				\tkzTabVar{-/$-\infty$,+/$1$,-D+/$-\infty$/$+\infty$,-/$5$,+/$+\infty$} 
			\end{tikzpicture}
			\end{center}
			Đồ  thị
			\begin{center}
				\begin{tikzpicture}[line join=round,line cap=round,>=stealth,scale=0.7]
					\tikzset{ethicknode/.style={scale=0.9}} 
					\draw[->]  (-2.6,0)--(6.1,0) node[below left] {$x$};
					\draw[->] (0,-1.6)--(0,7.1) node[below left] {$y$}; 
					\draw (0,0) node [below left] {$O$}; 
					\foreach  \x/\nx in {-2/-2,-1/,1/1,2/,3/3,4/4,5/5}
					\draw[thin] (\x,1pt)--(\x,-1pt) node [below] {$\nx$}; 
					\draw (-1,0) node  [above]{$-1$}
						 (2,0) node[below right]{$2$};
					\foreach \y/\ny in  {-1/-1,1/1,2/2,3/3,4/4,5/5,6/6}
					 \draw[thin] (1pt,\y)--(-1pt,\y) node [left] {$\ny$};
					\draw[dashed,thin](1,0)--(1,1)--(0,1); 
					\draw[dashed,thin](3,0)--(3,5)--(0,5); 
					\draw[dashed,thin]  (2,-1.5)--(2,7) node[pos=0.9,above,rotate=90]{$x=2$};
					\begin{scope}
						 \clip (-2.5,-1.5) rectangle (6,7);
						 \draw[samples=200,domain=-2.1:1.99,smooth,variable=\x] plot (\x,{(1*((\x)^2)+-1*(\x)+-1)/(1*(\x)+-2)});
						 \draw[samples=200,domain=2.01:6,smooth,variable=\x] plot (\x,{(1*((\x)^2)+-1*(\x)+-1)/(1*(\x)+-2)});
						\draw[dashed,thin] (-2.2,-1.2)--(6.1,7.1) node[pos=0.8,below,rotate=45]{$y=x+1$}; 
					\end{scope}
					\draw  (2,-2) node{\textbf{\textit{Hình 1.31}}};
				\end{tikzpicture}
			\end{center}
			\begin{itemize}
				\item Giao điểm của đồ thị hàm số với  trục tung là điểm $\left(0; \dfrac{1}{2}\right)$.
				\item Ta có $y=0 \Leftrightarrow \dfrac{x^{2}-x-1}{x-2}=0 \Leftrightarrow x=\dfrac{1-\sqrt{5}}{2}$ hoặc $x=\dfrac{1+\sqrt{5}}{2}$. Do đó giao điểm của  đồ thị hàm số với trục hoành là các điểm $\left(\dfrac{1-\sqrt{5}}{2}; 0\right)$ và $\left(\dfrac{1+\sqrt{5}}{2}; 0\right)$.
				\item Đồ thị hàm số nhận giao điểm $I(2; 3)$ của hai đường tiệm cận làm tâm đối xứng và nhận hai đường phân giác của các góc  tạo bởi hai đường tiệm cận này làm các trục đối xứng.
			\end{itemize}
			\item Hàm  số $y=\dfrac{x^2+x-2}{x+1}$;\\
			Tập xác định của hàm số: $\mathbb{R} \setminus\{-1\}$.\\ 
			Sự biến thiên: 
			\begin{itemize}
				\item Viết $y=x-\dfrac{2}{x+1}$, ta có $y^{\prime}=1+\dfrac{2}{(x+1)^{2}}>0$  với mọi $x \neq-1$.
				\item Hàm số đồng biến trên từng khoảng $(-\infty;-1)$ và $(-1;+\infty)$. 
				\item  Hàm số không có cực trị.
				 \item $\lim\limits _{x \rightarrow-\infty} y=\lim\limits _{x \rightarrow-\infty} \dfrac{x^{2}+x-2}{x+1}=\lim\limits _{x \rightarrow-\infty} \dfrac{x+1-\dfrac{2}{x}}{1+\dfrac{1}{x}}=-\infty$;\\ $\lim\limits _{x \rightarrow+\infty} y=\lim\limits _{x \rightarrow+\infty} \dfrac{x^{2}+x-2}{x+1}=\lim\limits _{x \rightarrow+\infty} \dfrac{x+1-\dfrac{2}{x}}{1+\dfrac{1}{x}}=+\infty$.
				\item Tiệm  cận:
				\begin{align*}
					&\lim\limits _{x \rightarrow-1^{-}} y=\lim\limits  _{x \rightarrow-1^{-}}\left(x-\dfrac{2}{x+1}\right)=+\infty; \lim\limits _{x \rightarrow-1^{+}} y=\lim\limits _{x \rightarrow-1^{+}}\left(x-\dfrac{2}{x+1}\right)=-\infty;\\
					&\lim\limits  _{x \rightarrow+\infty}[y-x]=\lim\limits _{x \rightarrow+\infty}\left(-\dfrac{2}{x+1}\right)=0; \lim\limits _{x \rightarrow-\infty}[y-x]=\lim\limits _{x \rightarrow-\infty}\left(-\dfrac{2}{x+1}\right)=0.
				\end{align*}
			\end{itemize}
			Do đó, đồ thị hàm số có tiệm cận  đứng là đường thẳng $x=-1$, tiệm cận xiên là đường thẳng $y=x$.\\
			 Bảng biến thiên:
			\begin{center}
				\begin{tikzpicture}[line join=round,font=\footnotesize,>=stealth,scale=1]
					 \tkzTabInit[nocadre=false,lgt=1.2,espcl=2.5,deltacl=0.6]
						 {$x$ /0.6,$y'$ /0.6,$y$ /2}
						 {$-\infty$,$-1$,$+\infty$}
					\tkzTabLine{,+,d,+,} 
					 \tkzTabVar{-/$-\infty$,+D-/$+\infty$/$-\infty$,+/$+\infty$}
				\end{tikzpicture}
			\end{center}
			Đồ thị 
			\begin{center}
				\begin{tikzpicture}[line cap=round,line join=round,font=\footnotesize,>=stealth,scale=1,x=0.5cm,y=0.5cm]
					\tikzset{label  style/.style={font=\footnotesize}}
					 \draw[->] (-6.1,0)--(6.1,0) node[below] {$x$};
					\draw[->] (0,-5.1)--(0,5.1) node[left]  {$y$};
					\draw  (0,0) node [above left] {\scriptsize$O$};
					\foreach \x/\nx  in {-6/-6,-5/-5,-4/-4,-3/-3,-2/-2,-1/,1/1,2/2,3/3,4/4,5/5}
					\draw[thin] (\x,1pt)--(\x,-1pt) node  [below] {\scriptsize$\nx$};
					\foreach \y/\ny  in {-4/right,-3/right,-2/right,-1/left,1/left,2/left,3/left,4/left}
					\draw[thin] (1pt,\y)--(-1pt,\y) node [\ny]  {\scriptsize $\y$};
					\draw[dashed,thin]  (-1,-5)--(-1,5) node[above,rotate=90,pos=0.1]{$x=-1$};
					\draw (-1,0) node[above left]{$-1$}; 
					\begin{scope}
						\clip (-6,-5) rectangle  (6,5);
						\draw[samples=200,domain=-6:-1.01,smooth,variable=\x,magenta]  plot (\x,{(1*((\x)^2)+1*(\x)+-2)/(1*(\x)+1)});
						 \draw[samples=200,domain=-0.99:6,smooth,variable=\x,magenta] plot (\x,{(1*((\x)^2)+1*(\x)+-2)/(1*(\x)+1)});
						 \draw[dashed,thin] (-6.1,-6.1)--(6.1,6.1) node[above,rotate=45,pos=0.8]{$y=x$};
					\end{scope}
					 \draw (0,-5.5) node{\textbf{\textit{Hình 1.32}}};
				\end{tikzpicture}
			\end{center}
			\begin{itemize}
				\item Giao điểm của đồ thị hàm số với  trục tung là điểm $(0;-2)$.
				\item Ta  có $y=0 \Leftrightarrow \dfrac{x^{2}+x-2}{x+1}=0 \Leftrightarrow x=-2$ hoặc $x=1$.
				Do đó giao điểm của đồ thị hàm số với trục  hoành là các điểm $(-2; 0)$ và $(1; 0)$.
				\item Đồ  thị hàm số nhận giao điểm $I(-1;-1)$ của hai đường tiệm cận làm tâm đối xứng và nhận hai đường phân giác của góc tạo bởi hai đường tiệm cận này làm các trục đối xứng.
			\end{itemize}
			\item Hàm  số $y=\dfrac{2x^2-x+4}{x-1}$;\\
			Tập xác định của hàm số: $\mathscr{D}=\mathbb{R}  \setminus\{1\}$.\\
			Sự biến thiên: 
			\begin{itemize}
				\item Viết $y=2x+1+\dfrac{5}{x-1}$, ta  có $y^{\prime}=2-\dfrac{5}{(x-1)^{2}}$ với mọi $x \neq 1$.
				Suy ra 
				\[y'=0\Leftrightarrow x=1\pm  \dfrac{\sqrt{10}}{2}. \]
				\item Trên các khoảng $\left( -\infty;1-\dfrac{\sqrt{10}}{2} \right)$  và $\left(1+\dfrac{\sqrt{10}}{2};+\infty \right)$ thì $y'>0$ nên hàm số đồng biến trên các khoảng này.\\
				Trên các khoảng $\left(1-\dfrac{\sqrt{10}}{2};1 \right)$ và $\left( 1;1+\dfrac{\sqrt{10}}{2} \right)$ thì  $y'<0$ nên hàm số nghịch biến trên các khoảng này.
				\item Đồ thị hàm số có điểm cực đại là $(x_{\text{ cđ}};y_{\text{ cđ}})=\left(1-\dfrac{\sqrt{10}}{2};3-2\sqrt{10} \right)$ và  điểm cực tiểu là $(x_{\text{ ct}},y_{\text{ ct}})=\left(1+\dfrac{\sqrt{10}}{2}; 3+2\sqrt{10} \right)$.
				\item  $\lim\limits _{x \rightarrow-\infty} y=-\infty$; $\lim\limits _{x \rightarrow+\infty} y=+\infty$.
				 \item Tiệm cận:
				\begin{align*}
					&\lim\limits _{x \rightarrow 1^{-}} y=-\infty; \lim\limits _{x \rightarrow 1^{+}} y=+\infty;\\ 
					&\lim\limits _{x \rightarrow+\infty}[y-(2x+1)]=  \lim\limits _{x \rightarrow-\infty}[y-(2x+1)]=0.
				\end{align*}
				Do đó, đồ thị hàm số có tiệm cận đứng là đường thẳng $x=1$, tiệm cận xiên là đường  thẳng $y=2x+1$.\\
				Bảng biến thiên 
				\begin{center}
					\begin{tikzpicture}[line join=round,font=\footnotesize,>=stealth,scale=1]
						 \tkzTabInit[nocadre=false,lgt=1.2,espcl=2.5,deltacl=0.6]
						{$x$ /1,$y'$ /1,$y$  /3}
						{$-\infty$,$1-\dfrac{\sqrt{10}}{2}$,$1$,$1+\dfrac{\sqrt{10}}{2}$,$+\infty$} 
						 \tkzTabLine{,+,0,-,d,-,0,+}
						\tkzTabVar{-/$-\infty$,+/$3-2\sqrt{10}$,-D+/$-\infty$/$+\infty$,-/$3+2\sqrt{10}$,+/$+\infty$} 
					\end{tikzpicture}
				\end{center}
			\end{itemize}
			Đồ thị 
			\begin{center}
				\begin{tikzpicture}[line cap=round,line join=round,font=\footnotesize,>=stealth,scale=1,x=0.5cm,y=0.25cm]
					\tikzset{label  style/.style={font=\footnotesize}}
					\draw[->] (-4,0)--(6,0) node[below] {$x$}; 
					\draw[->] (0,-8)--(0,15) node[left]  {$y$};
					\draw[smooth,samples=100,color=red] plot[domain=-4:0.5]  (\x,{(2*(\x)^2-(\x)+4)/(\x-1)});
					\draw[smooth,samples=100,color=red] plot[domain=1.5:6] (\x,{(2*(\x)^2-(\x)+4)/(\x-1)}); 
					\draw[dashed] (1,-8) node  [right]{$x=1$}--(1,15)
						plot[domain=-4:6](\x,{2*(\x)+1})  node[rotate=45,below]{$y=2x+1$}
						({1-sqrt  (10)/2},0) node[above]{$x_{\text{ cđ}}$}--({1-sqrt (10)/2},{3-2*sqrt (10)})--(0,{3-2*sqrt (10)}) node[right]{$y_{\text{ cđ}}$}
						({1+sqrt (10)/2},0) node[below]{$x_{\text{ ct}}$}--({1+sqrt  (10)/2},{3+2*sqrt (10)})--(0,{3+2*sqrt (10)}) node[left]{$y_{\text{ ct}}$};
				\end{tikzpicture}
			\end{center}
			\begin{itemize}
				\item Giao  điểm của đồ thị hàm số với trục tung là điểm $(0;-4)$.
				\item Ta có $y=0 \Leftrightarrow 2x^{2}-x+4=0$. Phương trình này vô nghiệm nên đồ thị hàm số không cắt  trục hoành.
				\item Đồ thị hàm số nhận giao điểm $I(1;3)$ của hai đường tiệm cận làm tâm đối xứng và nhận hai đường phân giác của góc tạo  bởi hai đường tiệm cận này làm các trục đối xứng.
			\end{itemize}
			 \item Hàm số $y=\dfrac{x^2+2x+1}{x+3}$;\\
			Tập  xác định của hàm số: $\mathscr{D}=\mathbb{R} \setminus\{-3\}$.\\
			 Sự biến thiên:
			\begin{itemize}
				\item Viết $y=x-1+\dfrac{4}{x+3}$, ta có $y^{\prime}=1-\dfrac{4}{(x+3)^{2}}$ với mọi  $x \neq -3$.
				Suy ra  $ y'=0\Leftrightarrow x=-5$ hoặc $x=-1$.
				\item Trên các khoảng $\left( -\infty;-5 \right)$ và $\left(-1;+\infty \right)$ thì $y'>0$  nên hàm số đồng biến trên các khoảng này.
				Trên các khoảng $\left(-5;-3 \right)$ và $\left( -3;-1\right)$ thì $y'<0$ nên hàm số nghịch biến trên các khoảng này. 
				\item Đồ thị hàm số  có điểm cực đại là $(x_{\text{ cđ}};y_{\text{ cđ}})=\left(-5;-8\right)$ và điểm cực tiểu là $(x_{\text{ ct}},y_{\text{ ct}})=\left(-1; 0 \right)$.
				\item $\lim\limits _{x \rightarrow-\infty} y=-\infty$; $\lim\limits _{x \rightarrow+\infty} y=+\infty$. 
				\item Tiệm  cận:
				\begin{align*}
					&\lim\limits _{x \rightarrow  (-3)^{-}} y=-\infty; \lim\limits _{x \rightarrow (-3)^{+}} y=+\infty;\\
					&\lim\limits  _{x \rightarrow+\infty}[y-(x-1)]= \lim\limits _{x \rightarrow-\infty}[y-(x-1)]=0.
				\end{align*}
				Do đó, đồ thị hàm số có tiệm cận đứng là đường thẳng $x=1$, tiệm cận xiên là đường thẳng $y=x-1$.\\ 
				Bảng biến  thiên
				\begin{center}
					\begin{tikzpicture}[line join=round,font=\footnotesize,>=stealth,scale=1]
						 \tkzTabInit[nocadre=false,lgt=1.2,espcl=2.5,deltacl=0.6]
						{$x$  /1,$y'$ /1,$y$ /3}
						{$-\infty$,$-5$,$-3$,$-1$,$+\infty$} 
						\tkzTabLine{,+,0,-,d,-,0,+} 
						\tkzTabVar{-/$-\infty$,+/$-8$,-D+/$-\infty$/$+\infty$,-/$0$,+/$+\infty$} 
					\end{tikzpicture}
				\end{center}
			\end{itemize}
			Đồ  thị
			\begin{center}
				\begin{tikzpicture}[line cap=round,line join=round,font=\footnotesize,>=stealth,scale=1,x=0.3cm,y=0.3cm]
					\tikzset{label  style/.style={font=\footnotesize}}
					 \draw[->] (-10,0)--(8,0) node[below] {$x$};
					\draw[->] (0,-12)--(0,7) node[left]  {$y$};
					\draw[smooth,samples=100,color=red]  plot[domain=-10:-3.5] (\x,{((\x)^2+2*(\x)+1)/(\x+3)});
					\draw[smooth,samples=100,color=red]  plot[domain=-2.6:8] (\x,{((\x)^2+2*(\x)+1)/(\x+3)});
					\draw[dashed] (-3,-12)--(-3,7) node [above left,rotate=90]{$x=-3$} 
						plot[domain=-10:8](\x,{(\x)-1})  node[rotate=45,below]{$y=x-1$}
						 (-5,0) node[above]{$x_{\text{ cđ}}$}--(-5,-8)--(0,-8) node[right]{$y_{\text{ cđ}}$};
				\end{tikzpicture}
			\end{center}
			\begin{itemize}
				\item  Giao điểm của đồ thị hàm số với trục tung là điểm $\left(0;\dfrac{1}{3}\right)$.
				\item Ta có $$y=0 \Leftrightarrow x^2+2x+1=0  \Leftrightarrow x=-1.$$
				Đồ thị hàm số cắt trục hoành  tại điểm $(0,-1)$.
				\item Đồ thị hàm số nhận giao điểm $I(-3;-4)$ của hai đường tiệm cận làm tâm đối xứng và nhận hai đường phân giác của góc tạo bởi hai đường tiệm cận  này làm các trục đối xứng.
			\end{itemize}
		\end{enumerate}
	} 
\end{bt}

% nguồn: Phan2_Lop_12_Dai_so_Chuong_1
\begin{bt}%[2D1H5-1]%[ma:377ece2a]
	Khảo sát  sự biến thiên và vẽ đồ thị các hàm số sau
	\begin{multicols}{2}
		\begin{enumerate}
			\item  $y = \dfrac{x+1}{x-2}$.
			 \item $y = \dfrac{2x+1}{x+1}$.
			\item $y  = \dfrac{x+3}{1-x}$.
			\item $y =  \dfrac{5+x}{2-x}$.
		\end{enumerate}
	\end{multicols}
	\loigiai{ 
		\begin{enumerate}
			\item  Xét hàm số $y = \dfrac{x+1}{x-2}$.
			\begin{itemize}
				 \item Tập xác định $\mathscr{D}=\mathbb{R} \setminus \{2\}$.
				\item Đạo hàm $y' = \dfrac{-3}{(x-2)^2} < 0$ với mọi  $x \ne 2$. \\
				Hàm  số nghịch biến trên từng khoảng $(-\infty; 2)$ và $(2; +\infty)$ và không có cực trị.
				\item Tiệm cận 
				\begin{itemize}
					 \item Ta có $\lim\limits_{x\to 2^-}y=\lim\limits_{x\to 2^-}\dfrac{x+1}{x-2}=-\infty$; $\lim\limits_{x\to 2^+}y=\lim\limits_{x\to 2^+}\dfrac{x+1}{x-2}=+\infty
					\Rightarrow$ Tiệm cận đứng của đồ thị hàm số là  đường thẳng $x = 2$.
					\item $\lim\limits_{x\to+\infty}y=\lim\limits_{x\to+\infty}\dfrac{x+1}{x-2}=1$; $\lim\limits_{x\to-\infty}y=\lim\limits_{x\to-\infty}\dfrac{x+1}{x-2}=1\Rightarrow$ Tiệm cận ngang của đồ thị hàm số là đường thẳng $y =  1$.
				\end{itemize}
				\item Bảng  biến thiên
				\begin{center}
					\begin{tikzpicture}
						\tkzTabInit[nocadre=false,lgt=1.5,espcl=4,deltacl=0.8] 
						{$x$ /1,$f'(x)$ /1,$f (x)$  /2}
						{$-\infty$,$2$,$+\infty$} 
						 \tkzTabLine{,-,d,-,}
						\tkzTabVar{+/ $1$,-D+/$-\infty$ / $+\infty$,-/  $1$}
					\end{tikzpicture}
				\end{center}
				\item Giao điểm với trục tung $x = 0 \Rightarrow y = \dfrac{1}{-2} = -\dfrac{1}{2} \Rightarrow \left(0; -\dfrac{1}{2}\right)$. 
				\item Giao điểm với trục hoành $y = 0 \Rightarrow x +  1 = 0 \Rightarrow x = -1 \Rightarrow (-1; 0)$.
				\item Nhận xét: Đồ thị nhận điểm $I(2;1)$ là tâm đối xứng và đồ  thị nhận hai đường phân giác làm trục đối xứng.
			\end{itemize}
			\begin{center}
				\begin{tikzpicture}[scale=0.5]
					\draw[->] (-6,0)--(9,0)  node[below]{$x$};
					\draw[->]  (0,-7)--(0,7) node[left]{$y$};
					\draw[dashed] (2,-7)--(2,7) node[below left] {$x=2$}; 
					\draw[dashed] (-6,1)--(9,1) node[below left]  {$y=1$};
					\draw[domain=-6:1.62,smooth,variable=\x]  plot ({\x},{(\x+1)/(\x-2)});
					 \draw[domain=2.5:9,smooth,variable=\x] plot ({\x},{(\x+1)/(\x-2)});
					\fill (0,0) circle (1.3pt) node[above  right]{$O$};
				\end{tikzpicture}
			\end{center}
			\item  Xét hàm số $y = \dfrac{2x+1}{x+1}$.
			\begin{itemize}
				\item Tập  xác định $\mathscr{D}=\mathbb{R} \setminus \{-1\}$.
				\item Đạo hàm $y' =  \dfrac{1}{(x+1)^2} > 0$ với mọi $x \ne -1$. \\
				Hàm số đồng biến trên từng khoảng $(-\infty; -1)$ và $(-1;  +\infty)$ và không có cực trị.
				\item Tiệm cận 
				\begin{itemize}
					\item $\lim\limits_{x\to (-1)^-}y=\lim\limits_{x\to (-1)^-}\dfrac{2x+1}{x+1}=+\infty$; $\lim\limits_{x\to (-1)^+}y=\lim\limits_{x\to  (-1)^+}\dfrac{2x+1}{x+1}=-\infty
					\Rightarrow$ Tiệm cận đứng của đồ thị hàm số là đường thẳng  $x = -1$.
					\item $\lim\limits_{x\to+\infty}y=\lim\limits_{x\to+\infty}\dfrac{2x+1}{x+1}=2$; $\lim\limits_{x\to-\infty}y=\lim\limits_{x\to-\infty}\dfrac{2x+1}{x+1}=2\Rightarrow$ Tiệm cận ngang của đồ thị hàm số là  đường thẳng $y = 2$.
				\end{itemize}
				 \item Bảng biến thiên
				\begin{center}
					\begin{tikzpicture}
						\tkzTabInit[nocadre=false,lgt=1.5,espcl=4,deltacl=0.8] 
						 {$x$ /1,$f'(x)$ /1,$f (x)$ /2}
						 {$-\infty$,$-1$,$+\infty$}
						\tkzTabLine{,+,d,+,} 
						\tkzTabVar{-/ $2$,+D-/$+\infty$ / $-\infty$,+/ $2$} 
					\end{tikzpicture}
				\end{center}
				\item Giao  điểm với trục tung $x = 0 \Rightarrow y = \dfrac{1}{1} = 1 \Rightarrow (0;1)$.
				\item Giao điểm với trục hoành $y = 0 \Rightarrow 2x + 1 = 0 \Rightarrow x =  -\dfrac{1}{2} \Rightarrow \left(-\dfrac{1}{2}; 0\right)$.
				\item Nhận xét: Đồ thị nhận điểm $I(-1;2)$ là tâm đối xứng và đồ thị nhận hai đường  phân giác làm trục đối xứng.
			\end{itemize}
			\begin{center}
				\begin{tikzpicture}[scale=0.8]
					 \draw[->] (-6,0)--(4,0) node[below]{$x$};
					 \draw[->] (0,-3)--(0,6) node[left]{$y$};
					\draw[dashed] (-1,6) -- (-1,-3) node[above  left]{$x=-1$};
					\draw[dashed] (-6,2)--(4,2)  node[above left]{$y=2$};
					\draw[domain=-6:-1.26,smooth,variable=\x] plot ({\x},{(2*\x+1)/(\x+1)}); 
					\draw[domain=-0.8:4,smooth,variable=\x] plot ({\x},{(2*\x+1)/(\x+1)}); 
					\fill (0,0)  circle (1.4pt) node[below right]{$O$};
				\end{tikzpicture}
			\end{center}
			\item Xét hàm số $y  = \dfrac{x+3}{1-x}$.
			\begin{itemize}
				\item Tập xác định $\mathscr{D} = \mathbb{R} \setminus \{1\}$. 
				 \item Đạo hàm $y' = \dfrac{4}{(1-x)^2} > 0$ với mọi $x \ne 1$. \\
				Hàm số đồng biến trên các khoảng $(-\infty;1)$ và $(1;+\infty)$ và không có  cực trị.
				\item  Tiệm cận
				\begin{itemize}
					\item $\lim\limits_{x\to 1^-}y=\lim\limits_{x\to 1^-}\dfrac{x+3}{1-x}=+\infty$;  $\lim\limits_{x\to 1^+}y=\lim\limits_{x\to 1^+}\dfrac{x+3}{1-x}=-\infty
					\Rightarrow$  Tiệm cận đứng của đồ thị hàm số là đường thẳng $x = 1$.
					\item $\lim\limits_{x\to+\infty}y=\lim\limits_{x\to+\infty}\dfrac{x+3}{1-x}=-1$; $\lim\limits_{x\to-\infty}y=\lim\limits_{x\to-\infty}\dfrac{x+3}{1-x}=-1\Rightarrow$ Tiệm cận ngang của đồ thị hàm số là  đường thẳng $y = -1$.
				\end{itemize}
				\item Bảng biến thiên 
				\begin{center}
					\begin{tikzpicture}
						\tkzTabInit[nocadre=false,lgt=1.5,espcl=4,deltacl=0.8] 
						{$x$ /1,$f'(x)$ /1,$f  (x)$ /2}
						{$-\infty$,$1$,$+\infty$} 
						 \tkzTabLine{,+,d,+,}
						\tkzTabVar{-/ $-1$,+D-/$+\infty$ / $-\infty$,+/ $-1$} 
					\end{tikzpicture}
				\end{center}
				\item Giao điểm với trục tung $(0;3)$, giao điểm với  trục hoành $(-3;0)$.
				\item Đồ thị  nhận giao điểm $I(1;-1)$ làm tâm đối xứng.
			\end{itemize}
			\begin{center}
				\begin{tikzpicture}[scale=0.5]
					\draw[->] (-7,0)--(9,0)  node[below]{$x$};
					\draw[->] (0,-7)--(0,7)  node[left]{$y$};
					\draw[dashed] (1,-7)--(1,7) node[below  right]{$x=1$};
					\draw[dashed] (9,-1)--(-7,-1) node[below right]{$y=-1$}; 
					\draw[domain=-7:0.5,smooth,variable=\x] plot  ({\x},{(\x+3)/(1-\x)});
					\draw[domain=1.7:9,smooth,variable=\x] plot ({\x},{(\x+3)/(1-\x)}); 
					\fill (0,0)  circle (1.4pt) node[above right]{$O$};
				\end{tikzpicture}
			\end{center}
			\item Xét  hàm số $y = \dfrac{5+x}{2-x}$.
			\begin{itemize}
				\item Tập xác định $\mathscr{D}=\mathbb{R}  \setminus \{2\}$.
				\item Đạo hàm $y' = \dfrac{7}{(-x+2)^2} > 0$  với mọi $x \ne 2$. \\ Hàm số đồng biến trên các khoảng $(-\infty;2)$ và $(2;+\infty)$.
				 \item Tiệm cận
				\begin{itemize}
					\item $\lim\limits_{x\to 2^-}y=\lim\limits_{x\to 2^-}\dfrac{5+x}{2-x}=+\infty$; $\lim\limits_{x\to 2^+}y=\lim\limits_{x\to 2^+}\dfrac{5+x}{2-x}=-\infty 
					\Rightarrow$ Tiệm cận đứng của đồ thị hàm số  là đường thẳng $x = 2$.
					\item $\lim\limits_{x\to+\infty}y=\lim\limits_{x\to+\infty}\dfrac{5+x}{2-x}=-1$; $\lim\limits_{x\to-\infty}y=\lim\limits_{x\to-\infty}\dfrac{5+x}{2-x}=-1\Rightarrow$ Tiệm cận ngang của đồ thị hàm số là  đường thẳng $y = -1$.
				\end{itemize}
				 \item Bảng biến thiên
				\begin{center}
					\begin{tikzpicture}
						 \tkzTabInit[nocadre=false,lgt=1.5,espcl=4,deltacl=0.8]
						{$x$  /1,$f'(x)$ /1,$f (x)$ /2}
						 {$-\infty$,$2$,$+\infty$}
						 \tkzTabLine{,+,d,+,}
						\tkzTabVar{-/ $-1$,+D-/$+\infty$ /  $-\infty$,+/ $-1$}
					\end{tikzpicture}
				\end{center}
				\item Giao  điểm với trục tung $(0;3)$, giao điểm với trục hoành $(-5;0)$.
				\item Đồ thị nhận  điểm $I(2;-1)$ là tâm đối xứng.
			\end{itemize}
			\begin{center}
				\begin{tikzpicture}[scale=0.4]
					 \draw[->] (-8,0)--(11,0) node[below]{$x$};
					\draw[->] (0,-8)--(0,8) node[left]{$y$}; 
					\draw[dashed] (2,-8)--(2,8) node[below right]{$x=2$}; 
					\draw[dashed] (11,-1)--(-8,-1)  node[below right]{$y=-1$};
					\draw[domain=-8:1.25,smooth,variable=\x] plot  ({\x},{(5+\x)/(2-\x)});
					\draw[domain=2.95:11,smooth,variable=\x] plot  ({\x},{(5+\x)/(2-\x)});
					\fill (0,0) circle (1.4pt) node[above right]{$O$}; 
				\end{tikzpicture}
			\end{center}
		\end{enumerate}
	} 
\end{bt}

% nguồn: Phan2_Lop_12_Dai_so_Chuong_1
\begin{bt}%[2D1V5-1]%[ma:e8b2e572]
	 \immini{
	Cho hàm số bậc ba $f(x)=x^3+bx^2+cx+d$. Biết đồ thị của hàm số $y=f'(x)$ như hình vẽ. Tính $\dfrac{c}{b}$  (làm tròn đến hàng phần chục).
	}{ 
	\begin{tikzpicture}[scale=1,>=stealth,font=\footnotesize,line join=round,line cap=round]
		\def\a{1} \def\b{-2} \def\c{0.75} 
		 \def\xmin{-1} \def\xmax{2.5}
		 \def\ymin{-1} \def\ymax{2}
		 \draw[->] (\xmin,0)--(\xmax,0) node [below]{$x$};
		\draw[->] (0,\ymin)--(0,\ymax)  node [left]{$y$};
		\node at (0,0) [below  left]{$O$};
		\fill (1,0) node[above]{$1$} circle (1.2pt); 
		 \fill (0.5,0) node[below]{$\dfrac{1}{2}$} circle (1.2pt);
		\fill (1.5,0) node[below]{$\dfrac{3}{2}$}  circle (1.2pt);
		\clip (\xmin+0.1,\ymin+0.1)  rectangle (\xmax-0.5,\ymax-0.1);
		 \draw[smooth,samples=300] plot (\x,{\a*(\x)^2+\b*(\x)+\c});
	\end{tikzpicture}
	} 
	 \loigiai{
		Tập xác định $\mathscr{D}=\mathbb{R}$.\\ 
		Đạo hàm cấp $1$ của $f(x)$ là  $f'(x)=3ax^2+2bx+c$.\\
		Dựa vào đồ thị của hàm số $f'(x)$ ta có bảng biến thiên  của hàm số $f(x)$ như sau:
		\begin{center}
			\begin{tikzpicture}
				\tkzTabInit[espcl=2.5,lgt=1.5,nocadre=false] 
					{$x$/1,$f'(x)$/0.7,$f  (x)$/2.1}
					{$-\infty$,$\dfrac{1}{2}$,$\dfrac{3}{2}$,$+\infty$} 
				 \tkzTabLine{,+,0,-,0,+,}
				 \tkzTabVar{-/$-\infty$,+/$f\left(\tfrac{1}{2} \right)$,-/$f\left(\tfrac{3}{2} \right)$,+/$+\infty$}
			\end{tikzpicture}
		\end{center}
		Ta có $f'\left(\dfrac{1}{2}\right)=\dfrac{3a}{4}+b+c$  và $f'\left(\dfrac{3}{2}\right)=\dfrac{27a}{4}+3b+c$.\\
		Dựa vào bảng  biến thiên ta có
		 $$\heva{&3a+4b+4c=0\\&27a+12b+4c=0} \Leftrightarrow \heva{&27a+36b+36c=0\\&27a+12b+4c=0}
		\Rightarrow 24b+32c=0 \Rightarrow  \dfrac{c}{b}=-\dfrac{3}{4}.$$\\
		Vậy $\dfrac{c}{b}=-\dfrac{3}{4}\approx-0{,}8$. 
	} 
\end{bt}
